% L'exercice sur texte — Philosophie 1re Tech — Cours signé OnBuch+
% Programme officiel MINESEC. Compiler avec : tectonic N21-exercice-texte.tex
\newcommand{\DOCMATIERE}{Philosophie}
\newcommand{\DOCNIVEAU}{1re Tech}
\newcommand{\DOCMODULE}{Philosophie – classe de Première technique}
\newcommand{\DOCLECON}{N21}
\newcommand{\DOCTITRE}{L'exercice sur texte}
% OnBuch+ — préambule « cours signés » ENRICHI (compilé avec tectonic / XeLaTeX)
% Système visuel « Pop » d'OnBuch+ : fond crème (jamais blanc), bordures encre
% épaisses, ombres « dures » décalées, titres Archivo Black en capitales.
% Version MATHÉMATIQUES : graphes (pgfplots/tikz), boîtes pédagogiques
% (définition, propriété, méthode, attention, le savais-tu, exercice résolu,
% exercice, TP, bilan), page de garde et signature OnBuch+ ancrée en bas.
%
% Macros attendues dans main.tex AVANT \input{preamble} :
%   \DOCMATIERE  \DOCNIVEAU  \DOCMODULE  \DOCLECON (numéro)  \DOCTITRE
\documentclass[11pt]{article}

\usepackage{fontspec}
\usepackage[french]{babel}
\usepackage[a4paper,margin=2cm,top=3.4cm,bottom=2.8cm,headheight=1.4cm,footskip=1.3cm]{geometry}
\usepackage{amsmath,amssymb,amsfonts}
\usepackage{mathtools} % \xmapsto, \coloneqq... souvent utilisés par les modèles
\usepackage{cancel} % simplifications barrées (bilans de forces)
% \abs{z} (module, valeur absolue) et \norm{u} (norme), utilisés par les modèles
\providecommand{\abs}{}\let\abs\relax\DeclarePairedDelimiter{\abs}{\lvert}{\rvert}
\providecommand{\norm}{}\let\norm\relax\DeclarePairedDelimiter{\norm}{\lVert}{\rVert}
\usepackage[table]{xcolor} % [table] : fournit aussi \rowcolors (alternance de lignes)
\usepackage{pagecolor}
\usepackage{tikz}
\usetikzlibrary{arrows.meta,positioning,calc,shapes.geometric,shapes.misc,shapes.arrows,decorations.pathmorphing,decorations.pathreplacing,decorations.markings,patterns,shadows,mindmap,trees,fit,backgrounds,matrix,intersections,angles,quotes}
\usepackage{pgfplots}
\usepackage[version=4]{mhchem} % équations nucléaires \ce{^{235}_{92}U}
\usepackage[european,siunitx]{circuitikz} % schémas électriques
\pgfplotsset{compat=1.18}
\usepgfplotslibrary{fillbetween,groupplots} % aires entre courbes (calcul intégral)
\usepackage{enumitem}
\usepackage{fancyhdr}
% [most] charge toutes les bibliothèques tcolorbox (listings, minted, raster,
% documentation...) et gonfle inutilement la mémoire TeX sur un document long
% (~300 boîtes) : on ne charge que ce qui est réellement utilisé.
\usepackage[skins,breakable]{tcolorbox}
\usepackage{graphicx}
\usepackage{colortbl}
\usepackage{booktabs}
\usepackage{tabularx}
\usepackage{multirow}
\usepackage{diagbox} % case d'en-tête diagonale des tableaux à double entrée
% Colonnes extensibles centrée / gauche / droite (C, L, R), souvent utilisées par les modèles
\newcolumntype{C}{>{\centering\arraybackslash}X}
\newcolumntype{L}{>{\raggedright\arraybackslash}X}
\newcolumntype{R}{>{\raggedleft\arraybackslash}X}
\usepackage{array}
\usepackage{multicol}
\usepackage{siunitx}
% parse-numbers=false : accepte aussi \SI{2,5 \times 10^{-3}}{...}
\sisetup{locale=FR,output-decimal-marker={,},group-minimum-digits=4,parse-numbers=false}
\DeclareSIUnit\ohm{\ensuremath{\symup{\Omega}}} % U+2126 absent de la police
\DeclareSIUnit\division{div} % réglages d'oscilloscope (ms/div, V/div)
\DeclareSIUnit\tr{tr} % tours (vitesse de rotation, ex. \SI{1500}{\tr\per\minute})
\DeclareSIUnit\molar{M} % concentration molaire (ex. \SI{3}{\molar}, \SI{1}{\micro\molar})
\DeclareSIUnit\an{an} % année (le modèle confond parfois avec \year, un compteur TeX) — ex. \SI{5}{\kilo\meter\per\an}
\DeclareSIUnit\FCFA{FCFA} % franc CFA (ex. \SI{500}{\FCFA\per\kilo\gram})
\DeclareSIUnit\degreeBrix{\textdegree Brix} % teneur en sucre d'un sirop/jus (réfractométrie)
\DeclareSIUnit\month{mois} % le modèle utilise parfois \month, un compteur TeX, comme unité de durée
\DeclareSIUnit\microliter{\micro\liter} % le modèle écrit parfois \microliter en un seul token
\DeclareSIUnit\million{millions} % dénombrement (ex. \SI{15}{\million\per\mL} spermatozoïdes)
\usepackage{qrcode}
% \degree : commande fréquemment utilisée par les modèles pour un angle ou une
% température en dehors de \SI{}{} (non fournie par les paquets chargés).
\providecommand{\degree}{^{\circ}}
% \qed (fin de démonstration), utilisé par les modèles sans amsthm
\providecommand{\qed}{\hfill$\square$}
% \barre : double barre des valeurs interdites dans un tableau de variations (tkz-tab)
\providecommand{\barre}{\ensuremath{\|}}
% environnement proof (amsthm non chargé) utilisé par les modèles
\ifdefined\proof\else\newenvironment{proof}[1][Démonstration]{\par\noindent\textit{#1.}\ }{\qed\par}\fi
\usepackage{lastpage}
\usepackage{hyperref}
\hypersetup{hidelinks}

% ============================================================
% Palette « Pop » OnBuch+ — mêmes hex que la classe P (plus_ui.dart)
% ============================================================
\definecolor{poporange}{HTML}{F59321}
\definecolor{popdark}{HTML}{E07A0C}
\definecolor{popcream}{HTML}{FFF6E8}
\definecolor{popink}{HTML}{1C1714}
\definecolor{poppurple}{HTML}{7A5AE0}
\definecolor{popgreen}{HTML}{2E9E6A}
\definecolor{poppink}{HTML}{E0457B}
\definecolor{popblue}{HTML}{2F7DE1}
\definecolor{popgold}{HTML}{C98A12}
\definecolor{popmuted}{HTML}{8A7F76}
\definecolor{popline}{HTML}{EADFD2}
\definecolor{popgray}{HTML}{8A7F76}
\colorlet{popgrey}{popgray}
% Alias tolérants (le modèle nomme parfois les couleurs différemment)
\colorlet{popwhite}{white}
\colorlet{popblack}{popink}
\colorlet{popred}{poppink}
\colorlet{popyellow}{popgold}
% Teintes claires (fonds de tableaux/encadrés) — le modèle nomme parfois
% les couleurs avec un suffixe L (ex. popblueL) sans les avoir définies.
\colorlet{poporangeL}{poporange!20}
\colorlet{popdarkL}{popdark!20}
\colorlet{poppurpleL}{poppurple!20}
\colorlet{popgreenL}{popgreen!20}
\colorlet{poppinkL}{poppink!20}
\colorlet{popblueL}{popblue!20}
\colorlet{popgoldL}{popgold!20}
\colorlet{popredL}{popred!20}
\colorlet{popgrayL}{popgray!20}
% Teintes claires pour fonds de tableaux / zones
\colorlet{popblueL}{popblue!10!white}
\colorlet{poporangeL}{poporange!14!white}
\colorlet{popgreenL}{popgreen!12!white}
\colorlet{poppurpleL}{poppurple!12!white}
\colorlet{poppinkL}{poppink!12!white}
\colorlet{popgoldL}{popgold!14!white}

\pagecolor{popcream}

% ============================================================
% Typographie
% ============================================================
\defaultfontfeatures{Ligatures=TeX}
\setmainfont{PlusJakartaSans-Medium.ttf}[Path=../../../fonts/,BoldFont=PlusJakartaSans-Bold.ttf]
% Maths en sans-serif (Fira Math) avec chiffres et lettres droites de Plus
% Jakarta, pour que les formules se fondent dans le texte.
\usepackage[math-style=ISO,bold-style=ISO]{unicode-math}
\setmathfont{FiraMath-Regular.otf}[Path=../../../fonts/]
\setmathfont{PlusJakartaSans-Medium.ttf}[Path=../../../fonts/,range={up/num,up/{Latin,latin},\mathup}]
\usepackage{icomma} % virgule décimale sans espace en mode math
% Caractères mathématiques Unicode absents des polices de texte (Plus Jakarta,
% Archivo Black) : rendus par leur équivalent mathématique, en texte comme en
% formule — sinon ils s'affichent en carré vide (ex. « Arithmétique dans ℤ »).
\usepackage{newunicodechar}
\newunicodechar{ℝ}{\ensuremath{\mathbb{R}}}
\newunicodechar{ℤ}{\ensuremath{\mathbb{Z}}}
\newunicodechar{ℂ}{\ensuremath{\mathbb{C}}}
\newunicodechar{ℕ}{\ensuremath{\mathbb{N}}}
\newunicodechar{ℚ}{\ensuremath{\mathbb{Q}}}
\newunicodechar{𝔻}{\ensuremath{\mathbb{D}}}
\newunicodechar{α}{\ensuremath{\alpha}}
\newunicodechar{β}{\ensuremath{\beta}}
\newunicodechar{γ}{\ensuremath{\gamma}}
\newunicodechar{δ}{\ensuremath{\delta}}
\newunicodechar{ε}{\ensuremath{\varepsilon}}
\newunicodechar{θ}{\ensuremath{\theta}}
\newunicodechar{λ}{\ensuremath{\lambda}}
\newunicodechar{μ}{\ensuremath{\mu}}
\newunicodechar{σ}{\ensuremath{\sigma}}
\newunicodechar{τ}{\ensuremath{\tau}}
\newunicodechar{φ}{\ensuremath{\varphi}}
\newunicodechar{ψ}{\ensuremath{\psi}}
\newunicodechar{ω}{\ensuremath{\omega}}
\newunicodechar{Σ}{\ensuremath{\Sigma}}
% majuscules produites par \MakeUppercase dans les titres (α-aminés -> Α)
\newunicodechar{Α}{\ensuremath{\alpha}}
\newunicodechar{Β}{\ensuremath{\beta}}
\newunicodechar{Γ}{\ensuremath{\Gamma}}
\newunicodechar{Δ}{\ensuremath{\Delta}}
\newunicodechar{Ε}{\ensuremath{\varepsilon}}
\newunicodechar{Θ}{\ensuremath{\Theta}}
\newunicodechar{Λ}{\ensuremath{\Lambda}}
\newunicodechar{Μ}{\ensuremath{\mu}}
\newunicodechar{Τ}{\ensuremath{\tau}}
\newunicodechar{Φ}{\ensuremath{\Phi}}
\newunicodechar{Ψ}{\ensuremath{\Psi}}
\newunicodechar{Ω}{\ensuremath{\Omega}}
\newunicodechar{↦}{\ensuremath{\mapsto}}
\newunicodechar{∈}{\ensuremath{\in}}
\newunicodechar{∉}{\ensuremath{\notin}}
\newunicodechar{∧}{\ensuremath{\wedge}}
\newunicodechar{⇔}{\ensuremath{\Leftrightarrow}}
\newunicodechar{⟺}{\ensuremath{\Longleftrightarrow}}
\newunicodechar{⇒}{\ensuremath{\Rightarrow}}
\newunicodechar{⟹}{\ensuremath{\Longrightarrow}}
\newunicodechar{∘}{\ensuremath{\circ}}
\newunicodechar{∪}{\ensuremath{\cup}}
\newunicodechar{∩}{\ensuremath{\cap}}
\newunicodechar{≡}{\ensuremath{\equiv}}
\newunicodechar{⊂}{\ensuremath{\subset}}
\newunicodechar{⊥}{\ensuremath{\perp}}
\newunicodechar{∃}{\ensuremath{\exists}}
\newunicodechar{∀}{\ensuremath{\forall}}
\newunicodechar{∛}{\ensuremath{\sqrt[3]{\,}}}
\newunicodechar{∜}{\ensuremath{\sqrt[4]{\,}}}
\newunicodechar{⃗}{\ensuremath{{}^{\to}}}
\newunicodechar{ⁿ}{\ifmmode^{n}\else\textsuperscript{\ensuremath{n}}\fi}
\newunicodechar{ˣ}{\ifmmode^{x}\else\textsuperscript{\ensuremath{x}}\fi}
\newunicodechar{ᵇ}{\ifmmode^{b}\else\textsuperscript{\ensuremath{b}}\fi}
\newunicodechar{ʳ}{\ifmmode^{r}\else\textsuperscript{\ensuremath{r}}\fi}
\newunicodechar{ᵘ}{\ifmmode^{u}\else\textsuperscript{\ensuremath{u}}\fi}
\newunicodechar{ʸ}{\ifmmode^{y}\else\textsuperscript{\ensuremath{y}}\fi}
\newunicodechar{ᵃ}{\ifmmode^{a}\else\textsuperscript{\ensuremath{a}}\fi}
\newunicodechar{ᵉ}{\ifmmode^{e}\else\textsuperscript{\ensuremath{e}}\fi}
\newunicodechar{ᵏ}{\ifmmode^{k}\else\textsuperscript{\ensuremath{k}}\fi}
\newunicodechar{ᵖ}{\ifmmode^{p}\else\textsuperscript{\ensuremath{p}}\fi}
\newunicodechar{ᴺ}{\ifmmode^{N}\else\textsuperscript{\ensuremath{N}}\fi}
\newunicodechar{ᶜ}{\ifmmode^{c}\else\textsuperscript{\ensuremath{c}}\fi}
\newunicodechar{ʰ}{\ifmmode^{h}\else\textsuperscript{\ensuremath{h}}\fi}
\newunicodechar{ᵠ}{\ifmmode^{q}\else\textsuperscript{\ensuremath{q}}\fi}
\newunicodechar{ₙ}{\ifmmode_{n}\else\textsubscript{\ensuremath{n}}\fi}
\newunicodechar{ₐ}{\ifmmode_{a}\else\textsubscript{\ensuremath{a}}\fi}
\newunicodechar{ᵢ}{\ifmmode_{i}\else\textsubscript{\ensuremath{i}}\fi}
\newunicodechar{ᵣ}{\ifmmode_{r}\else\textsubscript{\ensuremath{r}}\fi}
\newunicodechar{ₖ}{\ifmmode_{k}\else\textsubscript{\ensuremath{k}}\fi}
\newunicodechar{ᵦ}{\ifmmode_{\beta}\else\textsubscript{\ensuremath{\beta}}\fi}
\newunicodechar{ₓ}{\ifmmode_{x}\else\textsubscript{\ensuremath{x}}\fi}
\newfontfamily\popdisplay{ArchivoBlack-Regular.ttf}[Path=../../../fonts/]
\newfontfamily\poplogo{SpaceGrotesk-Bold.ttf}[Path=../../../fonts/]
\newfontfamily\popsemi{PlusJakartaSans-SemiBold.ttf}[Path=../../../fonts/]
\newfontfamily\popxbold{PlusJakartaSans-ExtraBold.ttf}[Path=../../../fonts/]
\newfontfamily\popxboldls{PlusJakartaSans-ExtraBold.ttf}[Path=../../../fonts/,LetterSpace=3.0]

\setlist[enumerate]{leftmargin=1.7em,itemsep=5pt,topsep=4pt}
\setlist[itemize]{leftmargin=1.7em,itemsep=4pt,topsep=4pt,label={\color{poporange}\textbullet}}
\setlength{\parindent}{0pt}
\setlength{\parskip}{5pt}
\renewcommand{\arraystretch}{1.25}

% pgfplots : style Pop par défaut
\pgfplotsset{
  every axis/.append style={axis line style={popink,line width=1pt},tick style={popink},
    grid style={popline,line width=0.5pt},label style={font=\small},tick label style={font=\footnotesize}},
  popcurve/.style={poporange,line width=1.8pt,no markers,smooth},
}
% Même style disponible dans un \draw TikZ simple (hors pgfplots)
\tikzset{popcurve/.style={poporange,line width=1.8pt}}
\tikzset{popgrid/.style={popline,very thin}}
\colorlet{popcurve}{poporange} % le nom du style est parfois utilisé comme couleur

% ============================================================
% Pill / squiggle
% ============================================================
\newcommand{\poppill}[3]{%
  \tikz[baseline=-0.55ex]\node[draw=popink,line width=1.2pt,rounded corners=2.6pt,%
    fill=#2,inner xsep=6.5pt,inner ysep=3pt]{%
    \popxboldls\fontsize{7}{7}\selectfont\color{#3}\MakeUppercase{#1}};%
}
\newcommand{\popsquiggle}[1]{%
  \tikz[baseline=-0.4ex]{
    \draw[#1,line width=2.6pt,line cap=round]
      plot [smooth] coordinates {(0,0.09) (0.34,0.01) (0.68,0.21) (1.02,0.01) (1.36,0.10)};
  }%
}
% Mot-clé surligné (marqueur orange sous le texte)
\newcommand{\cle}[1]{{\popxbold\color{popdark}#1}}

% ============================================================
% En-tête / pied
% ============================================================
\pagestyle{fancy}
\fancyhf{}
\renewcommand{\headrulewidth}{0pt}
\renewcommand{\footrulewidth}{0pt}
\lhead{\raisebox{-0.15ex}{\poplogo\fontsize{13}{13}\selectfont\color{popink}OnBuch\color{poporange}+}}
\rhead{\poppill{\DOCMATIERE\ \textbullet\ \DOCNIVEAU\ \textbullet\ Leçon \DOCLECON}{white}{popink}}
\lfoot{\popxbold\fontsize{7.6}{7.6}\selectfont\color{popmuted}\MakeUppercase{Cours signé OnBuch+}\ \textbullet\ {\popsemi\DOCTITRE}}
\rfoot{\tikz[baseline=-0.6ex]\node[rounded corners=8.5pt,fill=popink,minimum height=17pt,minimum width=17pt,inner xsep=5pt,inner sep=0]{\popxbold\fontsize{8}{8}\selectfont\color{popcream}\thepage\,/\,\pageref*{LastPage}};}

% ============================================================
% Titres
% ============================================================
\newcommand{\chaptitre}[2]{%
  \begin{tcolorbox}[enhanced,colback=poporange,colframe=popink,boxrule=2.6pt,arc=11pt,
      drop shadow=popink,left=15pt,right=15pt,top=12pt,bottom=11pt,before skip=3pt,after skip=18pt]
    \poppill{#2}{popink}{popcream}\par\vspace{9pt}
    {\raggedright\hyphenpenalty=10000\exhyphenpenalty=10000\popdisplay\fontsize{22}{24}\selectfont\color{popcream}\MakeUppercase{#1}\par}
  \end{tcolorbox}
}
% Section numérotée : pastille orange avec le numéro + titre Archivo
\newcounter{popsec}
\newcommand{\coursec}[1]{%
  \stepcounter{popsec}\setcounter{popsub}{0}%
  \par\vspace{16pt}\noindent
  \begin{minipage}{\linewidth}
  \tikz[baseline=-0.7ex]\node[draw=popink,line width=1.6pt,fill=poporange,rounded corners=4pt,minimum size=22pt,inner sep=0]{\popdisplay\fontsize{12}{12}\selectfont\color{popcream}\thepopsec};%
  \hspace{8pt}{\popdisplay\fontsize{15}{17}\selectfont\color{popink}\MakeUppercase{#1}}\par
  \vspace{4pt}\noindent{\color{poporange}\rule{36pt}{3.6pt}}
  \end{minipage}\par\nopagebreak\vspace{8pt}
}
\newcounter{popsub}
\newcommand{\courssub}[1]{%
  \stepcounter{popsub}%
  \par\vspace{9pt}\noindent{\popxbold\fontsize{11.5}{13}\selectfont\color{popink}{\color{poporange}\thepopsec.\thepopsub}\hspace{6pt}#1}\par\nopagebreak\vspace{4pt}
}

% ============================================================
% Boîtes pédagogiques — même recette : fond blanc, bordure épaisse colorée,
% ombre dure encre, badge plein. Titre optionnel [#1].
% ============================================================
\tcbset{popbase/.style={breakable,enhanced,colback=white,boxrule=2.3pt,arc=9pt,
  shadow={3pt}{-3pt}{0pt}{popink},left=14pt,right=14pt,top=11pt,bottom=11pt,before skip=11pt,after skip=15pt,
  fontupper=\popsemi\fontsize{10.6}{14.5}\selectfont\color{popink},
  fontlower=\popsemi\fontsize{10.6}{14.5}\selectfont\color{popink}}}
% Coche dessinée (le glyphe ✓ n'existe pas dans les polices du document)
\newcommand{\popcheck}[1]{\tikz[baseline=-0.5ex]\draw[#1,line width=1.6pt,line cap=round,line join=round](-0.13,0.0)--(-0.04,-0.09)--(0.13,0.1);}
\providecommand{\coche}{\popcheck{popgreen}}
\providecommand{\resultat}[1]{\par\smallskip\noindent\fcolorbox{popgreen}{popgreenL}{\parbox{\dimexpr\linewidth-2\fboxsep-2\fboxrule\relax}{\textbf{Résultat.} #1}}\par\smallskip}
\newcommand{\popboxhead}[3]{\poppill{#1}{#2}{white}\ifx\relax#3\relax\else\hspace{6pt}{\popxbold\fontsize{10.6}{12}\selectfont\color{popink}#3}\fi\par\vspace{7pt}}

\newtcolorbox{aretenir}[1][]{popbase,colframe=popgreen,before upper={\popboxhead{À retenir}{popgreen}{#1}}}
\newtcolorbox{exemplebox}[1][]{popbase,colframe=poporange,before upper={\popboxhead{Exemple}{poporange}{#1}}}
\newtcolorbox{definition}[1][]{popbase,colframe=popblue,before upper={\popboxhead{Définition}{popblue}{#1}}}
\newtcolorbox{propriete}[1][]{popbase,colframe=poppurple,before upper={\popboxhead{Propriété}{poppurple}{#1}}}
\newtcolorbox{methode}[1][]{popbase,colframe=popgold,before upper={\popboxhead{Méthode}{popgold}{#1}}}
\newtcolorbox{attention}[1][]{popbase,colframe=poppink,before upper={\popboxhead{Attention}{poppink}{#1}}}
\newtcolorbox{experience}[1][]{popbase,colframe=popblue,colback=popblueL,before upper={\popboxhead{Expérience}{popblue}{#1}}}
% « Le savais-tu ? » : carte violette pleine (contexte, culture, Cameroun)
\newtcolorbox{savaistu}[1][]{popbase,colback=poppurple,colframe=popink,
  fontupper=\popsemi\fontsize{10.4}{14.2}\selectfont\color{white},
  before upper={\poppill{Le savais-tu\,?}{white}{poppurple}\ifx\relax#1\relax\else\hspace{6pt}{\popxbold\color{white}#1}\fi\par\vspace{7pt}}}
% Exercice résolu : énoncé en haut, \tcblower puis la solution sur fond crème
\newcounter{popexo}\newcounter{popexr}
\newtcolorbox{exoresolu}[1][]{popbase,colframe=poporange,colbacklower=popcream,
  segmentation style={popink,line width=1.2pt,dashed},
  before upper={\stepcounter{popexr}\popboxhead{Exercice résolu \thepopexr}{poporange}{#1}},
  before lower={\poppill{Solution}{popink}{popcream}\par\vspace{6pt}}}
% Exercice (non résolu en place) — la correction est dans la partie Corrigés
\newtcolorbox{exercice}[1][]{popbase,colframe=popink,
  before upper={\stepcounter{popexo}\popboxhead{Exercice \thepopexo}{popink}{#1}}}
% Corrigé
\newtcolorbox{corrige}[1][]{popbase,colframe=popgreen,colback=popgreenL,
  before upper={\popboxhead{Corrigé}{popgreen}{#1}}}
% Objectifs / prérequis / bilan
\newtcolorbox{objectifs}{popbase,colframe=popink,colback=poporangeL,
  before upper={\popboxhead{Objectifs de la leçon}{popink}{}}}
\newtcolorbox{prerequis}{popbase,colframe=popink,colback=popblueL,
  before upper={\popboxhead{Prérequis}{popblue}{}}}
\newtcolorbox{bilan}{popbase,colframe=popink,colback=white,boxrule=2.8pt,
  before upper={\popboxhead{Fiche bilan}{poporange}{L'essentiel à connaître pour l'examen}}}
% Figure encadrée avec légende
\newtcolorbox{popfigure}[1][]{enhanced,breakable=false,colback=white,colframe=popline,boxrule=1.4pt,arc=8pt,
  left=10pt,right=10pt,top=10pt,bottom=8pt,before skip=10pt,after skip=12pt,halign=center,
  lower separated=false,halign lower=center,
  fontlower=\popsemi\fontsize{9}{11.5}\selectfont\color{popmuted},
  #1}
\newcommand{\legende}[1]{\par\vspace{6pt}{\popsemi\fontsize{9}{11.5}\selectfont\color{popmuted}#1}}

% ============================================================
% Signature OnBuch+
% ============================================================
\newcommand{\poplogoinline}[1]{{\poplogo\fontsize{#1}{#1}\selectfont\color{popink}OnBuch\color{poporange}+}}
\newcommand{\signatureonbuch}{%
  \begin{tcolorbox}[enhanced,colback=popink,colframe=popink,boxrule=2.2pt,arc=10pt,
      drop shadow=poporange,left=14pt,right=14pt,top=10pt,bottom=10pt,before skip=0pt,after skip=0pt]
    \tikz[baseline=-0.7ex]\node[circle,fill=popgreen,draw=popcream,line width=1.3pt,minimum size=22pt,inner sep=0]{\popcheck{white}};%
    \hspace{9pt}\begin{minipage}[c]{0.62\linewidth}
      {\popxbold\fontsize{12}{13}\selectfont\color{popcream}Signé\ \textbullet\ {\poplogo OnBuch\color{poporange}+}}\par\vspace{2pt}
      {\raggedright\popsemi\fontsize{8}{10.5}\selectfont\color{popline}\MakeUppercase{Équipe pédagogique OnBuch+}\\\MakeUppercase{Conforme au programme officiel MINESEC}\par}
    \end{minipage}\hfill
    {\poplogo\fontsize{22}{22}\selectfont\color{popcream}OnBuch\color{poporange}+}
  \end{tcolorbox}%
}

% ============================================================
% Page de garde
% #1 titre · #2 matière · #3 niveau · #4 module · #5 n° leçon · #6 durée ·
% #7 description / objectifs (texte ou itemize)
% ============================================================
\newcommand{\pagegarde}[7]{%
  \thispagestyle{empty}
  \newgeometry{a4paper,margin=1.7cm,top=1.6cm,bottom=1.4cm}%
  \begin{tikzpicture}[remember picture,overlay]
    \fill[poporange!18] ([xshift=-3.2cm,yshift=-3.6cm]current page.north east) circle (4.6cm);
    \fill[poppurple!14] ([xshift=2.4cm,yshift=7.6cm]current page.south west) circle (3.8cm);
  \end{tikzpicture}%
  {\poplogo\fontsize{30}{30}\selectfont\color{popink}OnBuch\color{poporange}+}\hfill
  \raisebox{0.6ex}{\poppill{Cours signé \textbullet\ Premium}{popink}{popcream}}\par
  \vspace{1pt}\popsquiggle{poppurple}\par
  \vspace{8pt}
  \begin{tcolorbox}[enhanced,colback=poporange,colframe=popink,boxrule=2.8pt,arc=13pt,
      drop shadow=popink,left=18pt,right=18pt,top=12pt,bottom=13pt,after skip=10pt]
    \poppill{#2 \textbullet\ #3}{popink}{popcream}\hspace{5pt}\poppill{#4}{white}{popink}\par\vspace{8pt}
    {\popdisplay\fontsize{14}{14}\selectfont\color{popink}LEÇON #5}\par\vspace{4pt}
    {\raggedright\hyphenpenalty=10000\exhyphenpenalty=10000\popdisplay\fontsize{25}{27}\selectfont\color{popcream}\MakeUppercase{#1}\par}
    \vspace{8pt}
    {\popxbold\fontsize{9.5}{9.5}\selectfont\color{popink}\MakeUppercase{Durée conseillée : #6}}
  \end{tcolorbox}
  \begin{tcolorbox}[enhanced,colback=white,colframe=popink,boxrule=2.2pt,arc=10pt,
      drop shadow=popink,left=16pt,right=16pt,top=10pt,bottom=10pt,after skip=8pt]
    \poppill{Dans cette leçon}{poppurple}{white}\par\vspace{5pt}
    \popsemi\fontsize{9.3}{12.6}\selectfont\color{popink}#7
  \end{tcolorbox}
  \noindent\begin{minipage}[t]{0.487\linewidth}
    \begin{tcolorbox}[enhanced,colback=popgreen,colframe=popink,boxrule=2.2pt,arc=10pt,
        drop shadow=popink,left=11pt,right=11pt,top=10pt,bottom=10pt,height=3.1cm,valign=center]
      \tcbox[colback=white,colframe=popink,boxrule=1.3pt,arc=5pt,boxsep=0pt,nobeforeafter,
        left=3pt,right=3pt,top=3pt,bottom=3pt,valign=center]{\qrcode[height=1.9cm]{https://play.google.com/store/apps/details?id=com.luvvix.onbuch&hl=fr}}%
      \hspace{8pt}\begin{minipage}[c]{0.5\linewidth}
        {\popdisplay\fontsize{11}{12.5}\selectfont\color{white}\MakeUppercase{Télécharge OnBuch}}\par\vspace{4pt}
        {\raggedright\popsemi\fontsize{8.4}{11}\selectfont\color{white}Gratuit \textbullet\ Google Play\\Tuteur IA, annales, concours, résultats\par}
      \end{minipage}
    \end{tcolorbox}
  \end{minipage}\hfill
  \begin{minipage}[t]{0.487\linewidth}
    \begin{tcolorbox}[enhanced,colback=poppurple,colframe=popink,boxrule=2.2pt,arc=10pt,
        drop shadow=popink,left=11pt,right=11pt,top=10pt,bottom=10pt,height=3.1cm,valign=center]
      \tikz[baseline=-0.7ex]\node[circle,fill=white,draw=popink,line width=1.3pt,minimum size=1.6cm,inner sep=0]{\popdisplay\fontsize{24}{24}\selectfont\color{poppurple}+};%
      \hspace{8pt}\begin{minipage}[c]{0.6\linewidth}
        {\popdisplay\fontsize{11}{12.5}\selectfont\color{white}\MakeUppercase{Passe à OnBuch+}}\par\vspace{4pt}
        {\raggedright\popsemi\fontsize{8.4}{11}\selectfont\color{white}Tous les cours signés, Tuteur IA illimité, fiches PDF — onglet OnBuch+ de l'app\par}
      \end{minipage}
    \end{tcolorbox}
  \end{minipage}
  \vspace{8pt}
  \popcontenu
  \vfill
  \signatureonbuch
  \restoregeometry
  \newpage
}

% Rangée « Ce cours contient » (page de garde)
\newcommand{\poptile}[3]{% #1 couleur #2 gros texte #3 légende
  \begin{tcolorbox}[enhanced,colback=#1,colframe=popink,boxrule=2pt,arc=9pt,shadow={2.5pt}{-2.5pt}{0pt}{popink},
      left=6pt,right=6pt,top=9pt,bottom=9pt,halign=center,valign=center,height=1.75cm,width=\linewidth,nobeforeafter]
    {\popdisplay\fontsize{12.5}{13}\selectfont\color{white}\MakeUppercase{#2}}\par\vspace{3pt}
    {\centering\popsemi\fontsize{7.8}{9.5}\selectfont\color{white}#3\par}
  \end{tcolorbox}}
\newcommand{\popcontenu}{%
  \par\noindent{\popxboldls\fontsize{7.5}{7.5}\selectfont\color{popmuted}CE COURS SIGNÉ CONTIENT}\par\vspace{7pt}
  \noindent\begin{minipage}[t]{0.235\linewidth}\poptile{popblue}{Cours}{complet, illustré, pas à pas}\end{minipage}\hfill
  \begin{minipage}[t]{0.235\linewidth}\poptile{poporange}{Méthodes}{et exercices résolus}\end{minipage}\hfill
  \begin{minipage}[t]{0.235\linewidth}\poptile{poppink}{Exercices}{type examen + corrigés}\end{minipage}\hfill
  \begin{minipage}[t]{0.235\linewidth}\poptile{popgreen}{Fiche bilan}{l'essentiel à retenir}\end{minipage}\par}

% Sommaire
\newtcolorbox{sommaire}{enhanced,colback=white,colframe=popink,boxrule=2.2pt,arc=10pt,
  drop shadow=popink,left=18pt,right=18pt,top=13pt,bottom=13pt,before skip=6pt,after skip=20pt,
  before upper={\poppill{Sommaire}{poppurple}{white}\par\vspace{9pt}\popsemi\fontsize{10.6}{14.5}\selectfont\color{popink}}}

% Signature de fin de cours — poussée tout en bas de la dernière page
\newcommand{\findecours}{%
  \par\vfill\nopagebreak
  \begin{center}\popsquiggle{poporange}\end{center}\vspace{4pt}
  \signatureonbuch
}
\AtBeginDocument{\def\llbracket{\mathopen{[\mkern-3.5mu[}}\def\rrbracket{\mathclose{]\mkern-3.5mu]}}} % intervalles d'entiers (glyphe ⟦ absent de la police)
% Glyphes absents de Fira Math : redéfinis à partir de glyphes disponibles
\AtBeginDocument{%
  \def\setminus{\mathbin{\backslash}}\let\smallsetminus\setminus
  \def\vdots{\mathord{\vbox{\baselineskip3.2pt\lineskiplimit0pt\kern4pt\hbox{.}\hbox{.}\hbox{.}}}}%
  \def\ddots{\mathinner{\mkern1mu\raise6.5pt\hbox{.}\mkern2mu\raise3.6pt\hbox{.}\mkern2mu\raise0.7pt\hbox{.}\mkern1mu}}%
  \def\triangle{\mathord{\scalebox{0.9}{$\symup{\Delta}$}}}%
  \def\nabla{\mathord{\raisebox{\depth}{\scalebox{1}[-1]{$\symup{\Delta}$}}}}%
  \def\checkmark{\popcheck{popdark}}%
  \def\star{\mathbin{\ast}}\def\bigstar{\mathord{\ast}}%
  \def\diamond{\mathbin{\scalebox{0.6}{$\Diamond$}}}%
  \def\bigcup{\mathop{\mathchoice{\raisebox{-0.3em}{\scalebox{1.7}{$\cup$}}}{\scalebox{1.25}{$\cup$}}{\cup}{\cup}}\displaylimits}%
  \def\bigcap{\mathop{\mathchoice{\raisebox{-0.3em}{\scalebox{1.7}{$\cap$}}}{\scalebox{1.25}{$\cap$}}{\cap}{\cap}}\displaylimits}%
  \def\lesseqqgtr{\mathrel{\lessgtr}}\def\bigoplus{\mathop{\oplus}\displaylimits}%
}
\providecommand{\ssi}{si et seulement si} % abréviation fréquente
\AtBeginDocument{\providecommand{\wideparen}{\overparen}} % arc de cercle

%% FIGFIX-BEGIN v5 — correctifs globaux des figures (docs/onbuch-generation/tools/figfix)
\usepackage{adjustbox}\usepackage{etoolbox}\tcbuselibrary{raster}
% toute figure TikZ/pgfplots plus large que la ligne est réduite à la largeur disponible
\BeforeBeginEnvironment{tikzpicture}{\begin{adjustbox}{max width=\linewidth}}
\AfterEndEnvironment{tikzpicture}{\end{adjustbox}}
% légendes pgfplots lisibles par-dessus les courbes
\pgfplotsset{every axis/.append style={legend style={fill=white,fill opacity=0.92,text opacity=1,draw=black!15,inner sep=3pt}}}
% étiquettes d'arêtes (syntaxe edge["..."]) avec fond blanc
\tikzset{every edge quotes/.append style={fill=white,inner sep=1.2pt}}
%% FIGFIX-END
\begin{document}

\pagegarde{L'exercice sur texte}{Philosophie}{1re Tech}{Philosophie – classe de Première technique}{N21}{3 séances (fiche de progression harmonisée)}{Cette leçon présente la méthode complète de l'exercice sur texte philosophique : lecture analytique, réponse aux questions de compréhension et rédaction de l'essai personnel. Elle outille l'élève de Première technique pour réussir cette épreuve incontournable du Probatoire et du Baccalauréat.
\vspace{4pt}
\begin{multicols}{2}\small
\begin{itemize}
\item À la fin de cette leçon, je sais définir l'exercice sur texte et distinguer ses trois étapes (généralités, réponses aux questions, essai personnel)
\item À la fin de cette leçon, je sais dégager le thème, la thèse et la problématique d'un texte philosophique
\item À la fin de cette leçon, je sais repérer la structure argumentative d'un texte (arguments, exemples, articulations logiques)
\item À la fin de cette leçon, je sais rédiger des réponses précises et justifiées aux questions de compréhension
\item À la fin de cette leçon, je sais construire un essai personnel structuré (introduction, développement, conclusion)
\item À la fin de cette leçon, je sais appliquer les notions philosophiques étudiées à des situations concrètes de la vie courante camerounaise
\end{itemize}\end{multicols}}

\begin{sommaire}
\begin{multicols}{2}
\begin{enumerate}
\item Généralités sur l'exercice sur texte
\item La lecture analytique du texte
\item Les réponses aux questions
\item L'essai personnel
\item Application guidée : un sujet complet corrigé
\item Activité d'intégration
\item Méthodes et exercices résolus
\item Exercices
\item Corrigés des exercices
\item Fiche bilan
\end{enumerate}
\end{multicols}
\end{sommaire}

\begin{objectifs}
\begin{itemize}
\item À la fin de cette leçon, je sais définir l'exercice sur texte et distinguer ses trois étapes (généralités, réponses aux questions, essai personnel)
\item À la fin de cette leçon, je sais dégager le thème, la thèse et la problématique d'un texte philosophique
\item À la fin de cette leçon, je sais repérer la structure argumentative d'un texte (arguments, exemples, articulations logiques)
\item À la fin de cette leçon, je sais rédiger des réponses précises et justifiées aux questions de compréhension
\item À la fin de cette leçon, je sais construire un essai personnel structuré (introduction, développement, conclusion)
\item À la fin de cette leçon, je sais appliquer les notions philosophiques étudiées à des situations concrètes de la vie courante camerounaise
\item À la fin de cette leçon, je sais justifier ma démarche, mes réponses et mes conclusions par des arguments et des références
\item À la fin de cette leçon, je sais éviter les erreurs fréquentes (paraphrase, hors-sujet, copier-coller du texte)
\end{itemize}
\end{objectifs}

\begin{prerequis}
\begin{itemize}
\item Notions de l'unité en cours (ex. : la conscience, la liberté, le devoir) étudiées en classe de Première
\item La problématique et la thèse : notions introduites dans les leçons précédentes
\item Techniques de lecture rapide et de prise de notes acquises en Seconde
\item Connaissance des grands auteurs et courants philosophiques du programme
\item Règles de base de la rédaction d'un devoir de philosophie (introduction, développement, conclusion)
\end{itemize}
\end{prerequis}

\newpage

\coursec{Généralités sur l'exercice sur texte}

Au lycée de Nkongsamba, Amina, élève de Première technique, reçoit au Bac blanc un texte de Kant sur la liberté. Elle le lit trois fois, mais recopie simplement des phrases du texte dans sa copie : elle obtient 6/20. Son voisin, qui a pourtant moins de connaissances, applique une méthode rigoureuse et décroche 14/20. Qu'est-ce qui fait la différence ? Ni l'intelligence, ni la mémoire, mais la \cle{méthode}. Comment transformer la lecture d'un texte philosophique en une copie structurée et valorisée ? Cette leçon donne les clés de l'exercice sur texte, épreuve reine du Probatoire et du Baccalauréat techniques.

\textbf{Problématique :} qu'est-ce que l'exercice sur texte, qu'attend exactement le correcteur, et comment organiser son temps et sa copie pour réussir cette épreuve ?

\courssub{Définition et nature de l'épreuve}

Commençons par une intuition simple. Quand tu lis un message sur ton téléphone, tu peux te contenter d'en saisir le sens global. Mais si l'on te demande ensuite : « Que veut dire précisément cette phrase ? Es-tu d'accord avec son auteur ? Pourquoi ? », tu dois passer d'une lecture passive à une lecture \emph{active} et \emph{argumentée}. L'exercice sur texte est exactement cela, appliqué à un extrait d'œuvre philosophique : on ne te demande pas de réciter un cours, mais de \textbf{comprendre} un texte, de le \textbf{questionner} et de \textbf{penser par toi-même} à partir de lui.

\begin{definition}[L'exercice sur texte]
L'\cle{exercice sur texte} est une épreuve consistant à \textbf{lire} un texte philosophique (un extrait d'œuvre d'un auteur au programme), à \textbf{répondre à des questions} portant sur sa compréhension (sens des termes, thèse de l'auteur, arguments, structure), puis à \textbf{rédiger un essai personnel} sur le problème soulevé par le texte. C'est un \emph{commentaire guidé} : le candidat n'est pas livré à lui-même, le sujet l'accompagne pas à pas.
\end{definition}

Trois traits distinguent donc cette épreuve :
\begin{itemize}
\item Elle part toujours d'un \cle{texte} précis, signé par un auteur (Platon, Descartes, Kant, Nietzsche, Senghor, etc.), tiré d'une œuvre identifiée.
\item Elle est \cle{guidée} : les questions indiquent le chemin à suivre ; elles sont une aide, non un piège.
\item Elle culmine dans un \cle{essai personnel} où le candidat prend position de manière argumentée sur le problème du texte.
\end{itemize}

\begin{exemplebox}[Un sujet type]
« Lisez le texte suivant de Kant, extrait de \emph{Qu'est-ce que les Lumières ?}, puis répondez aux questions : 1) Que signifie "minorité" dans ce texte ? 2) Quelle est la thèse de l'auteur ? 3) Dégager l'intérêt philosophique du texte. \textbf{Essai personnel :} La liberté de penser est-elle sans limites ? »
\end{exemplebox}

\courssub{Place de l'épreuve au Probatoire et au Baccalauréat techniques}

L'exercice sur texte n'est pas un exercice scolaire parmi d'autres : c'est l'une des épreuves \textbf{officielles} de philosophie au Probatoire et au Baccalauréat techniques, au même titre que la dissertation.

\begin{aretenir}[Repères officiels]
\begin{itemize}
\item \textbf{Durée de l'épreuve de philosophie :} 3 heures (au choix : une dissertation \emph{ou} un exercice sur texte).
\item \textbf{Coefficient :} la philosophie compte pour un coefficient significatif (généralement 2) dans les séries techniques : une bonne note y change une moyenne générale.
\item \textbf{Barème indicatif sur 20 :} généralités et présentation du texte (environ 4 points), réponses aux questions (environ 6 points), essai personnel (environ 8 points), qualité de la langue et de la présentation (environ 2 points).
\end{itemize}
\end{aretenir}

\begin{popfigure}
\begin{center}
\begin{tikzpicture}
% Manual pie chart: 20%, 30%, 40%, 10%
\colorlet{cGen}{popblue}
\colorlet{cQues}{poporange}
\colorlet{cEss}{popgreen}
\colorlet{cLang}{poppurple}
\fill[cGen] (0,0) -- (0:2.6) arc (0:72:2.6) -- cycle;
\fill[cQues] (0,0) -- (72:2.6) arc (72:180:2.6) -- cycle;
\fill[cEss] (0,0) -- (180:2.6) arc (180:324:2.6) -- cycle;
\fill[cLang] (0,0) -- (324:2.6) arc (324:360:2.6) -- cycle;
% White circle in center for readability (optional)
\fill[white] (0,0) circle (1.0);
% Labels inside slices (optional, but legend is in \legende)
\node[font=\small\bfseries, text=popdark] at (36:1.3) {20\%};
\node[font=\small\bfseries, text=popdark] at (126:1.3) {30\%};
\node[font=\small\bfseries, text=popdark] at (252:1.3) {40\%};
\node[font=\small\bfseries, text=popdark] at (342:0.8) {10\%};
\end{tikzpicture}
\end{center}
\legende{Répartition indicative des points sur 20 dans l'exercice sur texte au Probatoire et au Baccalauréat techniques. L'essai personnel pèse le plus lourd (40\%) : il faut lui réserver le plus de temps.}
\end{popfigure}

\begin{savaistu}[L'épreuve la plus choisie... et la moins réussie]
Au Baccalauréat technique camerounais, l'exercice sur texte est l'un des sujets les plus fréquemment choisis par les candidats, qui le croient « plus facile » que la dissertation parce que le texte est « sous les yeux ». Pourtant, c'est aussi l'un des moins bien réussis, faute de méthode : beaucoup de copies se contentent de paraphraser le texte ou ignorent purement et simplement les questions posées. La différence entre 6/20 et 14/20 tient presque toujours au respect de la démarche, non à la quantité de connaissances.
\end{savaistu}

\courssub{Les trois parties obligatoires de la copie}

Une copie d'exercice sur texte se construit comme un bâtiment à trois étages : on ne peut pas sauter l'un d'eux sans faire s'effondrer l'ensemble.

\begin{propriete}[Structure obligatoire]
Toute copie d'exercice sur texte comporte \textbf{trois parties obligatoires} :
\begin{enumerate}
\item Les \cle{généralités} : présentation de l'auteur, de l'œuvre, du thème et de la thèse du texte, ainsi que de sa structure (les grandes étapes du raisonnement).
\item Les \cle{réponses aux questions} : traitées \textbf{dans l'ordre}, chacune annoncée clairement, avec des réponses exactes, justifiées et appuyées sur le texte.
\item L'\cle{essai personnel} : une réflexion argumentée et structurée (introduction, développement, conclusion) sur le problème posé, où le candidat prend position.
\end{enumerate}
\end{propriete}

\begin{popfigure}
\begin{center}
\begin{tikzpicture}[node distance=0.9cm,
bloc/.style={rectangle, rounded corners, draw=popdark, thick, align=center, text width=4.2cm, minimum height=1.5cm, font=\small},
fleche/.style={-{Stealth[length=3mm]}, very thick, popdark}]
\node[bloc, fill=popblueL, align=center] (g){\textbf{1. Généralités}\\ Auteur, œuvre, thème, thèse, structure\\ \emph{(15--20 min)}};
\node[bloc, fill=poporangeL, below=of g, align=center] (q){\textbf{2. Réponses aux questions}\\ Dans l'ordre, exactes, justifiées\\ \emph{(45 min)}};
\node[bloc, fill=popgreenL, below=of q, align=center] (e){\textbf{3. Essai personnel}\\ Introduction, développement, conclusion\\ \emph{(1 h 30)}};
\draw[fleche] (g) -- (q);
\draw[fleche] (q) -- (e);
\end{tikzpicture}
\end{center}
\legende{Les trois blocs de la copie et leur répartition horaire conseillée sur 3 heures (il reste environ 15 minutes pour la relecture).}
\end{popfigure}

\begin{attention}[Le piège mortel : ignorer les questions]
L'erreur la plus fréquente est de confondre l'exercice sur texte avec une \textbf{dissertation libre} : le candidat rédige un long développement général sur « la liberté » ou « la vérité » sans jamais répondre aux questions posées. Or chaque question non traitée fait perdre des points précis au barème. Règle d'or : \textbf{les questions sont le contrat} ; on y répond toutes, dans l'ordre, et on les annonce (« Question 1 : ... »).
\end{attention}

\courssub{Les critères d'évaluation officiels}

Pour réussir, il faut savoir ce que le correcteur cherche. Quatre critères structurent la notation.

\begin{center}
\begin{tabularx}{\linewidth}{|l|X|}\hline
\rowcolor{popblueL} \textbf{Critère} & \textbf{Ce que le correcteur attend} \\ \hline
\cle{Fidélité au texte} & Les réponses s'appuient sur le texte (citations courtes, références précises), sans paraphrase ni hors-sujet. \\ \hline
\cle{Exactitude des réponses} & Les termes sont bien définis, la thèse est correctement identifiée, les arguments sont bien dégagés. \\ \hline
\cle{Cohérence de l'essai} & L'essai personnel est structuré, argumenté, prend position et reste en lien avec le problème du texte. \\ \hline
\cle{Qualité de la langue} & Orthographe, syntaxe, vocabulaire philosophique, présentation soignée de la copie. \\ \hline
\end{tabularx}
\end{center}

Autrement dit, le correcteur ne demande pas au candidat d'être « d'accord » avec l'auteur, mais d'être \textbf{fidèle} au texte dans l'analyse et \textbf{personnel} dans l'essai. On peut critiquer Kant à condition de l'avoir d'abord correctement compris.

\courssub{Gestion du temps : la répartition conseillée}

Trois heures, c'est long pour qui n'a pas de plan, court pour qui s'organise mal. La répartition suivante, éprouvée, doit devenir un automatisme.

\begin{methode}[Gérer ses 3 heures]
\begin{enumerate}
\item \textbf{Lecture (15 min)} : lire le texte \textbf{deux fois} avant d'écrire quoi que ce soit — une première fois pour le sens global, une seconde fois au crayon pour souligner la thèse, les mots-clés et les étapes du raisonnement. Lire aussi toutes les questions.
\item \textbf{Généralités (15--20 min)} : rédiger la présentation du texte au propre.
\item \textbf{Réponses aux questions (45 min)} : traiter les questions \textbf{dans l'ordre}, en justifiant chaque réponse par le texte.
\item \textbf{Essai personnel (1 h 30)} : réserver \textbf{au moins la moitié du temps total} à l'essai, qui pèse le plus lourd au barème : 10 minutes de plan au brouillon, puis rédaction.
\item \textbf{Relecture (10--15 min)} : corriger l'orthographe, vérifier qu'aucune question n'a été oubliée.
\end{enumerate}
\end{methode}

\courssub{La lecture analytique du texte}

La lecture analytique est l'étape fondatrice : c'est elle qui conditionne la qualité des généralités, des réponses aux questions et de l'essai. Elle ne se fait pas « en diagonale », mais au crayon, sur la copie du sujet (brouillon autorisé).

\begin{methode}[Les trois temps de la lecture analytique]
\begin{enumerate}
\item \textbf{Première lecture (globale) :} Lire sans s'arrêter pour saisir le \cle{sujet} (de quoi parle le texte ?), le \cle{thème} (quelle notion philosophique ?), la \cle{thèse} (que défend l'auteur ?) et le \cle{ton} (affirmatif, interrogatif, polémique, ironique ?).
\item \textbf{Deuxième lecture (détaillée, au crayon) :}
      \begin{itemize}
      \item Souligner les \cle{mots-clés} (notions techniques : liberté, conscience, État, etc.) et les \cle{connecteurs logiques} (donc, mais, cependant, car, si... alors...).
      \item Numéroter les \cle{étapes du raisonnement} (paragraphes ou groupes de phrases formant une unité logique).
      \item Repérer les \cle{citations} potentielles (expressions exactes à reprendre entre guillemets dans les réponses).
      \item Identifier les \cle{arguments} de l'auteur (exemples, analogies, réfutations d'objections, déductions).
      \end{itemize}
\item \textbf{Troisième lecture (croisée avec les questions) :} Relire le texte en ayant les questions devant soi. Pour chaque question, localiser précisément la zone du texte qui y répond et noter le numéro de ligne ou le paragraphe concerné.
\end{enumerate}
\end{methode}

\begin{aretenir}[Ce qu'il faut extraire de la lecture analytique]
\begin{itemize}
\item \textbf{Le thème} : la notion philosophique centrale (ex. : la liberté, le devoir, la technique).
\item \textbf{La thèse} : la réponse de l'auteur au problème (ex. : « La liberté est autonomie »).
\item \textbf{Le problème} : la tension ou la question implicite à laquelle le texte répond (ex. : « L'obéissance à la loi est-elle compatible avec la liberté ? »).
\item \textbf{La structure} : le découpage en 2, 3 ou 4 moments logiques (ex. : 1. Définition de la minorité, 2. Causes de la minorité, 3. Sortie de la minorité).
\item \textbf{L'intérêt philosophique} : pourquoi ce texte est-il un classique ? (nouveauté de la thèse, rigueur de l'argumentation, portée historique).
\end{itemize}
\end{aretenir}

\courssub{Les réponses aux questions : méthode et rédaction}

Les questions portent généralement sur : la définition d'un terme, l'identification de la thèse, l'explication d'un argument, la structure du texte, l'intérêt philosophique. Chaque réponse doit être une \textbf{mini-dissertation} : thèse (réponse directe), argument (justification par le texte), illustration (citation courte).

\begin{methode}[Répondre à une question type]
\begin{enumerate}
\item \textbf{Annoncer la question} : écrire « Question 1 : » ou « 1) » en marge ou en début de paragraphe.
\item \textbf{Répondre directement} (une phrase) : donner la définition, la thèse, l'argument principal.
\item \textbf{Justifier par le texte} : expliquer \emph{comment} le texte permet cette réponse, en citant \textbf{des mots ou groupes de mots exacts} (entre guillemets, intégrés à la phrase).
\item \textbf{Préciser la portée} (si nécessaire) : expliquer la signification philosophique du terme ou de l'argument dans la pensée de l'auteur.
\end{enumerate}
\end{methode}

\begin{exemplebox}[Modèle de réponse : « Que signifie 'minorité' chez Kant ? »]
\textbf{Mauvaise réponse :} « La minorité, c'est quand on est petit. » (Hors sujet, sens courant).
\textbf{Réponse insuffisante :} « C'est l'incapacité de se servir de son entendement. » (Paraphrase sans citation).
\textbf{Bonne réponse :} « La "minorité" désigne, pour Kant, l'\textbf{incapacité de se servir de son entendement sans la direction d'autrui} (citation). Ce n'est pas un état naturel (âge, QI), mais un état \textbf{volontaire} par "paresse" et "lâcheté", car l'individu refuse de penser par lui-même. »
\end{exemplebox}

\begin{attention}[Pièges à éviter dans les réponses]
\begin{itemize}
\item \textbf{La paraphrase :} recopier la phrase du texte en changeant deux mots ne prouve pas la compréhension.
\item \textbf{L'oubli de citation :} une réponse sans référence textuelle explicite (guillemets) perd des points de \cle{fidélité}.
\item \textbf{L'apport de connaissances extérieures :} les questions portent \textbf{sur le texte}. On n'y insère pas son cours sur Kant, sauf si la question demande explicitement « En quoi cette thèse est-elle caractéristique des Lumières ? ».
\item \textbf{Le désordre :} répondre Q3 avant Q1 trouble le correcteur et fait risquer l'oubli d'une question.
\end{itemize}
\end{attention}

\courssub{L'essai personnel : structure et argumentation}

L'essai personnel n'est pas une « dissertation bis » sur le thème général. C'est une \textbf{prise de position argumentée sur le problème précis soulevé par le texte}. Il doit dialoguer avec l'auteur : prolonger, nuancer, critiquer ou opposer une autre thèse.

\begin{propriete}[Structure type de l'essai personnel]
\begin{enumerate}
\item \textbf{Introduction} (10--15 lignes) :
      \begin{itemize}
      \item \cle{Accroche} : lien avec le texte (reprise de la thèse ou du problème).
      \item \cle{Problématisation} : formuler la question philosophique à laquelle l'essai va répondre (souvent donnée par le sujet, sinon la dégager).
      \item \cle{Annonce du plan} : 2 ou 3 grandes parties logiques (Thèse / Antithèse / Synthèse \emph{ou} Analyse / Critique / Dépassement).
      \end{itemize}
\item \textbf{Développement} (2 ou 3 parties équilibrées, chacune en 1--2 paragraphes) :
      \begin{itemize}
      \item Chaque partie = une \cle{idée directrice} annoncée en phrase de tête.
      \item Argumentation : \textbf{concept} (notion du cours) + \textbf{raisonnement} (pourquoi ?) + \textbf{exemple concret} (fait social, historique, littéraire, expérience vécue au Cameroun).
      \item \cle{Transition} explicite entre les parties.
      \end{itemize}
\item \textbf{Conclusion} (5--10 lignes) :
      \begin{itemize}
      \item \cle{Bilan} : réponse synthétique à la problématique.
      \item \cle{Ouverture} : prolongement (autre auteur, autre domaine, actualité) sans ouvrir un nouveau débat.
      \end{itemize}
\end{enumerate}
\end{propriete}

\begin{methode}[Construire son plan au brouillon (10 min)]
\begin{enumerate}
\item Formuler la \textbf{problématique} en une question claire.
\item Lister \textbf{tous les arguments} qui viennent (pour, contre, nuances) en s'appuyant sur le cours et des exemples camerounais.
\item Les regrouper en \textbf{2 ou 3 axes logiques} (pas en « arguments pour / arguments contre » juxtaposés, mais en étapes de pensée : ex. 1. La thèse du texte est valide car... 2. Elle rencontre une limite... 3. On peut la dépasser en...).
\item Rédiger les \textbf{phrases de tête} de chaque partie et les \textbf{transitions}.
\end{enumerate}
\end{methode}

\begin{savaistu}[Exemples camerounais pour illustrer]
\begin{itemize}
\item \textbf{Liberté / Autorité} : le port du casque à moto (contrainte légale vs choix individuel), le code de la route, l'autorité parentale vs autonomie du jeune.
\item \textbf{Technique / Nature} : l'agriculture intensive (engrais, OGM) vs agroécologie dans le Nord-Ouest ; l'exploitation forestière vs préservation de la biodiversité au Sud.
\item \textbf{Travail / Technique} : l'informel (call-box, moto-taxi) face à l'automatisation ; la formation technique comme émancipation.
\item \textbf{État / Société} : la décentralisation, le bilinguisme, la justice traditionnelle vs justice moderne.
\end{itemize}
Ces exemples ancrent la philosophie dans le réel et valent mieux que des références abstraites.
\end{savaistu}

\courssub{Exercice sur texte, commentaire, explication : ne pas confondre}

Au Bac, le vocabulaire des épreuves peut prêter à confusion. Distinguons clairement trois exercices voisins.

\begin{center}
\begin{tabularx}{\linewidth}{|l|X|X|}\hline
\rowcolor{popblueL} \textbf{Exercice} & \textbf{Principe} & \textbf{Spécificité} \\ \hline
\cle{Explication de texte} & Suivre le texte ligne par ligne pour en expliquer le détail du raisonnement. & Très proche du texte, presque sans recul personnel. \\ \hline
\cle{Commentaire de texte} & Dégager le sens, la thèse et l'enjeu du texte, puis le discuter. & Plus de recul critique, mais le plan reste libre. \\ \hline
\cle{Exercice sur texte} & Commentaire \textbf{guidé} : questions imposées + essai personnel. & Le sujet fixe le parcours ; l'essai final est obligatoire. \\ \hline
\end{tabularx}
\end{center}

\courssub{Application guidée : un sujet complet corrigé}

Voici un sujet type « Probatoire » traité intégralement pour montrer l'enchaînement des trois parties.

\begin{exoresolu}[Sujet complet : Kant, \emph{Qu'est-ce que les Lumières ?} (extrait)]
\textbf{Énoncé du sujet :}
\textit{« Les Lumières, c'est la sortie de l'homme hors de sa minorité dont il est lui-même responsable. La minorité, c'est l'incapacité de se servir de son entendement sans la direction d'autrui. [...] "Ayez le courage de vous servir de votre propre entendement !" Telle est la devise des Lumières.} \\
\textbf{Questions :}
1) Expliquez le sens du terme "minorité" dans ce texte.
2) Quelle est la thèse de Kant dans cet extrait ?
3) Dégagez les arguments par lesquels l'auteur soutient sa thèse.
\textbf{Essai personnel :} \textit{Obéir, est-ce renoncer à sa liberté ?}
\tcblower
\textbf{Solution détaillée}

\textbf{1. Généralités}
\textbf{Auteur :} Emmanuel Kant (1724--1804), philosophe allemand des Lumières, figure majeure du criticisme.
\textbf{Œuvre :} \emph{Qu'est-ce que les Lumières ?} (1784), essai court publié dans le \emph{Berlinische Monatsschrift}, texte manifeste des Lumières.
\textbf{Thème :} L'émancipation intellectuelle et morale de l'individu (l'autonomie).
\textbf{Thèse :} Les Lumières sont le processus par lequel l'homme sort de sa « minorité » (dépendance intellectuelle volontaire) par l'exercice courageux de sa propre raison, accédant ainsi à l'autonomie.
\textbf{Structure du texte :} Trois moments. \textbf{Moment 1 (Définition) :} Les Lumières = sortie de la minorité ; définition de la minorité (incapacité d'utiliser son entendement seul) et de sa cause (paresse, lâcheté). \textbf{Moment 2 (Conditions) :} La sortie exige « liberté » (usage public de la raison) et « courage » (devise : \emph{Sapere aude}). \textbf{Moment 3 (Portée) :} Distinction usage public / usage privé de la raison ; l'obéissance civile n'empêche pas la liberté de philosopher.

\textbf{2. Réponses aux questions}
\textbf{Question 1 :} La « minorité » ne désigne pas un âge (minorité civile) ni un déficit intellectuel inné. C'est un \textbf{état moral et intellectuel} : l'\textbf{incapacité de se servir de son entendement sans la direction d'autrui} (l. 2--3). Kant précise qu'on en est « soi-même responsable » (l. 1) car sa cause n'est pas le manque d'esprit, mais la « paresse et la lâcheté » (l. 4) : on préfère payer un médecin, un pasteur, un livre de morale pour ne pas avoir à penser ni à décider.
\textbf{Question 2 :} La thèse est que \textbf{les Lumières sont l'accession à l'autonomie par l'usage courageux de sa propre raison}. L'homme doit oser penser par lui-même (\emph{Sapere aude}) pour sortir de la tutelle volontaire. Cette sortie est un devoir et un droit.
\textbf{Question 3 :} Trois arguments principaux. \textbf{Argument de responsabilité :} la minorité est « dont il est lui-même responsable » (l. 1) : elle est choisie, non subie. \textbf{Argument psychologique :} les tuteurs (gardiens) entretiennent la minorité en montrant le danger de l'autonomie (« le pas le plus sûr est encore de rester mineur », l. 5--6) $\rightarrow$ peur entretenue. \textbf{Argument politique/juridique :} la sortie ne demande qu'une chose : « la liberté » (l. 7), entendue comme « la liberté de faire usage public de sa raison en tous domaines » (l. 8). Distinction cruciale : usage public (citoyen savant qui s'adresse au public) vs usage privé (fonctionnaire qui obéit). On peut obéir en tant que fonctionnaire (privé) et critiquer en tant que savant (public).

\textbf{3. Essai personnel : Obéir, est-ce renoncer à sa liberté ?}

\textbf{Introduction}
[Accroche] Kant vient de distinguer l'usage public et l'usage privé de la raison : le fonctionnaire doit obéir (usage privé), mais le citoyen doit pouvoir critiquer (usage public). [Problématique] Cette distinction suffit-elle à dire qu'obéir n'est pas renoncer à sa liberté ? N'y a-t-il pas des obéissances qui aliènent, et des désobéissances qui libèrent ? La question est de savoir si l'obéissance est structurellement contraire à la liberté ou si elle peut en être la condition. [Plan] Nous montrerons d'abord que l'obéissance à une loi rationnelle et consentie est l'expression même de la liberté (I), avant d'analyser les formes d'obéissance qui constituent une renonciation à la liberté : l'hétéronomie et la servitude volontaire (II). Nous verrons enfin que la vraie liberté réside dans l'autonomie critique, capable de désobéir aux lois injustes (III).

\textbf{Développement}
\textbf{I. Obéir à la loi qu'on s'est prescrite : la liberté comme autonomie (Kant, Rousseau).}
[Idée directrice] Pour Kant, dans \emph{Fondements de la métaphysique des mœurs}, la volonté est libre quand elle suit sa propre loi (autonomie), non un désir extérieur (hétéronomie). Obéir à la loi morale (l'impératif catégorique), c'est obéir à sa propre raison. [Exemple] Au Cameroun, un citoyen qui paie ses impôts non par peur de la prison, mais parce qu'il reconnaît la nécessité de financer les hôpitaux et les écoles, obéit librement : il est l'auteur de la loi via ses représentants.
[Transition] Mais toute obéissance n'est pas de cet ordre. Beaucoup relèvent de la soumission à une volonté étrangère.

\textbf{II. L'obéissance hétéronome : renoncement à sa liberté (La Boétie, Marx, Freud).}
[Idée directrice] La Boétie (\emph{Discours de la servitude volontaire}) montre que les tyrans n'ont de pouvoir que celui qu'on leur donne. Obéir par peur, par habitude, par intérêt matériel ou par aliénation idéologique, c'est renoncer à sa liberté. [Exemple] Le conducteur de moto-taxi à Douala qui porte le casque uniquement au barrage policier, et le retire ensuite, obéit par contrainte (hétéronomie), non par conviction. Il n'est pas libre, il subit. [Argument] Marx ajoute que l'obéissance au capital (salariat) peut être une aliénation structurelle : le travailleur « vend » sa liberté pour survivre.

\textbf{III. La liberté critique : savoir désobéir pour rester libre (Thoreau, Arendt, Droit positif).}
[Idée directrice] La liberté n'est pas l'indépendance absolue (impossible en société), mais la capacité d'évaluer les ordres. Hannah Arendt (\emph{Eichmann à Jérusalem}) montre que le mal radical vient de l'incapacité de penser (« banalité du mal ») : Eichmann obéissait aux lois de son État. [Exemple] Au Cameroun, la résistance à des lois coloniales injustes (ex. code de l'indigénat) ou la lutte pour le multipartisme dans les années 90 ont été des désobéissances civiles fondatrices de libertés nouvelles. [Synthèse] Le citoyen libre, c'est celui qui obéit aux lois justes (autonomie) et désobéit aux lois injustes (devoir de conscience), comme le prévoit d'ailleurs l'article 65 de la Constitution camerounaise (droit de résistance à l'oppression).

\textbf{Conclusion}
[Bilan] Obéir n'est pas \emph{nécessairement} renoncer à sa liberté : c'en est même l'acte fondateur quand la loi émane de la raison commune (autonomie). Mais l'obéissance aveugle, par peur ou paresse (la « minorité » kantienne), est bien une renonciation. [Ouverture] La liberté se conquiert donc dans l'exercice permanent du jugement critique : « Ayez le courage de vous servir de votre propre entendement », y compris face à la loi. C'est le prix de la majorité civique.

\end{exoresolu}

\begin{aretenir}[L'essentiel]
L'exercice sur texte est un \textbf{commentaire guidé} en trois parties obligatoires : généralités, réponses aux questions (dans l'ordre), essai personnel. Le correcteur évalue la fidélité au texte, l'exactitude des réponses, la cohérence de l'essai et la qualité de la langue. La réussite tient à la méthode — double lecture analytique, temps réservé à l'essai, aucune question ignorée — bien plus qu'à l'étendue des connaissances.
\end{aretenir}

\coursec{La lecture analytique du texte}

Dans la section précédente, nous avons présenté les généralités sur l'exercice sur texte : sa place à l'épreuve de philosophie du Probatoire et du Baccalauréat technique, sa structure (questions puis essai personnel) et ses exigences générales. Nous savons donc \emph{ce qu'est} cet exercice. Reste à apprendre \emph{comment} le réussir. Or tout commence par une étape décisive, trop souvent négligée par les candidats pressés : la \cle{lecture analytique} du texte. C'est elle qui fournit la matière de toutes les réponses ultérieures. Un candidat qui lit mal le texte répondra à côté, même avec beaucoup de connaissances ; un candidat qui lit bien le texte dispose déjà de l'essentiel de sa copie.

\courssub{Première lecture : le survol}

\begin{definition}[Lecture analytique]
La \cle{lecture analytique} est une lecture méthodique et active d'un texte philosophique, qui vise à en dégager le \cle{thème}, la \cle{problématique}, la \cle{thèse} et les \cle{arguments}, afin de comprendre exactement ce que l'auteur dit et comment il le démontre.
\end{definition}

La première lecture est un \cle{survol} : on lit le texte une première fois, sans stylo à la main, pour se faire une impression générale. Cette lecture rapide a trois objectifs précis :

\begin{enumerate}
\item \textbf{Identifier l'auteur.} Le nom de l'auteur figure toujours en dessous ou au-dessus du texte, souvent avec le titre de l'œuvre. Ce renseignement est précieux : savoir qu'un texte est de Kant, de Descartes ou de Nkrumah donne déjà des indications sur l'époque, le contexte et les grandes thèses de l'auteur.
\item \textbf{Identifier l'œuvre.} Le titre de l'ouvrage (par exemple \emph{Critique de la raison pratique}) éclaire le sujet traité.
\item \textbf{Saisir le sujet général.} Après cette première lecture, on doit pouvoir répondre en une phrase à la question : \og de quoi parle ce texte ? \fg{} (la liberté, le devoir, le bonheur, l'État, la technique...).
\end{enumerate}

\begin{attention}[Piège fréquent]
Ne commencez \textbf{jamais} à rédiger après une seule lecture. Le survol ne donne qu'une impression globale ; il ne suffit ni pour trouver la thèse exacte, ni pour repérer les arguments. Rédiger trop tôt, c'est risquer le hors-sujet, faute grave à l'examen.
\end{attention}

\courssub{Deuxième lecture : le soulignage actif}

La deuxième lecture se fait \emph{stylo en main}. Il s'agit d'une lecture lente, ligne par ligne, au cours de laquelle on souligne ou on annote trois catégories d'éléments :

\begin{itemize}
\item \textbf{Les mots-clés et les concepts} : les noms importants (liberté, loi, devoir, conscience, État...) et les expressions techniques de la philosophie. Ce sont eux qui indiquent le thème.
\item \textbf{Les articulations logiques} (ou connecteurs) : ce sont les mots qui relient les idées entre elles et qui révèlent la structure du raisonnement.
\item \textbf{Les verbes d'affirmation} : \og il faut \fg{}, \og on doit \fg{}, \og il est clair que \fg{}, \og nous soutenons que \fg{}... Ils signalent que l'auteur prend position.
\end{itemize}

\begin{center}
\begin{tabularx}{\linewidth}{|l|X|X|}
\hline
\rowcolor{popblueL} \textbf{Type de connecteur} & \textbf{Exemples} & \textbf{Fonction dans le texte} \\
\hline
Cause, explication & car, en effet, parce que & Introduit une justification \\
\hline
Conséquence, conclusion & donc, ainsi, par conséquent, c'est pourquoi & Introduit la thèse ou une conclusion partielle \\
\hline
Opposition, objection & mais, cependant, pourtant, néanmoins & Introduit une objection ou une restriction \\
\hline
Ajout & de plus, en outre, également & Introduit un argument supplémentaire \\
\hline
Illustration & par exemple, ainsi qu'on le voit & Introduit un exemple \\
\hline
\end{tabularx}
\end{center}

\begin{exemplebox}[Lire les connecteurs]
Dans la phrase : \og La liberté semble être le pouvoir de faire ce que l'on veut ; \emph{cependant}, obéir à tous ses désirs, c'est être esclave de ses passions ; \emph{donc} la vraie liberté consiste à obéir à la loi qu'on s'est prescrite \fg{}, le mot \emph{cependant} signale une objection à une première idée, et le mot \emph{donc} signale la thèse de l'auteur.
\end{exemplebox}

\courssub{Dégager le thème : de quoi parle le texte ?}

Le \cle{thème} est le sujet général du texte, c'est-à-dire le domaine ou la notion dont il est question. Pour le dégager, on procède ainsi :

\begin{enumerate}
\item On relève les mots qui reviennent le plus souvent (les répétitions sont des indices).
\item On regroupe les mots proches en un champ lexical (par exemple : \emph{loi, devoir, obligation, interdiction} forment le champ lexical de la morale).
\item On formule le thème en un ou deux mots, sous forme d'un nom ou d'une courte expression : \og la liberté \fg{}, \og le rapport entre la liberté et la loi \fg{}.
\end{enumerate}

\begin{attention}[Le thème n'est pas la thèse]
Erreur très fréquente à l'examen : confondre le \cle{thème} (le sujet général, ce dont on parle) avec la \cle{thèse} (la position précise défendue par l'auteur sur ce sujet). Dire \og le texte parle de la liberté \fg{}, c'est donner le thème ; dire \og l'auteur soutient que la liberté consiste à obéir à la loi morale \fg{}, c'est donner la thèse. Le thème se formule en quelques mots ; la thèse se formule en une phrase complète avec un verbe.
\end{attention}

\courssub{Dégager la thèse : que soutient l'auteur ?}

\begin{definition}[La thèse]
La \cle{thèse} est la position défendue par l'auteur, l'idée principale qu'il cherche à démontrer. Elle doit pouvoir être exprimée en \textbf{une seule phrase}, claire et complète, de la forme : \og Pour l'auteur, ... \fg{}
\end{definition}

\begin{methode}[Trouver la thèse en quatre étapes]
\begin{enumerate}
\item Repérer les \textbf{verbes d'affirmation} par lesquels l'auteur prend position : \og il faut \fg{}, \og on doit \fg{}, \og il est clair que \fg{}, \og nous affirmons \fg{}, \og il est évident que \fg{}.
\item Repérer les \textbf{conclusions} introduites par \og donc \fg{}, \og ainsi \fg{}, \og par conséquent \fg{}, \og c'est pourquoi \fg{} : la thèse se trouve souvent juste après ces mots.
\item Vérifier que la phrase retenue \textbf{résume l'ensemble du texte} : si une partie du texte contredit votre formulation, ce n'est pas la bonne thèse.
\item Reformuler la thèse \textbf{avec vos propres mots}, en une seule phrase, sans recopier mot à mot une longue phrase de l'auteur.
\end{enumerate}
\end{methode}

\courssub{Construire la problématique : quelle question le texte résout-il ?}

\begin{definition}[La problématique]
La \cle{problématique} est la question centrale à laquelle le texte apporte une réponse. Elle se formule sous forme interrogative et relie le thème à la thèse : la thèse est précisément la réponse de l'auteur à la problématique.
\end{definition}

Pour construire la problématique, on peut utiliser un raisonnement simple : \emph{si la thèse est la réponse, quelle était la question ?} Par exemple, si la thèse est \og la vraie liberté consiste à obéir à la loi qu'on s'est prescrite \fg{}, alors la problématique est : \og En quoi consiste la véritable liberté ? \fg{} ou encore \og La liberté est-elle l'absence de toute contrainte, ou suppose-t-elle l'obéissance à une loi ? \fg{}

\begin{propriete}[Le triangle thème -- problématique -- thèse]
Les trois éléments sont solidaires : le \textbf{thème} indique le sujet, la \textbf{problématique} pose la question sur ce sujet, la \textbf{thèse} répond à cette question. Une bonne analyse vérifie toujours cette cohérence : ma problématique porte-t-elle bien sur mon thème ? Ma thèse répond-elle bien à ma problématique ?
\end{propriete}

\begin{popfigure}
\begin{center}
\begin{tikzpicture}[
  node distance=0.55cm and 1.1cm,
  every node/.style={font=\small},
  box/.style={draw=popblue, thick, rounded corners=3pt, fill=popblueL, text width=3.1cm, align=center, minimum height=0.9cm},
  arg/.style={draw=poporange, thick, rounded corners=3pt, fill=poporangeL, text width=2.6cm, align=center, minimum height=0.8cm},
  ex/.style={draw=popgreen, thick, rounded corners=3pt, fill=popgreenL, text width=2.6cm, align=center, minimum height=0.7cm},
  fleche/.style={-{Stealth[length=2.5mm]}, thick, popdark}
]
\node[box, align=center] (theme){\textbf{Thème}\\ De quoi parle le texte ?};
\node[box, below=of theme, align=center] (pb){\textbf{Problématique}\\ Quelle question est posée ?};
\node[box, below=of pb, align=center] (these){\textbf{Thèse}\\ Que soutient l'auteur ?};
\node[arg, below left=0.7cm and -0.4cm of these] (a1) {\textbf{Argument 1}};
\node[arg, below=0.7cm of these] (a2) {\textbf{Argument 2}};
\node[arg, below right=0.7cm and -0.4cm of these] (a3) {\textbf{Argument 3}};
\node[ex, below=0.5cm of a1] (e1) {Exemple 1};
\node[ex, below=0.5cm of a2] (e2) {Exemple 2};
\node[ex, below=0.5cm of a3] (e3) {Exemple 3};
\draw[fleche] (theme) -- (pb);
\draw[fleche] (pb) -- (these);
\draw[fleche] (these) -- (a1);
\draw[fleche] (these) -- (a2);
\draw[fleche] (these) -- (a3);
\draw[fleche] (a1) -- (e1);
\draw[fleche] (a2) -- (e2);
\draw[fleche] (a3) -- (e3);
\end{tikzpicture}
\end{center}
\legende{Arbre d'analyse d'un texte philosophique : on descend du thème vers la problématique, puis vers la thèse, que soutiennent les arguments (A1, A2, A3), eux-mêmes illustrés par des exemples.}
\end{popfigure}

\courssub{Repérer les arguments et les exemples}

Une fois la thèse dégagée, il faut identifier \emph{comment} l'auteur la justifie. Un \cle{argument} est une raison avancée pour convaincre le lecteur que la thèse est vraie. Un \cle{exemple} est un cas particulier, concret, qui illustre l'argument.

\begin{aretenir}[Combien d'arguments dans un texte ?]
Un texte d'examen d'environ 15 lignes contient en moyenne \textbf{2 à 4 arguments principaux}. Les repérer \emph{tous} garantit l'exhaustivité de l'analyse : oublier un argument, c'est perdre des points sur les questions qui portent sur la structure du texte. Numérotez-les dans la marge (A1, A2, A3...) au fur et à mesure de votre lecture.
\end{aretenir}

Pour distinguer argument et exemple, on peut se poser la question : cet énoncé donne-t-il une \emph{raison générale} (argument) ou un \emph{cas particulier} (exemple) ? \og Celui qui obéit à tous ses désirs devient dépendant \fg{} est un argument ; \og ainsi, le fumeur qui ne peut s'empêcher de fumer n'est pas libre \fg{} est un exemple qui illustre cet argument.

\courssub{Repérer les objections et les réponses de l'auteur}

Les bons textes philosophiques ne se contentent pas d'affirmer : ils anticipent les \cle{objections}, c'est-à-dire les critiques qu'un adversaire pourrait adresser à la thèse, et ils y répondent. Repérer ce jeu thèse / objection / réponse est un signe de lecture fine, très valorisé par les correcteurs.

Les objections sont souvent introduites par des formules comme : \og on pourrait objecter que... \fg{}, \og certains diront que... \fg{}, \og mais alors... \fg{}, ou simplement par un \og cependant \fg{} suivi d'une idée contraire à la thèse. La réponse de l'auteur suit généralement, marquée par \og mais \fg{}, \og en réalité \fg{}, \og or \fg{}.

\begin{exemplebox}[Thèse, objection, réponse]
\og La liberté est le pouvoir absolu de l'individu [\emph{première idée, souvent celle du sens commun}]. On pourrait croire qu'elle exige le rejet de toute loi [\emph{objection implicite}]. Mais une liberté sans loi se détruit elle-même, car la liberté de chacun se heurte alors à celle des autres [\emph{réponse de l'auteur}]. \fg{}
\end{exemplebox}

\courssub{Application guidée : lecture analytique d'un texte de Kant sur la liberté}

Appliquons maintenant toute la méthode à un extrait adapté de Kant (\emph{Fondements de la métaphysique des mœurs}) :

\begin{exemplebox}[Texte à analyser]
\og La liberté ne consiste pas à faire tout ce que l'on veut. En effet, celui qui suit simplement ses désirs obéit à des penchants qu'il n'a pas choisis : il est l'esclave de ses inclinations. On pourrait objecter qu'obéir à une loi, c'est justement perdre sa liberté. Mais la loi morale n'est pas une loi étrangère : c'est la loi que notre propre raison se donne à elle-même. Ainsi, l'homme qui agit par devoir, et non par inclination, est seul véritablement libre. Par conséquent, la liberté est l'obéissance à la loi que l'on s'est prescrite. \fg{}
\end{exemplebox}

Voici comment le candidat annote ce texte dans ses marges :

\begin{popfigure}
\begin{center}
\begin{tikzpicture}
\node[draw=popline, thick, fill=popcream, text width=8.6cm, align=justify, font=\small, inner sep=8pt] (txt)
{\textcolor{popmuted}{\og} La liberté ne consiste pas à faire tout ce que l'on veut. \textbf{En effet}, celui qui suit simplement ses désirs obéit à des penchants qu'il n'a pas choisis : il est l'esclave de ses inclinations. On pourrait objecter qu'obéir à une loi, c'est justement perdre sa liberté. \textbf{Mais} la loi morale n'est pas une loi étrangère : c'est la loi que notre propre raison se donne à elle-même. Ainsi, l'homme qui agit par devoir, et non par inclination, est seul véritablement libre. \textbf{Par conséquent}, la liberté est l'obéissance à la loi que l'on s'est prescrite. \textcolor{popmuted}{\fg}};
\node[draw=popblue, fill=popblueL, rounded corners=2pt, text width=3.4cm, align=center, font=\footnotesize, right=0.5cm of txt.north east, anchor=north west] (n1) {\textbf{Thème}\\ la liberté};
\node[draw=poporange, fill=poporangeL, rounded corners=2pt, text width=3.4cm, align=center, font=\footnotesize, below=0.25cm of n1] (n2) {\textbf{Argument 1}\\ suivre ses désirs = être esclave};
\node[draw=poppurple, fill=poppurpleL, rounded corners=2pt, text width=3.4cm, align=center, font=\footnotesize, below=0.25cm of n2] (n3) {\textbf{Objection}\\ obéir = perdre sa liberté ?};
\node[draw=popgreen, fill=popgreenL, rounded corners=2pt, text width=3.4cm, align=center, font=\footnotesize, below=0.25cm of n3] (n4) {\textbf{Réponse}\\ la loi vient de ma raison};
\node[draw=popink, fill=poppinkL, rounded corners=2pt, text width=3.4cm, align=center, font=\footnotesize, below=0.25cm of n4] (n5) {\textbf{Thèse}\\ liberté = obéir à la loi qu'on s'est prescrite};
\draw[-{Stealth}, popblue] (n1.west) -- ++(-0.35,0);
\draw[-{Stealth}, poporange] (n2.west) -- ++(-0.35,0);
\draw[-{Stealth}, poppurple] (n3.west) -- ++(-0.35,0);
\draw[-{Stealth}, popgreen] (n4.west) -- ++(-0.35,0);
\draw[-{Stealth}, popink] (n5.west) -- ++(-0.35,0);
\end{tikzpicture}
\end{center}
\legende{Exemple de texte annoté : les connecteurs logiques sont mis en gras dans le texte ; les annotations en marge identifient le thème, l'argument, l'objection, la réponse et la thèse.}
\end{popfigure}

\begin{exoresolu}[Analyse complète de l'extrait de Kant]
Énoncé : après lecture analytique de l'extrait ci-dessus, dégagez le thème, la problématique, la thèse, les arguments, l'objection et sa réponse.
\tcblower
\textbf{Solution détaillée.}
\begin{itemize}
\item \textbf{Thème} : la liberté (et son rapport à la loi).
\item \textbf{Problématique} : la liberté consiste-t-elle à faire ce que l'on veut, ou suppose-t-elle l'obéissance à une loi ?
\item \textbf{Thèse} : pour Kant, la véritable liberté est l'obéissance à la loi morale que notre raison se prescrit à elle-même (signalée par \og par conséquent \fg{}).
\item \textbf{Argument 1} : celui qui suit ses désirs obéit à des penchants non choisis ; il est donc esclave de ses inclinations, non libre (introduit par \og en effet \fg{}).
\item \textbf{Argument 2} : la loi morale n'est pas une contrainte extérieure, car elle vient de notre propre raison ; y obéir, c'est donc s'obéir à soi-même.
\item \textbf{Objection} : obéir à une loi semble être une perte de liberté.
\item \textbf{Réponse de l'auteur} : cette objection tombe si la loi est celle que la raison se donne elle-même ; l'obéissance est alors autonomie, non soumission.
\end{itemize}
Vérification de cohérence : la thèse répond bien à la problématique, et les deux arguments soutiennent bien la thèse. L'analyse est complète : le texte compte ici 2 arguments principaux, ce qui correspond à la moyenne (2 à 4) pour un texte court.
\end{exoresolu}

\begin{savaistu}[Une voix africaine de la liberté : Kwame Nkrumah]
Le même exercice peut être mené sur des auteurs africains. Kwame Nkrumah (1909-1972), père de l'indépendance du Ghana, soutient dans ses écrits que la liberté politique d'un peuple exige son indépendance et l'unité de l'Afrique : \og Cherchez d'abord le royaume politique, et tout le reste vous sera donné par surcroît. \fg{} Sa thèse : sans liberté politique, aucun développement économique n'est possible. Cette idée a inspiré de nombreux penseurs africains, y compris au Cameroun, où la question de la liberté du citoyen reste au cœur des débats publics.
\end{savaistu}

\begin{exercice}[À vous de jouer]
Lisez deux fois le texte suivant (extrait adapté) : \og Beaucoup pensent que la liberté consiste à ne dépendre de rien ni de personne. Cependant, l'homme qui vit seul, sans règle ni loi, est en proie à toutes les craintes. En effet, sans loi commune, le plus fort écrase le plus faible. Donc la liberté véritable ne naît que dans une société réglée par des lois justes. \fg{}
Dégagez : 1) le thème ; 2) la problématique ; 3) la thèse ; 4) les arguments ; 5) l'objection éventuelle.
\end{exercice}

\begin{corrige}[Exercice : analyse du texte sur la liberté et la loi]
\begin{itemize}
\item \textbf{Thème} : la liberté et son rapport à la loi.
\item \textbf{Problématique} : la liberté consiste-t-elle dans l'indépendance absolue, ou dépend-elle des lois ?
\item \textbf{Thèse} : la liberté véritable ne naît que dans une société réglée par des lois justes (repérée grâce à \og donc \fg{}).
\item \textbf{Arguments} : A1 --- sans loi commune, le plus fort écrase le plus faible (introduit par \og en effet \fg{}) ; A2 --- l'homme sans règle vit dans la crainte permanente.
\item \textbf{Objection} : le sens commun pense que la liberté, c'est ne dépendre de rien (\og beaucoup pensent que... \fg{}, réfuté par \og cependant \fg{}).
\end{itemize}
\end{corrige}

\begin{aretenir}[L'essentiel de la lecture analytique]
\begin{enumerate}
\item Première lecture : survol pour identifier l'auteur, l'œuvre et le sujet général.
\item Deuxième lecture : souligner les mots-clés, les concepts et les connecteurs logiques (\emph{donc, mais, cependant, en effet}).
\item Dégager dans l'ordre : le \textbf{thème} (le sujet), la \textbf{problématique} (la question), la \textbf{thèse} (la réponse de l'auteur en une phrase).
\item Repérer tous les \textbf{arguments} (2 à 4 en moyenne), les \textbf{exemples}, les \textbf{objections} et les \textbf{réponses}.
\item Vérifier la cohérence : thème $\rightarrow$ problématique $\rightarrow$ thèse $\rightarrow$ arguments.
\end{enumerate}
Cette analyse achevée, vous êtes prêt pour l'étape suivante : la rédaction des réponses aux questions, que nous étudierons dans la prochaine section.
\end{aretenir}

\coursec{Les réponses aux questions}

Après avoir maîtrisé la lecture analytique qui permet de « faire parler le texte », nous abordons maintenant l'épreuve centrale : la rédaction des réponses aux questions. C'est ici que se joue une grande partie de la note au Probatoire. Contrairement à une idée reçue, il ne s'agit pas de « réciter le cours » ni de « dire ce qu'on pense », mais de \textbf{démontrer sa compréhension fine du texte} en suivant une méthode rigoureuse, codifiée par les correcteurs du MINESEC. La qualité d'une réponse ne se juge pas à sa longueur, mais à sa \textbf{pertinence}, sa \textbf{précision} et sa \textbf{justification textuelle}.

\courssub{Typologie des questions : identifier l'exigence pour y répondre}

Au Probatoire technique, les questions posées sur le texte ne sont pas homogènes. Elles obéissent à une progression logique : du fait (ce que dit le texte) vers le sens (comment il le dit) puis vers la portée (ce que cela vaut). Savoir classifier la question dès la lecture du sujet permet de choisir la bonne stratégie de réponse et de ne pas confondre \emph{explication} et \emph{appréciation personnelle}.

\begin{popfigure}
\centering
\begin{tikzpicture}[
    node distance=2mm and 4mm,
    header/.style={rectangle, draw=popblue, thick, fill=popblueL, text width=4.5cm, align=center, font=\bfseries\sffamily, minimum height=1.2cm},
    qtype/.style={rectangle, draw=popdark, fill=popcream, text width=4.5cm, align=center, font=\sffamily, minimum height=1.1cm},
    strat/.style={rectangle, draw=popgreen, fill=popgreenL, text width=6.5cm, align=left, font=\sffamily\small, minimum height=1.1cm},
    arrow/.style={-Stealth, popline, thick}
]
% En-têtes
\node[header] (h1) {Type de question};
\node[header, right=of h1] (h2) {Stratégie de réponse attendue};

% Ligne 1 : Compréhension
\node[qtype, below=of h1, align=center] (q1){Question de \textbf{compréhension}\\(Sens d'un passage, définition, thèse)};
\node[strat, right=of q1, align=center] (s1){\textbf{Reformuler} la réponse en ses mots.\\\textbf{Citer} le passage exact (guillemets).\\Ne pas interpréter, restituer.};

% Ligne 2 : Analyse
\node[qtype, below=of q1, align=center] (q2){Question d'\textbf{analyse}\\(Structure, arguments, procédés, concepts)};
\node[strat, right=of q2, align=center] (s2){\textbf{Identifier} le procédé ou l'enchaînement.\\\textbf{Expliquer} comment il fonctionne dans le texte.\\\textbf{Montrer} son rôle dans la démonstration.};

% Ligne 3 : Réflexion
\node[qtype, below=of q2, align=center] (q3){Question de \textbf{réflexion}\\(Portée, limites, actualité, critique)};
\node[strat, right=of q3, align=center] (s3){\textbf{Partir} du texte (thèse/argument).\\\textbf{Confrontation} : exemples précis, contre-exemples, auteurs connus.\\\textbf{Nuancer} : « Si... toutefois... en définitive ».};

% Flèches visuelles
\draw[arrow] (q1.east) -- (s1.west);
\draw[arrow] (q2.east) -- (s2.west);
\draw[arrow] (q3.east) -- (s3.west);
\end{tikzpicture}
\legende{Tableau de correspondance : Type de question $\rightarrow$ Stratégie de réponse. Ce tableau doit être mentalisé avant de rédiger.}
\end{popfigure}

\begin{definition}[Question de compréhension]
Elle vise à vérifier que le candidat a saisi le sens littéral d'un passage, la définition d'un concept clé utilisé par l'auteur, ou la thèse principale défendue. Elle appelle une réponse \textbf{fidèle au texte}, sans ajout extérieur.
\end{definition}

\begin{definition}[Question d'analyse]
Elle demande de mettre au jour la \textbf{structure argumentative} (enchaînement des idées, types de raisonnement : déduction, induction, analogie, \emph{reductio ad absurdum}) ou les \textbf{procédés rhétoriques/stylistiques} (ironie, comparaison, répétition, champ lexical) qui servent la thèse. Elle exige une réponse \textbf{technique} appuyée sur le texte.
\end{definition}

\begin{definition}[Question de réflexion (ou d'appréciation)]
Elle invite à discuter la portée de la thèse, ses limites, son actualité ou sa confrontation avec d'autres points de vue. C'est la seule où l'apport personnel (culture générale, exemples concrets camerounais, connaissances philosophiques) est \textbf{exigé}, mais il doit \textbf{toujours s'ancrer dans le texte} (point de départ ou point d'arrivée).
\end{definition}

\courssub{La règle d'or : répondre à la question posée, pas à côté}

La première cause d'échec à cet exercice est le \textbf{hors-sujet déguisé} : le candidat connaît le thème (ex. la liberté), il récite son cours sur la liberté, mais il ne répond pas à la question précise (ex. « Selon l'auteur, en quoi la contrainte extérieure est-elle compatible avec la liberté ? »).

\begin{methode}[Décryptage de l'intitulé de la question]
\begin{enumerate}
    \item \textbf{Surligner les mots-opérateurs} : « Selon l'auteur », « Dans le texte », « Expliquez », « Montrez », « Discutez », « Illustrer ».
    \item \textbf{Identifier la contrainte de source} :
    \begin{itemize}
        \item « Selon le texte / l'auteur / dans le passage \emph{l. 5 à 10} » $\rightarrow$ \textbf{Réponse intra-textuelle uniquement}. Aucun exemple personnel, aucune connaissance extérieure.
        \item « À partir du texte, discutez... » / « Qu'en pensez-vous ? » $\rightarrow$ \textbf{Réponse hybride} : point de départ textuel + apport personnel argumenté.
    \end{itemize}
    \item \textbf{Reformuler la question en une phrase affirmative} qui deviendra la première phrase de votre réponse (l'amorce).
\end{enumerate}
\end{methode}

\begin{exemplebox}[Exemple de reformulation]
\textbf{Sujet :} « Expliquez pourquoi l'auteur affirme que "la technique n'est pas un simple outil" (l. 12). »
\par\medskip
\textbf{Mauvaise amorce :} « La technique est très importante dans nos sociétés modernes... » (Hors-sujet, cours général).
\par\medskip
\textbf{Bonne amorce :} « L'auteur affirme que la technique n'est pas un simple outil \textbf{parce qu'elle structure notre rapport au monde et transforme l'homme lui-même}, comme le montre l'argumentation développée aux lignes 10-15. » (Réponse directe, annonce la justification textuelle).
\end{exemplebox}

\courssub{Rédaction des réponses : phrases complètes, citations courtes, justifications précises}

Une réponse n'est pas une liste de mots-clés. C'est un \textbf{paragraphe argumenté} rédigé en français correct, structuré par des connecteurs logiques (\emph{donc, car, en effet, par conséquent, toutefois}).

\begin{popfigure}
\centering
\begin{tikzpicture}[
    node distance=8mm and 4mm,
    box/.style={rectangle, draw=popblue, thick, fill=popblueL, rounded corners, text width=5cm, align=center, font=\sffamily\bfseries, minimum height=1.5cm},
    arrow/.style={-Stealth, popblue, thick},
    label/.style={font=\sffamily\small, text=popdark, align=center, midway, above, sloped}
]
\node[box, align=center] (r1){1. \textbf{REFORMULATION}\\\textit{« La question porte sur... »\\\textit{« L'auteur montre que... »}}};
\node[box, right=of r1, align=center] (r2){2. \textbf{RÉPONSE DIRECTE}\\\textit{« Il affirme que... »\\\textit{« La thèse est que... »}}};
\node[box, right=of r2, align=center] (r3){3. \textbf{JUSTIFICATION TEXTUELLE}\\\textit{« Comme l'indique la citation :\\\textit{« "... citation courte ... " (l. X) »}}};
\node[box, right=of r3, align=center] (r4){4. \textbf{EXPLICATION / ANALYSE}\\\textit{« Cela signifie que... »\\\textit{« Ce procédé renforce... »}}};

\draw[arrow] (r1) -- (r2);
\draw[arrow] (r2) -- (r3);
\draw[arrow] (r3) -- (r4);
\end{tikzpicture}
\legende{Schéma de la structure standard d'une bonne réponse aux questions (Compréhension / Analyse). Pour la Réflexion, l'étape 4 s'élargit à la discussion personnelle.}
\end{popfigure}

\begin{propriete}[Règle de la citation efficace]
Toute affirmation sur le texte doit être \textbf{ancrée} par une citation. Une citation efficace :
\begin{itemize}
    \item Est \textbf{courte} (quelques mots à une ligne maximum).
    \item Est encadrée de \textbf{guillemets français} (\og ... \fg{}).
    \item Est suivie de la \textbf{référence de ligne} : (l. 5) ou (l. 12-14).
    \item Est \textbf{intégrée à la phrase} (pas de citation « parachutée » isolée entre deux paragraphes).
\end{itemize}
\emph{Exemple :} L'auteur définit la technique comme un « mode de révélation » (l. 8), ce qui implique qu'elle n'est pas neutre.
\end{propriete}

\begin{attention}[Piège sanctionné : La paraphrase ou le « copier-coller »]
Recopier de longs passages du texte (3 lignes ou plus) \textbf{ne constitue pas une réponse}. C'est de la \textbf{paraphrase}. Le correcteur y voit une absence de compréhension (l'élève « noie le poisson ») et sanctionne sévèrement (souvent 0 ou 1 sur 4-5 points).
\par\medskip
\textbf{Règle :} On cite pour \textbf{prouver}, on explique pour \textbf{montrer qu'on a compris}. La citation est la \emph{preuve}, l'explication est la \emph{valeur ajoutée}.
\end{attention}

\begin{exemplebox}[Citation efficace vs Paraphrase]
\textbf{Texte (extrait) :} « L'homme moderne a délégué à la machine le soin de sa mémoire, de son calcul, et bientôt de son jugement. Il devient ainsi le serviteur de son propre outil. »
\par\medskip
\textbf{Question :} « Quel risque l'auteur identifie-t-il dans la délégation aux machines ? »
\par\medskip
\textbf{\textcolor{red}{Mauvaise réponse (Paraphrase) :}} L'auteur dit que « L'homme moderne a délégué à la machine le soin de sa mémoire, de son calcul, et bientôt de son jugement. Il devient ainsi le serviteur de son propre outil. » C'est très grave.
\par\medskip
\textbf{\textcolor{popgreen}{Bonne réponse :}} L'auteur identifie le risque d'une \textbf{inversion des rôles} : l'homme devient le « serviteur de son propre outil » (l. 3) en lui déléguant ses facultés cognitives essentielles (« mémoire », « calcul », « jugement », l. 1-2).
\end{exemplebox}

\begin{exemplebox}[Exemple chiffré : La longueur de la citation]
Sur une copie type (lignes de 10-12 mots), une citation de \textbf{15 à 20 mots} (environ 1 ligne et demie) est le maximum toléré. Au-delà, le correcteur considère que le candidat « remplit sa copie » au lieu d'analyser.
\par\medskip
\textbf{Astuce :} Utilisez les \textbf{points de suspension} entre crochets \og [...] \fg{} pour ne garder que les mots-clés d'une longue phrase.
\par\medskip
\textit{Exemple :} « L'homme [...] devient le serviteur de son propre outil » (l. 1-3).
\end{exemplebox}

\courssub{Hiérarchisation et gestion du temps : l'ordre stratégique}

Les questions ne sont pas numérotées par hasard. Elles suivent l'ordre du texte et la progression de difficulté. La gestion du temps (environ 1h30 pour l'ensemble de l'exercice sur texte au Probatoire) impose une discipline.

\begin{methode}[Ordre de traitement recommandé]
\begin{enumerate}
    \item \textbf{Questions factuelles (Compréhension / Analyse) :} Traitez-les \textbf{en premier}. Elles sont « faciles » (la réponse est \emph{dans} le texte), rapportent des points sûrs, et vous obligent à relire le texte en profondeur, ce qui prépare la réflexion.
    \item \textbf{Questions de réflexion (Discussion / Portée) :} Traitez-les \textbf{en dernier}. Elles demandent du recul, de la rédaction soignée, des exemples personnels. Ne les sacrifiez pas par manque de temps : elles valent souvent le plus de points (6 à 8 points sur 20).
    \item \textbf{Relisez} vos réponses aux questions factuelles avant d'attaquer la réflexion pour vous assurer qu'elles sont cohérentes avec votre discussion finale.
\end{enumerate}
\end{methode}

\begin{propriete}[Barème implicite : La justification vaut la réponse]
Dans la grande majorité des grilles de correction du MINESEC pour le Probatoire technique :
\begin{itemize}
    \item Une réponse \textbf{juste mais non justifiée} (pas de citation, pas de renvoi précis) = \textbf{50\% des points} (ex: 1/2, 2/4).
    \item Une réponse \textbf{juste et justifiée} (citation courte + n° de ligne + explication) = \textbf{100\% des points}.
    \item Une réponse \textbf{fausse mais justifiée par une citation mal comprise} = 0 point (la justification ne valide pas l'erreur).
\end{itemize}
\emph{Moralité :} Mieux vaut une réponse courte, juste, avec une citation précise, qu'une longue dissertation vague sans référence textuelle.
\end{propriete}

\courssub{Entraînement guidé : Sujet type Probatoire Technique (Session 2023 - Adapté)}

Appliquons la méthode sur un sujet complet simplifié. Le texte est un extrait fictif mais représentatif d'un auteur type programme (ex. Hans Jonas, Günther Anders, ou un texte sur l'éthique de la technologie).

\begin{exoresolu}[Sujet d'entraînement : L'homme face à la technique]
\textbf{Texte (Extrait) :}
\og Nous avons cru que la technique était un moyen neutre, un instrument docile à notre volonté. Mais l'histoire récente nous enseigne que les moyens finissent par déterminer les fins. En automatisant la production, nous n'avons pas seulement libéré l'homme de la peine ; nous avons modifié la nature du travail, réduit l'autonomie du geste, et soumis le rythme humain à la cadence de la machine. L'efficacité devient la valeur suprême, reléguant au second plan la justice, la beauté ou le sens. Ainsi, ce que nous pensions maîtriser nous maîtrise : le système technique s'auto-entretient, indifférent aux destins humains. \fg{}

\textbf{Questions :}
\begin{enumerate}
    \item \textbf{(Compréhension - 3 pts)} Selon l'auteur, quelle illusion l'homme a-t-il eue sur la technique ?
    \item \textbf{(Analyse - 4 pts)} Montrez, en vous appuyant sur le texte, comment l'auteur structure son argumentation pour démontrer que « les moyens déterminent les fins ».
    \item \textbf{(Réflexion - 8 pts)} « La technique nous maîtrise plus que nous ne la maîtrisons. » Discutez cette thèse en vous appuyant sur le texte et sur des exemples concrets.
\end{enumerate}

\tcblower

\textbf{Solution détaillée et commentée}

\par\medskip\noindent\textbf{Question 1 (Compréhension - 3 pts)}
\par\medskip
\emph{Analyse de la consigne :} « Selon l'auteur » $\rightarrow$ Réponse strictement intra-textuelle. « Quelle illusion » $\rightarrow$ Chercher une fausse croyance initiale.
\par\medskip
\textbf{Réponse type :}
L'illusion identifiée par l'auteur est de croire que \textbf{la technique est un « moyen neutre, un instrument docile à notre volonté »} (l. 1). L'homme a pensé qu'il pouvait utiliser la technique sans en être transformé, comme un simple outil passif.
\par\medskip
\emph{Commentaire prof :} Reformulation de la question en amorce. Citation courte exacte (l. 1). Explication immédiate en ses mots (« outil passif »). 3/3 points assurés.

\par\medskip\noindent\textbf{Question 2 (Analyse - 4 pts)}
\par\medskip
\emph{Analyse de la consigne :} « Montrez... en vous appuyant sur le texte » $\rightarrow$ Analyse structurelle + citations. « Comment l'auteur structure » $\rightarrow$ Repérer les étapes logiques (constat, exemples, conséquence).
\par\medskip
\textbf{Réponse type :}
L'auteur structure sa démonstration en \textbf{trois temps logiques} :
\begin{enumerate}
    \item \textbf{Le constat d'échec :} Il oppose l'illusion initiale (« moyen neutre », l. 1) à la réalité historique (« l'histoire récente nous enseigne », l. 1-2).
    \item \textbf{L'argumentation par l'exemple concret (l'automatisation) :} Il détaille les effets réels non voulus : modification de la nature du travail, réduction de l'autonomie, soumission au rythme machine (l. 2-4). C'est une \textbf{argumentation par les faits} (induction) : les moyens (automatisation) ont produit des fins (aliénation, cadence) non choisies.
    \item \textbf{La conclusion systémique :} Il généralise : « L'efficacité devient la valeur suprême » (l. 4), le système technique « s'auto-entretient » (l. 5). Les moyens (la technique) ont inversé la relation : ils déterminent désormais les fins (l'efficacité, l'indifférence au destin humain).
\end{enumerate}
\par\medskip
\emph{Commentaire prof :} Structure claire (1, 2, 3). Vocabulaire technique (induction, argumentation par les faits, généralisation). Citations intégrées (l. 1, l. 2-4, l. 4-5). 4/4 points.

\par\medskip\noindent\textbf{Question 3 (Réflexion - 8 pts)}
\par\medskip
\emph{Analyse de la consigne :} « Discutez cette thèse » $\rightarrow$ Dialectique (Pour / Contre / Nuance). « En vous appuyant sur le texte ET sur des exemples concrets » $\rightarrow$ Obligation d'aller puiser dans le texte ET d'apporter sa culture personnelle (exemples camerounais bienvenus).
\par\medskip
\textbf{Plan de rédaction suggéré :}
\begin{itemize}
    \item \textbf{Amorce :} Reprise de la thèse du texte : l'inversion moyen/fin (l. 2) et l'autonomie du système technique (l. 5).
    \item \textbf{Thèse (Appui au texte) :} La technique impose sa logique. Exemple texte : automatisation $\rightarrow$ cadence. Exemple personnel : \textbf{L'algorithme de gestion des flux (ex. applications de livraison type Jumia Food / Glovo au Cameroun)} : le livreur est noté, géolocalisé, pénalisé par l'algo ; il subit la cadence, il ne la choisit pas. La technique (l'algo) maîtrise l'homme.
    \item \textbf{Antithèse (Limites / Reprise de main) :} L'homme reste concepteur, législateur, utilisateur critique. Exemple : \textbf{La régulation de la monnaie mobile (Mobile Money) au Cameroun (BEAC/COBAC)} : l'État encadre la technique financière pour protéger l'usager (plafonds, KYC). La technique est remise dans le cadre du droit (fin humaine).
    \item \textbf{Synthèse / Nuance :} La maîtrise n'est pas binaire (0 ou 1). C'est une \textbf{lutte permanente} (Hans Jonas : « principe responsabilité »). La technique n'est pas un destin, mais un \textbf{défi éthique et politique}. Il faut une « éthique de la technique » (formation, régulation, design éthique) pour que l'homme reste le « berger de l'être » (Heidegger) et non le berger de la machine.
\end{itemize}

\par\medskip
\textbf{Extrait de rédaction attendue (Début de la thèse) :}
\og Le texte a raison de souligner que « les moyens finissent par déterminer les fins » (l. 2). Prenons l'exemple des plateformes de livraison à Douala ou Yaoundé. Le livreur, censé être l'utilisateur de l'application, se trouve en réalité \textbf{soumis à la cadence algorithmique} : l'optimisation du trajet, la notation client, la pénalité de retard sont imposées par le code informatique. Ici, le moyen (l'algorithme logistique) dicte la fin (la rentabilité de la plateforme) au détriment de la fin humaine (dignité du travail, autonomie). Cela illustre parfaitement l'idée du texte selon laquelle « le système technique s'auto-entretient, indifférent aux destins humains » (l. 5-6). \fg{}
\par\medskip
\emph{Commentaire prof :} Ancre textuelle (citations l. 2 et l. 5-6). Exemple concret, localisé, précis (Douala/Yaoundé, algorithme). Articulation logique (Illustration $\rightarrow$ Confirmation thèse). C'est ce niveau d'exigence qui vaut les 7-8/8.

\end{exoresolu}

\begin{aretenir}[Check-list avant de rendre la copie (Questions)]
\begin{itemize}
    \item [$\square$] Ai-je \textbf{numéroté} mes réponses exactement comme les questions (1, 2, 3...) ?
    \item [$\square$] Chaque réponse commence-t-elle par une \textbf{reformulation de la question} ?
    \item [$\square$] Mes citations sont-elles \textbf{courtes}, entre \og guillemets \fg{}, avec \textbf{numéro de ligne} ?
    \item [$\square$] Ai-je \textbf{expliqué} chaque citation (pas de citation orpheline) ?
    \item [$\square$] Pour les questions « Selon le texte » : \textbf{Zéro exemple personnel}, zéro connaissance extérieure ?
    \item [$\square$] Pour la question de réflexion : \textbf{Thèse / Antithèse / Synthèse} visible ? Exemples concrets (Cameroun si possible) ? Lien final avec le texte ?
    \item [$\square$] Orthographe, syntaxe, lisibilité (saut de ligne entre chaque réponse) ?
\end{itemize}
\end{aretenir}

\courssub{Transition vers l'Essai personnel}

Vous maîtrisez désormais l'art de répondre aux questions : vous savez lire, analyser, citer, argumenter. Mais l'exercice sur texte ne s'arrête pas là. La dernière épreuve, l'\textbf{Essai personnel}, demande de faire un pas de plus : quitter le texte pour construire \textbf{votre propre pensée}, tout en gardant la rigueur philosophique acquise ici. C'est l'objet de la prochaine section.

\coursec{L'essai personnel}

Après avoir répondu aux questions de lecture, qui vérifient la compréhension exacte du texte, le candidat aborde la dernière et la plus importante partie de l'exercice sur texte : l'\cle{essai personnel}. C'est ici que l'examinateur évalue la capacité de l'élève à \emph{penser par lui-même}. Les réponses aux questions montraient que vous avez compris l'auteur ; l'essai personnel montre que vous êtes capable de discuter avec lui, de le critiquer, de le dépasser. C'est l'exercice le plus proche de la dissertation, et il constitue souvent la moitié des points de l'épreuve au Probatoire et au Baccalauréat technique.

\courssub{Définition et finalité de l'essai personnel}

\begin{definition}[L'essai personnel]
L'\cle{essai personnel} est un texte argumentatif dans lequel le candidat discute, avec ses propres arguments, ses propres exemples et ses propres références philosophiques, le problème posé par le texte étudié. Il ne s'agit ni de résumer le texte, ni de répéter la thèse de l'auteur, mais de prendre position de manière réfléchie, nuancée et justifiée sur la question que le texte soulève.
\end{definition}

Pour bien comprendre la différence avec les réponses aux questions, prenons une image simple. Répondre aux questions, c'est comme décrire un plat que le cuisinier (l'auteur) a préparé : vous dites quels ingrédients il a mis et comment il a procédé. Rédiger l'essai personnel, c'est \emph{goûter le plat et donner votre avis argumenté} : est-il réussi ? Manque-t-il quelque chose ? Auriez-vous cuisiné autrement ? Le correcteur attend trois qualités :

\begin{itemize}
\item la \cle{personnalité} : vous parlez en votre nom (\og je pense que\fg{}, \og il me semble\fg{}), sans pour autant raconter votre vie ;
\item l'\cle{esprit critique} : vous examinez les forces et les faiblesses de la thèse de l'auteur, vous envisagez les objections ;
\item l'\cle{argumentation} : chaque jugement est soutenu par un raisonnement, un exemple et, si possible, une référence philosophique étudiée dans l'année.
\end{itemize}

\begin{attention}[L'erreur la plus fréquente]
L'erreur la plus sanctionnée consiste à \textbf{répéter la thèse de l'auteur sans la discuter}. Beaucoup de candidats recopient, avec d'autres mots, ce que le texte a déjà dit : le correcteur n'y trouve alors aucune pensée personnelle et attribue une note très faible, même si la copie est longue. Retenez cette règle d'or : l'essai doit être \cle{personnel}, \cle{critique} et \cle{argumenté}. L'auteur du texte est votre partenaire de dialogue, pas votre maître à penser.
\end{attention}

\begin{exemplebox}[Un essai réussi, c'est quelle longueur ?]
Au Baccalauréat technique camerounais, un essai personnel réussi compte en général \textbf{2 à 3 pages} de copie, soit environ \textbf{25 à 40 lignes}. En dessous de 20 lignes, la réflexion est presque toujours jugée insuffisamment développée ; au-delà de 4 pages, le candidat risque de se disperser et de manquer de temps pour relire. La qualité de l'argumentation prime toujours sur la quantité.
\end{exemplebox}

\courssub{La structure de l'essai personnel}

Comme toute dissertation philosophique, l'essai personnel obéit à une architecture en trois moments : une \cle{introduction}, un \cle{développement} en deux ou trois parties, et une \cle{conclusion}. Cette structure n'est pas une formalité arbitraire : elle reproduit le mouvement naturel de la pensée qui pose un problème, l'examine sous plusieurs angles, puis tranche.

\begin{popfigure}
\begin{center}
\begin{tikzpicture}[
  node distance=0.55cm,
  bloc/.style={draw=popblue, thick, fill=popblueL, rounded corners=3pt, text width=10.2cm, align=left, inner sep=6pt, font=\small},
  titre/.style={font=\small\bfseries, text=popdark},
  fleche/.style={-{Stealth[length=2.5mm]}, very thick, poporange}
]
\node[bloc] (intro) {\textbf{INTRODUCTION} \;---\; Amorce (accroche concrète) $\rightarrow$ présentation du problème $\rightarrow$ \textbf{problématique} $\rightarrow$ annonce du plan};
\node[bloc, below=of intro] (p1) {\textbf{PARTIE 1} \;---\; Première idée : accord nuancé avec la thèse du texte (argument + exemple + référence)};
\node[bloc, below=of p1] (p2) {\textbf{PARTIE 2} \;---\; Deuxième idée : discussion, objection, limite de la thèse (argument + exemple + référence)};
\node[bloc, below=of p2] (p3) {\textbf{PARTIE 3 (éventuelle)} \;---\; Dépassement : votre position synthétique (argument + exemple + référence)};
\node[bloc, below=of p3, fill=popgreenL, draw=popgreen] (ccl) {\textbf{CONCLUSION} \;---\; Bilan de la réflexion $\rightarrow$ réponse claire à la problématique $\rightarrow$ ouverture éventuelle};
\draw[fleche] (intro) -- (p1);
\draw[fleche] (p1) -- (p2);
\draw[fleche] (p2) -- (p3);
\draw[fleche] (p3) -- (ccl);
\end{tikzpicture}
\end{center}
\legende{Structure générale de l'essai personnel : chaque flèche correspond à une transition rédigée (connecteur logique) entre les parties.}
\end{popfigure}

\courssub{L'introduction : poser le problème}

L'introduction est la première chose que lit le correcteur : elle donne le ton de toute la copie. Elle comporte quatre étapes, dans l'ordre.

\begin{enumerate}
\item \textbf{L'amorce (ou accroche).} Une ou deux phrases qui partent d'un fait général, d'une expérience commune ou d'un débat actuel, et qui mènent naturellement vers le sujet. Évitez les amorces vides du type \og Depuis la nuit des temps, l'homme...\fg{}. Préférez un point de départ précis : une scène de la vie quotidienne, un débat public camerounais, un usage numérique.
\item \textbf{La présentation du problème.} Vous montrez en quoi le texte étudié soulève une difficulté : une tension entre deux idées, une opinion commune qu'il faut interroger. C'est le moment de mobiliser le texte : \og C'est précisément le problème que soulève l'auteur lorsqu'il affirme que...\fg{}.
\item \textbf{La problématique.} Une question unique, claire, formulée en une phrase interrogative, qui annonce l'enjeu de votre réflexion. Elle reprend le problème du texte mais avec vos mots.
\item \textbf{L'annonce du plan.} Vous indiquez brièvement les étapes de votre réflexion : \og Nous examinerons d'abord..., puis nous verrons que..., avant de montrer que...\fg{}.
\end{enumerate}

\begin{exemplebox}[Une introduction rédigée (sujet : la liberté)]
\og Au marché comme sur les réseaux sociaux, chacun revendique le droit de faire ce qu'il veut : la liberté semble être le bien le plus précieux et le plus réclamé de notre époque. Mais suffit-il de suivre tous ses désirs pour être libre ? Le texte étudié soutient que la véritable liberté suppose la maîtrise de soi, et non l'obéissance aux passions. Cette thèse, exigeante, mérite d'être discutée : \textbf{la liberté consiste-t-elle à faire ce que l'on veut ?} Nous verrons d'abord que l'opinion commune identifie liberté et absence de contrainte ; nous montrerons ensuite que cette conception est intenable ; nous défendrons enfin l'idée d'une liberté qui se construit par la raison.\fg{}
\end{exemplebox}

\courssub{Le développement : deux ou trois parties argumentées}

Le développement est le cœur de l'essai. Il comporte en général \textbf{deux ou trois parties}, qui correspondent à des étapes de la réflexion et non à de simples morceaux de texte.

\begin{methode}[Construire une partie du développement]
Chaque partie obéit à la formule : \textbf{une partie = une idée = un argument + un exemple + une référence philosophique}.
\begin{enumerate}
\item Formulez \textbf{l'idée directrice} de la partie en une phrase (c'est la petite thèse de la partie).
\item Développez \textbf{l'argument} : un raisonnement logique qui justifie cette idée (enchaînement de causes et de conséquences, analyse d'une notion, réfutation d'une objection).
\item Illustrez par \textbf{un exemple concret}, de préférence tiré de la vie courante camerounaise : débats publics, usages numériques, vie scolaire, traditions.
\item Appuyez-vous sur \textbf{une référence philosophique} étudiée dans l'année (un auteur, une doctrine, une notion du programme), citée brièvement et correctement.
\item Terminez par une \textbf{phrase de transition} qui prépare la partie suivante.
\end{enumerate}
\end{methode}

Le mouvement classique et le plus sûr pour un essai personnel est le suivant :

\begin{itemize}
\item \textbf{Partie 1 — Accord nuancé.} Vous montrez d'abord ce que la thèse de l'auteur (ou l'opinion commune) a de vrai et de convaincant. Cela prouve que vous l'avez comprise.
\item \textbf{Partie 2 — Discussion.} Vous soulevez une objection, une limite, un cas où la thèse semble mise en difficulté. C'est le moment critique de l'essai.
\item \textbf{Partie 3 — Dépassement.} Vous proposez votre propre position, qui dépasse la simple opposition : souvent une distinction (par exemple entre deux sens d'un même mot) permet de concilier ce que les deux premières parties avaient de juste.
\end{itemize}

\begin{attention}[Pièges du développement]
\begin{itemize}
\item Ne multipliez pas les idées dans une même partie : une partie, une seule idée.
\item Une référence philosophique mal comprise vaut moins qu'un bon exemple bien analysé : ne citez jamais un auteur que vous ne maîtrisez pas.
\item L'exemple ne remplace pas l'argument : il l'illustre. Annoncez-le (\og Ainsi...\fg{}, \og On le voit bien lorsque...\fg{}) et tirez-en la leçon.
\item Restez dans le sujet : chaque partie doit répondre, à sa manière, à la problématique annoncée.
\end{itemize}
\end{attention}

\courssub{Mobiliser les notions de l'unité et les références de l'année}

L'essai personnel n'est pas une conversation libre : c'est un exercice scolaire qui doit montrer que vous savez \emph{appliquer les notions de l'unité à des situations concrètes}. Avant de rédiger, faites mentalement l'inventaire de ce que le cours vous a appris sur le thème du texte : définitions des notions, distinctions importantes, thèses des auteurs étudiés. Ces acquis sont votre boîte à outils.

\begin{itemize}
\item Si le texte porte sur la \cle{liberté}, mobilisez la distinction entre la liberté comme indépendance (faire ce qu'on veut) et la liberté comme autonomie (se donner à soi-même sa loi), ainsi que les analyses de la conscience et du devoir vues dans l'année.
\item Si le texte porte sur la \cle{vérité}, rappelez la différence entre l'opinion et le savoir, et le rôle de la démonstration.
\item Si le texte porte sur l'\cle{État} ou la \cle{justice}, distinguez la loi positive (les règles votées) et l'idée de justice qui permet de les juger.
\end{itemize}

Les exemples concrets, eux, doivent être puisés dans votre expérience et dans l'actualité : un débat à la radio ou sur les réseaux sociaux, une règle de vie au lycée, une coutume familiale ou villageoise, une discussion sur la gouvernance ou sur la jeunesse. L'exemple local, bien choisi et bien analysé, prouve au correcteur que la philosophie n'est pas pour vous une matière abstraite.

\begin{savaistu}[Ce que valorisent les correcteurs camerounais]
Les correcteurs du Baccalauréat camerounais valorisent explicitement les exemples tirés du contexte local — débats sur la gouvernance, place de la jeunesse, usages des réseaux sociaux, tensions entre modernité et traditions — \emph{lorsqu'ils illustrent correctement l'argument}. Un exemple camerounais bien analysé vaut mieux qu'une référence étrangère récitée sans être comprise. Attention cependant : l'exemple local doit rester un \emph{exemplebox}, au service d'un raisonnement général, et non devenir un simple commentaire d'actualité ou un règlement de comptes.
\end{savaistu}

\courssub{Exemple de plan détaillé rédigé}

Voici, à titre de modèle, un plan détaillé complet sur la question : \og \textbf{La liberté consiste-t-elle à faire ce que l'on veut ?}\fg{} — question typique qu'un texte sur la liberté peut soulever à l'exercice sur texte.

\begin{exemplebox}[Plan détaillé : \og La liberté consiste-t-elle à faire ce que l'on veut ?\fg]
\textbf{Introduction.} Amorce : chacun réclame la liberté (expression, déplacement, choix de vie). Problème : l'opinion identifie liberté et absence de toute contrainte ; mais celui qui cède à tous ses désirs est-il vraiment libre ? Problématique : la liberté consiste-t-elle à faire ce que l'on veut ? Plan : I. l'opinion commune ; II. sa critique ; III. la liberté comme maîtrise de soi.

\textbf{Partie 1 — Faire ce que l'on veut semble être la définition même de la liberté.}
Argument : être libre, c'est d'abord ne pas subir la volonté d'autrui ; la contrainte extérieure (prison, interdit, censure) est l'ennemie évidente de la liberté. Exemple : l'élève qui, une fois le portail du lycée franchi, se sent \og libre\fg{} parce qu'il échappe au règlement ; ou l'internaute qui revendique la liberté d'expression sur les réseaux sociaux. Référence : la conception classique de la liberté comme indépendance à l'égard d'autrui. Transition : mais cette liberté-là suffit-elle à rendre un homme réellement libre ?

\textbf{Partie 2 — Pourtant, faire tout ce que l'on veut peut être une forme d'esclavage.}
Argument : celui qui obéit à chacun de ses désirs ne se gouverne plus lui-même ; il devient l'esclave de ses passions, de ses habitudes, des influences extérieures (publicité, groupe, réseaux sociaux). L'absence de contrainte extérieure cache alors une contrainte intérieure plus redoutable. Exemple : le jeune qui passe ses nuits sur son téléphone \og parce qu'il en a envie\fg{} et qui n'arrive plus à travailler ; le consommateur endetté qui achète sous l'effet de la publicité. Référence : l'idée, présente chez Platon et reprise par toute la tradition philosophique, que l'homme dominé par ses passions est un esclave qui se croit libre. Transition : il faut donc chercher une conception plus exigeante de la liberté.

\textbf{Partie 3 — La véritable liberté est autonomie : se donner à soi-même sa loi.}
Argument : être libre, ce n'est pas supprimer toute règle, c'est choisir la règle que l'on s'impose, par la raison et la volonté. La discipline n'est pas le contraire de la liberté : elle en est l'instrument. Exemple : le sportif ou l'apprenti artisan de Douala qui accepte un entraînement sévère pour devenir capable ; le citoyen qui respecte une loi qu'il juge juste. Référence : la notion d'autonomie (se donner à soi-même sa loi) étudiée dans le cours sur la liberté et le devoir. 

\textbf{Conclusion.} Bilan : l'opinion avait raison de voir dans la contrainte extérieure un obstacle à la liberté, mais tort d'en conclure que la liberté est le règne du caprice. Réponse à la problématique : la liberté ne consiste pas à faire ce que l'on veut, mais à vouloir ce que l'on fait, en connaissance de cause. Ouverture : cette liberté intérieure suffit-elle lorsque les conditions sociales et économiques pèsent sur les choix ?
\end{exemplebox}

\courssub{La conclusion : trancher et ouvrir}

La conclusion est souvent négligée, alors qu'elle laisse au correcteur sa dernière impression. Elle comporte trois moments :

\begin{enumerate}
\item \textbf{Le bilan de la réflexion.} Rappelez en deux ou trois phrases le chemin parcouru : ce que chaque partie a établi, et comment la discussion a fait progresser le problème.
\item \textbf{La réponse claire à la problématique.} C'est le point essentiel : vous devez répondre \emph{explicitement} à la question posée dans l'introduction, par oui, par non, ou — le plus souvent en philosophie — par une réponse nuancée mais ferme (\og la liberté ne consiste pas à faire ce que l'on veut, mais...\fg{}). Une conclusion qui ne tranche pas est une conclusion manquée.
\item \textbf{L'ouverture éventuelle.} Vous pouvez terminer par une question nouvelle qui prolonge la réflexion, à condition qu'elle découle logiquement de votre sujet. L'ouverture est un plus, non une obligation : mieux vaut une conclusion nette sans ouverture qu'une ouverture artificielle.
\end{enumerate}

\courssub{Les qualités rédactionnelles}

Un essai philosophiquement riche mais mal écrit perd une grande partie de sa valeur. Le correcteur juge aussi la langue et la présentation.

\begin{itemize}
\item \textbf{Clarté.} Phrases courtes, une idée par phrase. Relisez chaque phrase en vous demandant : \og un lecteur qui n'a pas le texte sous les yeux comprend-il ?\fg{}
\item \textbf{Orthographe et grammaire.} Les fautes répétées (accords du participe passé, homophones \og a/à\fg{}, \og se/ce\fg{}, \og son/sont\fg{}) fatiguent le correcteur et font baisser la note. Gardez cinq minutes pour la relecture.
\item \textbf{Paragraphes liés par des connecteurs logiques.} Chaque partie et chaque étape doivent être enchaînées par des mots de liaison qui indiquent le sens du raisonnement :
\begin{itemize}
\item pour ajouter : \emph{de plus, en outre, également} ;
\item pour opposer : \emph{cependant, pourtant, en revanche, néanmoins} ;
\item pour conclure : \emph{donc, ainsi, par conséquent, c'est pourquoi} ;
\item pour ordonner : \emph{d'abord, ensuite, enfin}.
\end{itemize}
\item \textbf{Présentation.} Alinéa à chaque partie, pas de ratures excessives, écriture lisible.
\end{itemize}

\begin{exoresolu}[Transformer un paragraphe faible en paragraphe d'essai]
\textbf{Énoncé.} Voici le paragraphe d'un candidat sur le sujet \og La liberté consiste-t-elle à faire ce que l'on veut ?\fg{} : \og L'auteur dit que la liberté ce n'est pas faire ce qu'on veut. Il a raison. Moi aussi je pense pareil. Par exemple les jeunes qui font n'importe quoi sur les réseaux sociaux ne sont pas libres. Donc la liberté c'est autre chose.\fg{} Identifiez les défauts de ce paragraphe et réécrivez-le selon la méthode de l'essai personnel.
\tcblower
\textbf{Solution détaillée.} Ce paragraphe présente quatre défauts : (1) il se contente de \emph{répéter} la thèse de l'auteur sans la discuter ; (2) l'accord (\og il a raison, moi aussi je pense pareil\fg{}) n'est justifié par aucun argument ; (3) l'exemple est jeté sans analyse ; (4) le style est oral et sans connecteurs logiques.

Réécriture possible : \og Cependant, la conception commune de la liberté se révèle intenable dès qu'on l'examine de près. En effet, celui qui suit chacun de ses désirs ne se gouverne plus lui-même : il obéit à des forces qu'il ne contrôle pas, et l'absence de contrainte extérieure masque alors une servitude intérieure. On le voit bien avec certains usages des réseaux sociaux au Cameroun : le lycéen qui consulte son téléphone toute la nuit \og parce qu'il en a envie\fg{} croit exercer sa liberté, alors qu'il est en réalité devenu incapable de s'en détacher ; son envie commande à sa place. C'est précisément ce que la tradition philosophique, de Platon aux moralistes modernes, appelle l'esclavage des passions : l'homme le moins libre est souvent celui qui se croit le plus libre. Il faut donc rechercher, au-delà de la simple indépendance, une liberté qui se construit par la maîtrise de soi.\fg{} Ce paragraphe contient une idée (la critique de l'opinion), un argument (l'obéissance aux désirs est une servitude cachée), un exemple local analysé et une référence, le tout articulé par des connecteurs (\og cependant\fg{}, \og en effet\fg{}, \og on le voit bien\fg{}, \og il faut donc\fg{}).
\end{exoresolu}

\begin{exercice}[À vous de rédiger]
Un texte soutient que \og la tradition est un trésor qu'il faut transmettre sans y toucher\fg{}. Rédigez l'introduction complète (amorce, problème, problématique, annonce du plan) et le plan détaillé (trois parties : accord nuancé, discussion, dépassement) d'un essai personnel sur ce sujet, en prévoyant pour chaque partie un argument, un exemple camerounais (coutumes, vie familiale, école, numérique) et une référence du cours.
\end{exercice}

\begin{corrige}[Exercice : la tradition]
\textbf{Introduction possible.} Amorce : au village comme en ville, les aînés rappellent sans cesse aux jeunes l'importance des coutumes (salutations, rites, langues). Problème : la tradition transmet des valeurs précieuses, mais faut-il la recevoir sans examen, au risque de conserver aussi ce qu'elle a d'injuste ou de dépassé ? Problématique : la fidélité à la tradition exige-t-elle de la transmettre sans y toucher ? Plan : I. la tradition comme trésor ; II. les limites de la transmission sans examen ; III. une fidélité critique.

\textbf{Partie 1.} Idée : la tradition est bien un trésor, car elle transmet l'expérience accumulée des générations et donne à chacun une identité. Argument : sans héritage, chaque génération devrait tout recommencer ; la mémoire collective est une richesse. Exemple : les langues locales, les contes, les rites de solidarité villageoise (tontines, travaux collectifs). Référence : la notion de culture comme héritage, vue dans le cours.

\textbf{Partie 2.} Idée : cependant, transmettre \og sans y toucher\fg{} est impossible et dangereux. Argument : toute tradition contient des éléments injustes ou inadaptés ; les conserver par simple habitude, c'est renoncer à juger. Exemple : certaines pratiques coutumières défavorisant les filles (mariages précoces, exclusion de l'héritage) sont de plus en plus contestées, y compris au sein des communautés. Référence : la distinction entre coutume et justice.

\textbf{Partie 3.} Idée : la vraie fidélité est critique : on honore la tradition en la transmettant transformée. Argument : une tradition vivante est celle que chaque génération réinterprète ; la figer, c'est la tuer. Exemple : les jeunes artistes camerounais qui réactualisent les rythmes traditionnels dans la musique urbaine. Référence : le rôle de la raison dans le jugement des valeurs.

\textbf{Conclusion attendue.} Réponse nuancée : la tradition est un trésor, mais un trésor qu'il faut trier pour le transmettre. Ouverture possible : qui, dans la société, est légitime pour opérer ce tri ?
\end{corrige}

\begin{aretenir}[L'essentiel de l'essai personnel]
\begin{itemize}
\item L'essai personnel est un texte argumentatif où vous \textbf{discutez} le problème du texte avec vos propres arguments, exemples et références : jamais un résumé, jamais une répétition de l'auteur.
\item Structure : introduction (amorce, problème, problématique, plan), développement en 2 ou 3 parties (accord nuancé, discussion, dépassement), conclusion (bilan, réponse claire, ouverture éventuelle).
\item Chaque partie = une idée + un argument + un exemple concret + une référence philosophique de l'année.
\item Longueur visée : 2 à 3 pages, soit 25 à 40 lignes, rédigées clairement, avec des connecteurs logiques et une relecture soignée.
\end{itemize}
\end{aretenir}

\coursec{Application guidée : un sujet complet corrigé}

Nous savons désormais ce qu'est un essai personnel : une prise de position argumentée, structurée et personnelle, qui prolonge les réponses aux questions sans les répéter. Mais connaître les règles d'un exercice ne suffit pas : il faut s'entraîner sur un sujet complet, dans les conditions réelles de l'examen. Cette dernière séquence met en application tout ce que nous avons appris, étape par étape, sur un sujet type Probatoire technique. Suivez la démarche comme si vous étiez en salle d'examen : c'est en refaisant soi-même chaque étape que la méthode devient un réflexe.

\courssub{Présentation du sujet type Probatoire technique}

\begin{exemplebox}[Sujet d'exercice sur texte (type Probatoire technique)]
\textbf{Texte} : « Tout homme a une conscience. Mais que signifie ce mot ? La conscience est d'abord cette voix intérieure qui nous avertit, avant même que nous ayons agi, que l'acte que nous nous apprêtons à commettre est bon ou mauvais. Le voleur qui s'introduit dans une maison la nuit sait, au moment même où il agit, qu'il fait le mal : sa conscience le lui dit, et c'est pour l'étouffer qu'il se répète que tout le monde ferait pareil à sa place. Mais la conscience n'est pas seulement ce juge qui condamne ou approuve nos actes passés ou futurs. Elle est aussi cette présence à soi-même qui fait que nous savons que nous existons, que nous pensons, que nous souffrons. Un homme sans conscience ne serait pas un homme : il serait une machine. Ainsi, la conscience est à la fois le témoin de notre existence et le juge de nos actions. C'est pourquoi on peut dire qu'elle est ce qu'il y a de plus profondément humain en l'homme. Celui qui écoute sa conscience ne se trompe jamais complètement, car elle lui rappelle sans cesse la différence entre ce qu'il fait et ce qu'il devrait faire. »

\textbf{Questions} :
\begin{enumerate}
\item Dégagez le thème et la thèse de ce texte.
\item Quelles sont les deux dimensions de la conscience présentées par l'auteur ? Justifiez votre réponse par des citations du texte.
\item Expliquez la phrase suivante : « Un homme sans conscience ne serait pas un homme : il serait une machine. »
\item Selon l'auteur, pourquoi celui qui écoute sa conscience « ne se trompe jamais complètement » ?
\end{enumerate}

\textbf{Essai personnel} : La conscience morale suffit-elle pour bien agir ? Vous prendrez position sur cette question en vous appuyant sur le texte, sur le cours et sur des exemples tirés de votre expérience.
\end{exemplebox}

Ce sujet est représentatif de ce qui attend le candidat au Probatoire / Baccalauréat technique : un texte court (15 à 20 lignes) portant sur une notion du programme, quatre questions de difficulté croissante, puis une consigne d'essai personnel qui invite à dépasser le texte.

\begin{savaistu}[Les sujets d'exercice sur texte à l'examen]
Les sujets d'exercice sur texte du Probatoire / Baccalauréat technique portent le plus souvent sur les notions du programme de Première : la \cle{conscience}, l'\cle{inconscient}, la \cle{liberté}, le \cle{devoir}, le \cle{bonheur}, l'\cle{État}. Réviser ces notions, c'est donc se préparer directement à l'exercice sur texte. Un élève qui maîtrise les définitions et les thèses des auteurs étudiés en classe possède déjà l'essentiel des outils pour réussir.
\end{savaistu}

\courssub{Étape 1 : la lecture analytique collective}

Avant toute rédaction, la classe procède à une lecture analytique du texte, selon la méthode étudiée dans la section précédente. Trois questions guident cette lecture : de quoi parle le texte (le \cle{thème}) ? Que soutient l'auteur (la \cle{thèse}) ? Quel problème le texte pose-t-il (la \cle{problématique}) ?

\begin{methode}[Dégager thème, thèse et problématique]
\begin{enumerate}
\item \textbf{Première lecture silencieuse} : lire le texte en entier, sans écrire, pour saisir le sens global.
\item \textbf{Deuxième lecture active} : souligner les mots-clés (ici : « conscience », « voix intérieure », « juge », « témoin de notre existence ») et entourer les articulations logiques (« d'abord », « mais », « ainsi », « c'est pourquoi »).
\item \textbf{Formuler le thème} en un ou deux mots : ici, \emph{la conscience}.
\item \textbf{Formuler la thèse} en une phrase complète : ici, \emph{la conscience est à la fois le témoin de notre existence et le juge de nos actions ; elle est ce qu'il y a de plus profondément humain en l'homme}.
\item \textbf{Formuler la problématique} sous forme de question : \emph{qu'est-ce que la conscience et pourquoi est-elle le propre de l'homme ?}
\end{enumerate}
\end{methode}

Appliquée collectivement, cette méthode permet de vérifier que tous les élèves ont compris le texte de la même manière. Le professeur note au tableau les propositions des élèves, les confronte, et la classe retient la formulation la plus exacte. Cette étape dure environ dix minutes : c'est un investissement rentable, car une erreur de lecture à ce stade fausserait toutes les réponses.

\courssub{Étape 2 : rédaction guidée des réponses aux questions 1 et 2}

La rédaction des réponses obéit à une règle simple : répondre d'abord, justifier ensuite, citer toujours. Chaque réponse doit être une phrase complète, rédigée, et non un mot isolé.

\begin{exoresolu}[Rédaction guidée des questions 1 et 2]
\textbf{Énoncé} : Rédiger les réponses aux questions 1 et 2 du sujet ci-dessus.
\tcblower
\textbf{Solution détaillée.}

\emph{Question 1} : Le thème de ce texte est \textbf{la conscience}. La thèse de l'auteur est que la conscience possède une double fonction : elle est à la fois « le témoin de notre existence » et « le juge de nos actions », et c'est elle qui constitue « ce qu'il y a de plus profondément humain en l'homme ».

\emph{Question 2} : L'auteur distingue deux dimensions de la conscience. La première est la \textbf{conscience morale} : elle est « cette voix intérieure qui nous avertit, avant même que nous ayons agi, que l'acte que nous nous apprêtons à commettre est bon ou mauvais » ; elle juge nos actes, comme le montre l'exemple du voleur qui « sait, au moment même où il agit, qu'il fait le mal ». La seconde est la \textbf{conscience psychologique} (ou conscience de soi) : elle est « cette présence à soi-même qui fait que nous savons que nous existons, que nous pensons, que nous souffrons ». La première dimension concerne nos actions, la seconde concerne notre existence même.
\end{exoresolu}

Remarquez la structure de chaque réponse : une phrase qui répond directement à la question, puis une ou deux citations entre guillemets qui prouvent que la réponse est bien dans le texte. C'est cette fidélité au texte que le correcteur attend en premier lieu.

\begin{attention}[Erreur fréquente : négliger la relecture]
Beaucoup de candidats rendent leur copie dès la dernière phrase écrite, sans relire. C'est une erreur coûteuse : \textbf{dix minutes de relecture} permettent de corriger les fautes d'orthographe, de compléter une phrase inachevée et surtout de découvrir une question oubliée. Prévoyez toujours ce temps de relecture dans votre gestion de l'épreuve : il vaut souvent un ou deux points.
\end{attention}

\begin{exemplebox}[Le coût d'une question non traitée]
Sur un sujet comportant quatre questions, chaque question non traitée fait perdre en moyenne \textbf{1,5 à 2 points sur 20}. Un candidat qui néglige deux questions part donc avec un handicap de 3 à 4 points : même excellent sur le reste, il ne peut plus viser la moyenne confortable. Mieux vaut une réponse courte et honnête à chaque question qu'une réponse longue à une seule. Traitez toutes les questions, même partiellement.
\end{exemplebox}

\courssub{Étape 3 : construction collective du plan de l'essai personnel}

L'essai personnel demande : « La conscience morale suffit-elle pour bien agir ? » La formulation « suffit-elle » invite à une réponse nuancée : la conscience morale est nécessaire, mais peut-être pas suffisante. La classe construit ensemble le plan au tableau.

\begin{methode}[Construire le plan de l'essai personnel]
\begin{enumerate}
\item \textbf{Analyser la consigne} : identifier la notion interrogée (la conscience morale) et le type de question (« suffit-elle » appelle une discussion).
\item \textbf{Première partie} : montrer ce que la conscience morale apporte. Elle nous avertit du bien et du mal (le texte), elle nous rend responsables de nos actes. Sans elle, l'action serait aveugle. Exemple : l'écolier de Yaoundé qui rend le portefeuille trouvé dans la rue parce que sa conscience le lui commande.
\item \textbf{Deuxième partie} : montrer les limites de la conscience morale. Elle peut être mal formée, étouffée par l'habitude ou les préjugés ; le voleur du texte « l'étouffe » en se répétant que « tout le monde ferait pareil ». De plus, savoir le bien ne suffit pas : il faut encore la volonté de le faire.
\item \textbf{Conclusion} : prendre position. La conscience morale est une condition nécessaire du bien agir, mais elle doit être éclairée par la raison, éduquée par la société et accompagnée par la volonté.
\end{enumerate}
\end{methode}

Ce plan en deux parties (thèse puis limites, ou l'inverse) est le plus sûr pour une question de type « suffit-il » : il permet de traiter les deux faces du problème avant de conclure par une position personnelle et argumentée.

\courssub{Étape 4 : rédaction individuelle de l'introduction et de la première partie}

Chaque élève rédige ensuite seul, en vingt minutes, l'introduction et la première partie de l'essai. L'introduction suit le schéma classique : amorce, présentation du problème, annonce de la problématique et du plan.

\begin{exemplebox}[Exemple d'introduction rédigée]
Nous entendons souvent dire : « écoute ta conscience ». Cette injonction suppose que la conscience morale est un guide sûr, capable de nous indiquer en toute circonstance ce qu'il faut faire. Mais cette confiance est-elle justifiée ? Le texte étudié présente la conscience comme « cette voix intérieure qui nous avertit » du bien et du mal, et affirme que celui qui l'écoute « ne se trompe jamais complètement ». Faut-il en conclure qu'elle suffit pour bien agir ? Nous montrerons d'abord que la conscience morale est un guide indispensable, puis qu'elle présente cependant des limites qui obligent à la compléter par la raison et la volonté.
\end{exemplebox}

La première partie développe ensuite un seul argument principal, illustré par un exemple et appuyé sur le texte ou le cours. Un paragraphe = une idée = un exemple : cette règle simple garantit la clarté de la copie.

\courssub{Correction croisée et grille d'évaluation}

Une fois la rédaction terminée, les élèves échangent leurs copies et les corrigent mutuellement à l'aide d'une grille inspirée des critères officiels de correction. Cette correction croisée est formatrice à double titre : celui qui corrige apprend à repérer les erreurs, et celui qui est corrigé reçoit un regard neuf sur son travail.

\begin{popfigure}
\begin{center}
\renewcommand{\arraystretch}{1.4}
\begin{tabularx}{\linewidth}{|l|X|X|X|}
\hline
\rowcolor{popblueL} \textbf{Critère} & \textbf{Maîtrise insuffisante (0--0,5)} & \textbf{Maîtrise moyenne (1--1,5)} & \textbf{Bonne maîtrise (2)} \\
\hline
\cellcolor{popcream}\textbf{Fidélité au texte} & Réponses hors sujet, aucune citation, texte ignoré & Réponses partielles, citations rares ou mal choisies & Réponses exactes, appuyées par des citations pertinentes \\
\hline
\cellcolor{popcream}\textbf{Exactitude} & Erreurs sur les notions, confusion entre conscience morale et conscience de soi & Notions globalement comprises mais définitions imprécises & Notions définies correctement, vocabulaire philosophique juste \\
\hline
\cellcolor{popcream}\textbf{Cohérence} & Essai sans plan, contradictions, répétitions & Plan visible mais parties déséquilibrées & Introduction, développement structuré, conclusion qui prend position \\
\hline
\cellcolor{popcream}\textbf{Langue} & Fautes nombreuses gênant la lecture, phrases inachevées & Quelques fautes, expression parfois maladroite & Orthographe correcte, phrases claires et bien construites \\
\hline
\end{tabularx}
\end{center}
\legende{Grille d'évaluation simplifiée pour l'entraînement (4 critères $\times$ 2 points = 8 points), inspirée des critères officiels : fidélité au texte, exactitude, cohérence, langue.}
\end{popfigure}

Pour rendre la correction concrète, le correcteur annote la copie dans la marge : il signale les points forts par une mention positive et les erreurs par un renvoi précis au critère concerné.

\begin{popfigure}
\begin{center}
\begin{tikzpicture}
% Feuille de copie
\draw[fill=white,draw=popdark,thick] (0,0) rectangle (7.2,8.6);
% Marge
\draw[popline,dashed] (5.3,0) -- (5.3,8.6);
% Lignes d'écriture simulées
\foreach \y in {7.9,7.5,7.1,6.7,6.3,5.9,5.5,5.1,4.7,4.3,3.9,3.5,3.1,2.7,2.3,1.9,1.5,1.1,0.7}{
  \draw[popmuted!40] (0.3,\y) -- (5.1,\y);
}
% Titre de la copie
\node[anchor=west,font=\small\bfseries] at (0.3,8.3) {\textcolor{popink}{Essai personnel : brouillon d'élève}};
% Passage fort surligné
\draw[popgreen,line width=5pt,opacity=0.35] (0.3,6.7) -- (5.1,6.7);
\draw[popgreen,line width=5pt,opacity=0.35] (0.3,6.3) -- (5.1,6.3);
% Passage fautif souligné
\draw[poporange,thick,decorate,decoration={snake,amplitude=0.4pt}] (0.3,3.85) -- (5.1,3.85);
\draw[poporange,thick,decorate,decoration={snake,amplitude=0.4pt}] (0.3,3.45) -- (5.1,3.45);
% Annotations marginales
\node[anchor=west,font=\scriptsize,text=popgreen!60!black,align=left] at (5.45,6.5) {\textbf{Point fort :}\\ citation exacte\\ du texte, bien\\ introduite (critère\\ fidélité : 2/2)};
\node[anchor=west,font=\scriptsize,text=poporange!80!black,align=left] at (5.45,3.6) {\textbf{Erreur :}\\ confusion entre\\ conscience morale\\ et conscience\\ de soi (critère\\ exactitude : 1/2)};
\node[anchor=west,font=\scriptsize,text=poppurple,align=left] at (5.45,1.2) {\textbf{Remarque :}\\ conclusion absente,\\ penser à prendre\\ position (critère\\ cohérence)};
% Flèches d'annotation
\draw[-{Stealth},popgreen!60!black] (5.4,6.5) -- (5.15,6.5);
\draw[-{Stealth},poporange!80!black] (5.4,3.6) -- (5.15,3.6);
\draw[-{Stealth},poppurple] (5.4,1.2) -- (4.6,0.9);
\end{tikzpicture}
\end{center}
\legende{Photocopie annotée d'une copie d'élève : en vert, un point fort (citation exacte du texte) ; en orange, une erreur de fond (confusion entre les deux dimensions de la conscience) ; en violet, une remarque de structure (conclusion manquante). Chaque annotation renvoie à un critère de la grille.}
\end{popfigure}

\begin{methode}[S'auto-corriger : les trois questions de la relecture]
Avant de rendre votre copie, relisez-la en vous posant trois questions, dans l'ordre :
\begin{enumerate}
\item \textbf{Ai-je répondu à toutes les questions ?} Vérifiez en cochant chaque question du sujet : une question sans réponse, c'est 1,5 à 2 points perdus.
\item \textbf{Ai-je cité le texte ?} Chaque réponse doit contenir au moins une citation courte entre guillemets. Sans citation, votre réponse n'est pas prouvée.
\item \textbf{Mon essai est-il personnel et structuré ?} Vérifiez la présence d'une introduction, de deux parties distinctes, d'exemples qui vous appartiennent et d'une conclusion où vous prenez clairement position.
\end{enumerate}
Si l'une de ces trois questions reçoit une réponse négative, corrigez immédiatement : il est encore temps.
\end{methode}

\courssub{Bilan : auto-évaluation des points forts et des erreurs à corriger}

La dernière étape de l'application guidée est le bilan personnel. Chaque élève rédige, sur une demi-page, son auto-évaluation : quels sont mes points forts (par exemple : « je cite bien le texte », « mon introduction est claire ») ? Quelles erreurs dois-je corriger (par exemple : « je confonds encore conscience morale et conscience de soi », « j'oublie la conclusion », « je perds du temps sur la question 1 et je bâcle la question 4 ») ?

Cet exercice d'auto-évaluation n'est pas une formalité : c'est lui qui transforme une simple correction en progrès durable. Un élève qui sait nommer ses erreurs est déjà à moitié guéri de ses erreurs. Conservez votre grille annotée et votre bilan dans votre cahier : avant l'examen, relisez-les pour ne pas reproduire les mêmes fautes.

\begin{aretenir}[L'essentiel de l'application guidée]
\begin{itemize}
\item Un sujet type comporte un texte court, quatre questions et une consigne d'essai personnel.
\item La lecture analytique (thème, thèse, problématique) conditionne la réussite de tout l'exercice.
\item Chaque réponse suit le schéma : répondre, justifier, citer le texte.
\item L'essai personnel suit un plan en deux parties et se conclut par une position argumentée.
\item La correction croisée à l'aide d'une grille (fidélité au texte, exactitude, cohérence, langue) et l'auto-évaluation finale transforment l'entraînement en progrès réel.
\end{itemize}
\end{aretenir}

\begin{exercice}[Entraînement à la correction croisée]
Échangez votre essai avec votre voisin. À l'aide de la grille d'évaluation, attribuez une note sur 2 pour chacun des quatre critères, justifiez chaque note par une annotation marginale précise, puis rédigez votre propre auto-évaluation en cinq lignes : deux points forts et deux erreurs à corriger avant la prochaine évaluation.
\end{exercice}

\newpage

\coursec{Activité d'intégration}

\coursec{Leçon N21 : L'exercice sur texte}

\courssub{Objectifs de l'activité}
Cette activité d'intégration vise à \cle{mobiliser l'ensemble de la méthode de l'exercice sur texte} dans des conditions simulées d'examen (Probatoire technique). Elle permet aux élèves de :
\begin{itemize}
    \item Maîtriser la \cle{lecture analytique} (identification du thème, de la thèse, de la problématique, des arguments, schématisation de la structure argumentative sous forme d'arbre).
    \item Traiter rigoureusement les \cle{questions} (explication, analyse, évaluation) en respectant les attendus de chaque type.
    \item Construire un \cle{plan détaillé d'essai personnel} (introduction, développement dialectique, conclusion) en intégrant des \cle{exemples concrets camerounais}.
    \item Développer l'\cle{esprit critique} et l'auto-évaluation via la \cle{grille officielle de correction} (échange de copies entre groupes, notation croisée, synthèse collective).
    \item Gérer le \cle{temps} (\SI{3}{h} simulées) et la \cle{présentation} (lisibilité, structure, langue).
\end{itemize}

\courssub{Situation}
\textbf{Contexte :} Vous êtes en situation d'examen du Probatoire de Philosophie, série Technique, session 2025. La durée de l'épreuve est de \SI{3}{h}. Le sujet porte sur la notion de \cle{Liberté} (Unité 3 du programme : \emph{La condition humaine}), plus précisément sur le rapport entre liberté, raison et contrainte.

\textbf{Sujet complet :}

\begin{center}
\textbf{RÉPUBLIQUE DU CAMEROUN \hfill PROBATOIRE TECHNIQUE 2025}\\
\textbf{ÉPREUVE DE PHILOSOPHIE \hfill Durée : 3 heures}\\
\textbf{Séries : Toutes séries techniques \hfill Coefficient : 2}
\end{center}

\vspace{0.5cm}

\textbf{TEXTE À ÉTUDIER}
\begin{quote}
\og On confond souvent la liberté avec l'indépendance. L'indépendant est celui qui ne dépend de personne, qui fait ce qu'il veut. Mais faire ce que l'on veut, est-ce être libre ? Le capricieux fait ce qu'il veut ; il est esclave de ses caprices. Le libertin fait ce qu'il veut ; il est esclave de ses passions. L'homme libre, au contraire, obéit à la raison. Obéir à la raison, ce n'est pas subir une contrainte extérieure, c'est se donner sa propre loi. C'est l'autonomie. Là où il y a contrainte, il y a un maître extérieur ; là où il y a obéissance à la raison, il n'y a que moi, moi qui me commande à moi-même.

La loi que je me donne n'est pas arbitraire ; elle est universelle. Elle est la loi de ma nature raisonnable. En la suivant, je réalise ma véritable nature. Je ne suis pas libre quand je cède à mon penchant pour la paresse ou la colère ; je suis libre quand je me lève pour travailler, quand je maîtrise ma colère pour juger sainement. La liberté n'est pas le pouvoir de tout faire ; elle est le pouvoir de faire le bien, c'est-à-dire l'action conforme à la raison.

Ainsi, la liberté vraie ne s'exerce que dans l'effort. Elle est conquête, non don. Elle suppose une éducation de la volonté, une ascèse. Celui qui n'a pas appris à se commander reste toute sa vie un enfant gâté, balloté par les désirs du moment. La société même ne peut exister sans cette liberté responsable : le citoyen libre n'est pas celui qui réclame tous ses droits sans devoirs, mais celui qui assume ses obligations envers la cité. \fg
\end{quote}
\textit{(Texte adapté d'Alain, \emph{Propos sur l'éducation} et \emph{Définitions})}

\vspace{0.5cm}

\textbf{QUESTIONS (12 points)}
\begin{enumerate}
    \item Quel est le thème de ce texte ? Justifiez votre réponse en vous appuyant sur des éléments précis du texte. (2 pts)
    \item Dégagez la thèse de l'auteur et les principaux arguments qui la soutiennent. (4 pts)
    \item Expliquez la distinction fondamentale que l'auteur établit entre \emph{indépendance} (ou caprice) et \emph{liberté}. (3 pts)
    \item \og La liberté suppose-t-elle l'absence de toute contrainte ? \fg Répondez à cette question en vous appuyant sur le texte et sur vos connaissances personnelles. (3 pts)
\end{enumerate}

\textbf{ESSAI PERSONNEL (8 points)}
\begin{quote}
\og Obéir à la raison, est-ce être libre ? \fg
Vous traiterez ce sujet en organisant votre devoir en trois parties (thèse, antithèse, synthèse) ou selon un plan dialectique pertinent. Vous illustrerez votre argumentation par au moins un exemple concret tiré de la vie sociale camerounaise (ex. : débat sur la liberté d'expression des jeunes sur les réseaux sociaux à Douala, respect du code de la route à Yaoundé, engagement associatif, etc.).
\end{quote}

\courssub{Consignes}
Travaillant par groupes de 4 élèves, vous disposez de \textbf{\SI{2,5}{h}} pour réaliser les tâches suivantes (les \SI{30}{min} restantes seront consacrées à l'échange et à l'évaluation) :

\begin{enumerate}
    \item \textbf{Fiche d'analyse du texte (\SI{30}{min}) :} Sur une feuille A4, réalisez la fiche complète : Thème -- Thèse -- Problématique implicite -- Arguments (numérotés, avec citations courtes) -- \textbf{Schéma de la structure argumentative sous forme d'arbre} (racine = thèse, branches = arguments/parties du texte).
    \item \textbf{Rédaction des réponses aux questions 2 et 4 (\SI{45}{min}) :} Rédigez \emph{in extenso} les réponses aux questions n°2 (Thèse/Arguments) et n°4 (Question d'évaluation). Respectez la méthode : Citation $\rightarrow$ Reformulation $\rightarrow$ Explication/Analyse $\rightarrow$ Illustration/Exemple personnel pour la Q4.
    \item \textbf{Plan détaillé de l'essai (\SI{45}{min}) :} Rédigez le plan complet de l'essai \og Obéir à la raison, est-ce être libre ? \fg.
    \begin{itemize}
        \item Introduction : Accroche, définition des termes, problématique, annonce du plan.
        \item Développement : 3 grandes parties (I, II, III), chacune divisée en 2 sous-parties (A, B) avec \textbf{arguments philosophiques} (auteurs, concepts) et \textbf{exemples concrets camerounais} (au moins un par grande partie).
        \item Conclusion : Bilan, réponse nuancée, ouverture.
    \end{itemize}
    \item \textbf{Évaluation croisée et synthèse (\SI{30}{min}) :} Échangez votre production avec un autre groupe. Évaluez-la à l'aide de la \textbf{grille de critères officiels} ci-dessous. Notez sur 20. Retournez la copie. En grand groupe, le professeur mène la synthèse : bonnes pratiques, erreurs fréquentes, gestion du temps.
\end{enumerate}

\textbf{Grille d'évaluation simplifiée (sur 20) :}
\begin{center}
\begin{tabularx}{\linewidth}{|l|X|c|}\hline
\textbf{Critères} & \textbf{Indicateurs} & \textbf{Pts} \\ \hline
Lecture analytique & Thème/Thèse/Arguments justes, Arbre cohérent & 4 \\ \hline
Questions & Réponses structurées, citations, exemples personnels (Q4) & 5 \\ \hline
Plan d'essai & Structure dialectique, arguments philosophiques, exemples camerounais, rédaction soignée & 8 \\ \hline
Présentation/Langue & Orthographe, syntaxe, connecteurs logiques, lisibilité & 3 \\ \hline
\end{tabularx}
\end{center}

\courssub{Ressources à mobiliser}
\begin{methode}[Mémento : Les 3 temps de l'exercice sur texte]
\textbf{1. Lecture analytique (\SIrange{30}{40}{min}) :}
\begin{enumerate}
    \item \textbf{Lectures successives :} 1ère : globale (sujet général). 2ème : repérage (mots-clés, connecteurs, ponctuation). 3ème : structurée (découpage en séquences/arguments).
    \item \textbf{Fiche :} Thème (1 mot/phrase), Thèse (phrase affirmative), Problématique (question tension), Arguments (Idée directrice + Citation + Explication).
    \item \textbf{Arbre argumentatif :} Thèse au centre ; branches principales = parties du texte ; sous-branches = arguments/exemples.
\end{enumerate}
\textbf{2. Questions (\SIrange{60}{70}{min}) :}
\begin{enumerate}
    \item \textbf{Type Explication (Q1, Q3) :} Citer $\rightarrow$ Reformuler $\rightarrow$ Expliquer le sens/portée $\rightarrow$ Articuler au reste du texte.
    \item \textbf{Type Analyse (Q2) :} Thèse + Arguments ordonnés (logique du texte).
    \item \textbf{Type Évaluation (Q4) :} Position nuancée (Oui/Non/Mais) $\rightarrow$ Argument 1 (texte) $\rightarrow$ Argument 2 (connaissance perso/auteur) $\rightarrow$ Exemple concret $\rightarrow$ Mini-synthèse.
\end{enumerate}
\textbf{3. Essai (\SIrange{80}{90}{min}) :}
\begin{enumerate}
    \item \textbf{Brouillon intensif :} Problématisation $\rightarrow$ Idées/Arguments/Exemples $\rightarrow$ Plan détaillé (I, A, B ; II, A, B ; III, A, B).
    \item \textbf{Rédaction :} Intro (AMPCI), Développement (1 paragraphe = 1 idée = Argument + Explication + Exemple + Transition), Conclusion (Bilan + Réponse + Ouverture).
    \item \textbf{Relecture :} Hors-sujet ? Structure ? Langue ? Exemples camerounais présents ?
\end{enumerate}
\end{methode}

\begin{aretenir}[Références philosophiques utiles (Programme 1re Tech)]
\begin{itemize}
    \item \textbf{Alain :} \cle{Liberté} = \cle{Autonomie} = Obéissance à la raison/loi morale (vs caprice/passion). \emph{Propos sur l'éducation}.
    \item \textbf{Kant :} \cle{Autonomie de la volonté} (Fondement de la métaphysique des mœurs). Liberté = obéissance à la loi qu'on se donne (\cle{Impératif catégorique}). \cle{Hétéronomie} = soumission aux inclinations.
    \item \textbf{Sartre :} L'homme est condamné à être libre (L'existentialisme est un humanisme). Liberté radicale, angoisse, responsabilité. Pas de nature humaine prédéfinie.
    \item \textbf{Spinoza :} Liberté = nécessité comprise (Éthique). Être libre, c'est connaître les causes qui nous déterminent (\cle{conatus}, raison).
    \item \textbf{Rousseau :} Obéir à la loi qu'on s'est prescrite, c'est la liberté (Du contrat social). Passage de l'indépendance naturelle à la liberté civile/morale.
\end{itemize}
\end{aretenir}

\begin{attention}[Pièges fréquents au Probatoire]
\begin{itemize}
    \item \textbf{Paraphrase :} Répéter le texte au lieu de l'analyser (Q2, Q3).
    \item \textbf{Hors-sujet essai :} Traiter \og La liberté \fg au lieu de \og Obéir à la raison \fg.
    \item \textbf{Plan "catalogue" :} Enchaîner des exemples sans articulation logique (thèse/antithèse/synthèse).
    \item \textbf{Oubli de l'exemple camerounais :} Consigne explicite non respectée = perte de points (critère "Ancrage contextuel").
    \item \textbf{Gestion du temps :} Passer \SI{2}{h} sur le texte $\rightarrow$ essai bâclé. \textbf{Règle :} \SI{1,5}{h} texte/questions, \SI{1,5}{h} essai.
\end{itemize}
\end{attention}

\courssub{Démarche et corrigé commenté}
\begin{exoresolu}[Corrigé complet de l'atelier -- Sujet Liberté/Alain]
\textbf{Énoncé (Rappel) :} Voir le sujet complet ci-dessus (Situation : Texte d'Alain, Questions 1 à 4, Sujet d'essai).
\tcblower
\textbf{ÉTAPE 1 : FICHE D'ANALYSE DU TEXTE (Modèle attendu)}

\textbf{Thème :} La nature de la vraie liberté (distinction liberté / indépendance / caprice).

\textbf{Thèse :} La vraie liberté n'est pas l'indépendance (faire ce qu'on veut), mais l'\cle{autonomie} (obéir à la raison/se donner sa propre loi universelle), ce qui suppose un effort d'\cle{éducation de la volonté}.

\textbf{Problématique implicite :} Comment distinguer la liberté authentique de l'illusion d'indépendance ? En quoi l'obéissance à la raison constitue-t-elle la condition de la liberté responsable ?

\textbf{Structure argumentative (Découpage en 3 séquences) :}

\begin{enumerate}
    \item \textbf{Séquence 1 (L.1-5) : Critique de l'indépendance/caprice.} 
    \begin{itemize}
        \item Arg 1 : Confusion courante liberté = indépendance (faire ce qu'on veut).
        \item Arg 2 : Contre-exemples : Capricieux (esclave caprices), Libertin (esclave passions).
        \item Arg 3 : Définition positive : Homme libre = obéit à la raison = \cle{Autonomie} (loi propre, pas maître extérieur).
    \end{itemize}
    \item \textbf{Séquence 2 (L.6-10) : Caractères de la loi rationnelle (Universalisation, Réalisation de soi).}
    \begin{itemize}
        \item Arg 4 : Loi non arbitraire, universelle (loi de la nature raisonnable).
        \item Arg 5 : Liberté = faire le bien (conforme à la raison) $\neq$ pouvoir tout faire.
        \item Arg 6 : Exemples concrets : Maîtrise de la paresse/colère = liberté ; Cession = servitude.
    \end{itemize}
    \item \textbf{Séquence 3 (L.11-16) : Genèse et dimension sociale de la liberté.}
    \begin{itemize}
        \item Arg 7 : Liberté = conquête/effort/\cle{ascèse} (éducation de la volonté), non don.
        \item Arg 8 : Sans autodiscipline $\rightarrow$ enfant gâté, esclave désirs immédiats.
        \item Arg 9 : Condition sociale : Citoyen libre = assume obligations/devoirs (pas seulement droits).
    \end{itemize}
\end{enumerate}

\textbf{Schéma arborescent (Arbre argumentatif) :}
\begin{popfigure}
\begin{center}\tcbox[colback=popink,colframe=popink,coltext=white,arc=9pt,boxsep=4pt,left=6pt,right=6pt]{\bfseries\small Thèse Liberté = Autonomie (Obéir à la raison)}\end{center}
\vspace{-2pt}
\begin{tcbraster}[raster columns=2,raster equal height=rows,raster column skip=6pt,raster row skip=6pt,enhanced,blankest]
\begin{tcolorbox}[enhanced,colback=white,colframe=popblue,boxrule=0.9pt,arc=3pt,title={Séquence 1 Critique Indépendance},fonttitle=\bfseries\scriptsize,coltitle=white,colbacktitle=popblue,left=4pt,right=4pt,top=2pt,bottom=2pt,boxsep=1pt,toptitle=1.5pt,bottomtitle=1.5pt,halign=left]
\begin{itemize}[nosep,leftmargin=*,itemsep=1pt,font=\scriptsize]
\item Arg 1: Confusion Liberté=Indépendance
\item Arg 2: Capricieux/Libertin = Esclaves passions
\item Arg 3: Libre = Obéit raison = Autonomie (Loi propre)
\end{itemize}
\end{tcolorbox}
\begin{tcolorbox}[enhanced,colback=white,colframe=poporange,boxrule=0.9pt,arc=3pt,title={Séquence 2 Nature Loi Rationnelle},fonttitle=\bfseries\scriptsize,coltitle=white,colbacktitle=poporange,left=4pt,right=4pt,top=2pt,bottom=2pt,boxsep=1pt,toptitle=1.5pt,bottomtitle=1.5pt,halign=left]
\begin{itemize}[nosep,leftmargin=*,itemsep=1pt,font=\scriptsize]
\item Arg 4: Loi universelle Nature raisonnable
\item Arg 5: Liberté = Faire le bien $\neq$ Tout faire
\item Arg 6: Ex: Maîtrise paresse/colère
\end{itemize}
\end{tcolorbox}
\begin{tcolorbox}[enhanced,colback=white,colframe=popgreen,boxrule=0.9pt,arc=3pt,title={Séquence 3 Genèse \& Social},fonttitle=\bfseries\scriptsize,coltitle=white,colbacktitle=popgreen,left=4pt,right=4pt,top=2pt,bottom=2pt,boxsep=1pt,toptitle=1.5pt,bottomtitle=1.5pt,halign=left]
\begin{itemize}[nosep,leftmargin=*,itemsep=1pt,font=\scriptsize]
\item Arg 7: Conquête/Effort Éducation volonté
\item Arg 8: Sans discipline = Enfant gâté
\item Arg 9: Citoyen libre = Devoirs + Droits
\end{itemize}
\end{tcolorbox}
\end{tcbraster}

\legende{Schéma arborescent de la structure argumentative du texte d'Alain (Thèse au centre, 3 séquences, 9 arguments).}
\end{popfigure}

\vspace{0.5cm}
\textbf{ÉTAPE 2 : RÉPONSES AUX QUESTIONS 2 ET 4 (Modèles de rédaction)}

\textbf{Question 2 : Dégagez la thèse de l'auteur et les principaux arguments qui la soutiennent. (4 pts)}

\emph{Modèle de réponse :}
\og \textbf{Thèse :} Alain soutient que la vraie liberté ne consiste pas dans l'indépendance (le fait de faire ce que l'on veut), mais dans l'\cle{autonomie}, c'est-à-dire l'obéissance à la raison qui nous fait nous donner notre propre loi universelle.

\textbf{Arguments (dans l'ordre du texte) :}
\begin{enumerate}
    \item \textbf{La confusion liberté/indépendance est une illusion.} L'indépendance (faire ce qu'on veut) caractérise le capricieux et le libertin qui sont, paradoxalement, des \og esclaves \fg de leurs caprices ou passions (L.2-4).
    \item \textbf{L'homme libre obéit à la raison, ce qui est de l'autonomie.} Obéir à la raison n'est pas subir une contrainte extérieure (\cle{hétéronomie}), mais se commander à soi-même : \og là où il y a obéissance à la raison, il n'y a que moi, moi qui me commande à moi-même \fg (L.4-5).
    \item \textbf{La loi rationnelle est universelle et réalise la nature humaine.} Elle n'est pas arbitraire ; c'est la \og loi de ma nature raisonnable \fg. En la suivant, \og je réalise ma véritable nature \fg (L.6-7). La liberté est le \og pouvoir de faire le bien \fg (conforme à la raison), non le pouvoir de tout faire (L.8-9).
    \item \textbf{La liberté est une conquête par l'effort (\cle{ascèse}).} Elle suppose une \og éducation de la volonté \fg. Sans cela, l'homme reste un \og enfant gâté \fg balloté par les désirs (L.11-13).
    \item \textbf{La liberté responsable est le fondement du lien social.} Le citoyen libre n'est pas celui qui réclame seulement des droits, mais celui qui \og assume ses obligations envers la cité \fg (L.14-16).
\end{enumerate} \fg

\textbf{Commentaire prof :} La réponse cite explicitement (L.2-4, L.4-5, etc.), reformule avec ses mots (\og hétéronomie \fg, \og universalisation \fg), structure par arguments numérotés. Note maximale si l'enchaînement est fluide.

\vspace{0.5cm}
\textbf{Question 4 : \og La liberté suppose-t-elle l'absence de toute contrainte ? \fg (3 pts)}

\emph{Modèle de réponse :}
\og \textbf{Non, la liberté authentique ne suppose pas l'absence de toute contrainte, mais la distinction entre contrainte extérieure (aliénation) et contrainte intérieure rationnelle (\cle{autonomie}).}
\begin{itemize}
    \item \textbf{D'après le texte :} Alain oppose explicitement la \og contrainte \fg (maître extérieur, hétéronomie) à l'\og obéissance à la raison \fg (autonomie). Le capricieux \og fait ce qu'il veut \fg (absence de contrainte extérieure) mais n'est \og pas libre \fg car il subit la contrainte intérieure de ses passions. À l'inverse, l'homme libre s'impose la contrainte de la raison (\og je me commande à moi-même \fg) pour être vraiment libre.
    \item \textbf{Apport personnel (Kant) :} Kant distingue l'\emph{hétéronomie} (volonté déterminée par des mobiles extérieurs/sensibles = contrainte) de l'\emph{autonomie} (volonté déterminée par la loi morale qu'elle se donne = liberté). La contrainte morale (devoir) est l'expression de ma liberté, non sa négation.
    \item \textbf{Exemple concret camerounais :} Prenons le \textbf{respect du Code de la route à Yaoundé}. Un conducteur qui \og fait ce qu'il veut \fg (grille un feu rouge, roule à contresens pour aller plus vite) croit exercer son indépendance. En réalité, il se met en danger et met autrui en danger ; il est esclave de son impatience (passion). Le conducteur qui s'arrête au feu rouge (contrainte légale extérieure \textbf{et} contrainte intérieure de la raison/sécurité) exerce sa liberté responsable : il assure sa sécurité et celle des autres, condition de sa mobilité future.
    \item \textbf{Nuance (Sartre/Spinoza) :} Sartre dirait que même en prison (contrainte physique maximale), je reste libre de \emph{signifier} ma situation (attitude face à l'emprisonnement). Spinoza ajouterait que comprendre la contrainte (nécessité) me rend libre par rapport à elle (\cle{conatus}, \cle{servitude} vs liberté).
\end{itemize}
\textbf{En conclusion,} la liberté suppose l'absence de contrainte \emph{aliénante} (extérieure, passionnelle), mais elle \emph{exige} la contrainte \emph{libérante} de la raison et de la loi (morale, juridique) que l'on s'impose à soi-même. \fg

\textbf{Commentaire prof :} Structure dialectique claire (Non $\rightarrow$ Texte $\rightarrow$ Auteur/Concept $\rightarrow$ Exemple Cameroun $\rightarrow$ Nuance $\rightarrow$ Synthèse). L'exemple du Code de la route à Yaoundé est parfait : concret, local, illustre contrainte extérieure/intérieure.

\vspace{0.5cm}
\textbf{ÉTAPE 3 : PLAN DÉTAILLÉ DE L'ESSAI \og Obéir à la raison, est-ce être libre ? \fg}

\textbf{INTRODUCTION}
\begin{itemize}
    \item \textbf{Accroche :} \og L'obéissance est la mère de la réussite, et la raison en est le guide \fg (proverbe bamiléké adapté) / Citation Kant : \og L'autonomie est le fondement de la dignité humaine \fg.
    \item \textbf{Définitions :} Obéir = se soumettre à un ordre. Raison = faculté de connaître, de juger, de se déterminer par des principes universels. Liberté = pouvoir de self-détermination (\cle{autonomie}) vs indépendance (absence de contrainte).
    \item \textbf{Problématique :} Obéir implique-t-il soumission (perte de liberté) ou, au contraire, l'obéissance à la raison (loi universelle, loi de ma nature) est-elle la condition \emph{sine qua non} de l'accès à la liberté véritable ?
    \item \textbf{Annonce du plan :} Nous verrons d'abord que l'obéissance à la raison semble contredire l'idée commune de liberté comme indépendance (I). Ensuite, que l'obéissance à la raison est la condition de l'autonomie, donc de la liberté vraie (II). Enfin, que cette liberté rationnelle ne nie pas les déterminismes mais les assume dans une responsabilité concrète, notamment citoyenne (III).
\end{itemize}

\textbf{DÉVELOPPEMENT}

\textbf{I. Obéir à la raison : une contrainte apparente, une hétéronomie ? (Thèse : Non, ce n'est pas être libre au sens vulgaire)}
\begin{itemize}
    \item \textbf{A. La raison comme contrainte extérieure (\cle{Hétéronomie}).} Si la raison m'est imposée (dogme, tradition, autorité scientifique non comprise), j'obéis sans être libre. Ex : L'élève technique qui applique une formule de résistance des matériaux sans en comprendre la démonstration physique ; il \og obéit \fg mais n'est pas libre intellectuellement (Alain : \og l'élève qui récite sa leçon sans la comprendre \fg).
    \item \textbf{B. L'illusion de l'indépendance (Caprice/Passion).} Vouloir \og faire ce qu'on veut \fg sans raison, c'est obéir à ses pulsions (Spinoza : \cle{servitude}). Ex camerounais : Un jeune artisan à Douala qui refuse les normes de sécurité (port de casque, gants) au nom de \og sa liberté \fg de travailler \og à sa guise \fg ; il subit la contrainte de l'accident ou de l'amende. Il n'est pas libre, il est esclave de son imprudence.
\end{itemize}

\textbf{II. Obéir à la raison : l'autonomie, fondement de la liberté vraie (Antithèse : Oui, c'est la condition de la liberté)}
\begin{itemize}
    \item \textbf{A. L'autonomie kantienne : se donner sa propre loi.} \og Autonomie \fg = \emph{autos} (soi) + \emph{nomos} (loi). Obéir à la raison (\cle{Impératif catégorique}), c'est obéir à la loi que ma raison législative universelle. Je suis l'auteur de la loi $\rightarrow$ Je suis libre. \textit{Exemple :} L'étudiant qui choisit de réviser son Probatoire le soir (contrainte raison) au lieu de sortir (pulsion) : il s'obéit à lui-même (projet de réussite).
    \item \textbf{B. La raison comme connaissance des causes (Spinoza).} \og La liberté est une nécessité comprise \fg. Obéir à la raison, c'est comprendre les déterminismes (physiques, psychologiques, sociaux) pour agir \emph{en connaissance de cause}. \textit{Exemple camerounais :} Un agriculteur du Nord qui, au lieu de subir la sécheresse (contrainte naturelle), adopte l'agroforesterie et les semences améliorées (raison scientifique). Il \og obéit \fg aux lois de l'agronomie pour \emph{gagner} sa liberté alimentaire.
\end{itemize}

\textbf{III. La liberté raisonnable : responsabilité, effort et engagement citoyen (Synthèse : Liberté comme conquête éthique et politique)}
\begin{itemize}
    \item \textbf{A. Liberté comme effort et éducation (Alain/Rousseau).} \og La liberté n'est pas un don, c'est une conquête \fg. Obéir à la raison demande une \cle{ascèse}, une \cle{éducation de la volonté}. \textit{Exemple :} La formation en école technique (ENIEG, Lycées techniques) : la discipline (port d'uniforme, horaires, rigueur du geste professionnel) est contrainte au début ; elle devient liberté professionnelle (compétence, employabilité, autonomie financière) à la fin.
    \item \textbf{B. Liberté et citoyenneté : le pacte social (Rousseau).} \og Obéir à la loi qu'on s'est prescrite, c'est la liberté \fg (\cle{Pacte social}). La loi rationnelle (votée par les représentants) n'est pas contrainte mais protection de la liberté de tous. \textit{Exemple concret :} \textbf{La liberté d'expression sur les réseaux sociaux à Douala.} Un jeune blogueur critique la gestion municipale. \textbf{Liberté irrationnelle (caprice) :} Insultes, fake news, appel à la haine $\rightarrow$ Poursuites judiciaires (Loi sur la cybercriminalité 2010), prison. Il a cru être libre, il a perdu sa liberté. \textbf{Liberté rationnelle (autonomie) :} Argumentation factuelle, respect de la déontologie, proposition constructive $\rightarrow$ Débat public enrichi, autorité interpellée, citoyenneté exercée. Il a \og obéi à la raison \fg (vérité, respect d'autrui, loi) et a \emph{gagné} en liberté effective (crédibilité, impact).
\end{itemize}

\textbf{CONCLUSION}
\begin{itemize}
    \item \textbf{Bilan :} Obéir à la raison n'est pas s'asservir ; c'est passer de l'hétéronomie (subir : nature, passions, autorité aveugle) à l'autonomie (s'autodéterminer par le savoir et la loi morale).
    \item \textbf{Réponse nuancée :} C'est être libre \emph{à la condition} que la raison obéie soit \emph{la mienne} (appropriée, comprise, universalisable) et non un dogme imposé.
    \item \textbf{Ouverture :} Cette liberté rationnelle est-elle accessible à tous ? Question de l'éducation (Alain), des conditions matérielles (Marx), de l'accès au numérique au Cameroun (fracture numérique). La technique (notre série) est-elle un outil de cette raison libératrice ou un nouveau déterminisme (IA, algorithmes) ?
\end{itemize}

\textbf{Commentaire prof sur le plan :} Structure dialectique parfaite (Thèse/Antithèse/Synthèse). Arguments philosophiques nommés (Kant, Spinoza, Rousseau, Alain). Exemples camerounais \textbf{variés} (artisan, agriculteur, élève technique, blogueur Douala) et \textbf{pertinents} (illustrent exactement l'argument). Plan rédigé en phrases complètes, connecteurs logiques (\textit{dès lors, par conséquent, en revanche, c'est pourquoi}). Prêt pour la rédaction en \SI{1,25}{h}.

\end{exoresolu}

\courssub{Bilan de l'activité}
Cette activité d'intégration a permis de mettre en évidence plusieurs points décisifs pour la réussite au Probatoire :

\begin{enumerate}
    \item \textbf{L'unité de la méthode :} Les élèves ont compris que la fiche d'analyse (Étape 1) \emph{est} le squelette des réponses aux questions (Étape 2) et du plan d'essai (Étape 3). L'arbre argumentatif construit au brouillon réapparaît dans la structure de l'essai (Séquence 1 $\rightarrow$ Partie I, Séquence 2 $\rightarrow$ Partie II, Séquence 3 $\rightarrow$ Partie III).
    \item \textbf{La distinction cruciale Explication / Analyse / Évaluation :} L'évaluation croisée a sanctionné les copies qui \og expliquaient \fg au lieu d'\og analyser \fg (Q2) ou qui donnaient un \og avis personnel \fg sans ancrage textuel ni conceptuel (Q4). La grille a objectivé la notation.
    \item \textbf{L'ancrage camerounais comme levier de pertinence :} Les groupes qui ont utilisé l'exemple du \textbf{blogueur à Douala} (liberté d'expression / cybercriminalité) ou de l'\textbf{artisan/agriculteur} ont produit des essais nettement supérieurs à ceux restés abstraits. Cela valide l'exigence du programme : \og Appliquer les notions à des situations concrètes de la vie courante \fg.
    \item \textbf{La gestion du temps :} La simulation chronométrée a révélé que \SI{45}{min} pour la fiche d'analyse + arbre est un maximum ; il faut viser \SI{30}{min} pour garder \SI{1,5}{h} pour l'essai. L'entraînement régulier (1 sujet par mois) est indispensable.
    \item \textbf{Les erreurs types à éradiquer :} (a) Arbre argumentatif \og plat \fg (liste) au lieu de hiérarchisé (Thèse $\rightarrow$ Arguments principaux $\rightarrow$ Sous-arguments). (b) Introduction d'essai sans \textbf{problématique explicite}. (c) Plan d'essai en \og I. Pour / II. Contre / III. Mon avis \fg au lieu d'une vraie dialectique (Thèse/Antithèse/Synthèse dépassant l'alternative). (d) Oubli des \textbf{connecteurs logiques} (\textit{dès lors, par conséquent, en revanche, c'est pourquoi}) qui font la cohérence de la démonstration.
\end{enumerate}

\emph{Recommandation finale :} Conservez ce corrigé commenté comme \textbf{modèle type} pour vos révisions. Refaites l'exercice seul(e) chez vous dans les \SI{3}{h} exactes, puis auto-évaluez-vous avec la grille. La maîtrise de cette méthode est la clé de la réussite en Philosophie au Probatoire Technique.

\newpage

\coursec{Méthodes et exercices résolus}

\coursec{L'exercice sur texte}

\courssub{Généralités sur l'exercice sur texte}

\begin{definition}[L'exercice sur texte (commentaire de texte)]
L'exercice sur texte, épreuve obligatoire du Probatoire / Baccalauréat technique en Philosophie, consiste à expliquer un extrait d'œuvre philosophique (auteur identifié, thème au programme) et à en discuter la thèse. Il évalue la capacité à \textbf{comprendre}, \textbf{analyser}, \textbf{argumenter} et \textbf{prendre position} de manière rigoureuse. La note repose sur la fidélité au texte, la clarté de l'explication, la pertinence des arguments personnels et la qualité de la langue.
\end{definition}

\begin{methode}[Organiser sa copie et gérer son temps (3 h à 4 h selon coefficient)]
\begin{enumerate}
\item \textbf{Repérage (15--20 min) :} lire le texte 3 fois, identifier auteur/thème/thèse, découper en moments, souligner concepts et arguments.
\item \textbf{Réponses aux questions (1 h -- 1 h 30) :} traiter dans l'ordre, une question = un paragraphe structuré (idée + citation + explication).
\item \textbf{Essai personnel (1 h 30 -- 2 h) :} brouillon obligatoire (plan détaillé : thèse, position, 2 arguments + exemples, objection, conclusion), puis rédaction soignée.
\item \textbf{Relecture (15 min) :} orthographe, syntaxe, numérotation des questions, aération de la copie (saut de ligne entre parties).
\end{enumerate}
\textbf{Consigne absolue :} ne jamais recopier le texte en entier. Citer court (quelques mots à une phrase) entre guillemets français « ... ».
\end{methode}

\courssub{La lecture analytique du texte}

\begin{methode}[Lire et analyser un texte philosophique]
Pour appliquer les notions de l'unité à une situation concrète, il faut maîtriser la lecture analytique en six étapes :
\begin{enumerate}
\item \textbf{Première lecture globale.} Lire deux fois sans écrire. Repérer : l'auteur, l'ouvrage (si indiqué), le thème général (liberté, devoir, État, travail, technique, etc.).
\item \textbf{Repérage de la thèse.} Souligner la phrase-clé (souvent en début ou fin de texte). Se demander : que veut démontrer l'auteur ? Contre quelle thèse adverse s'oppose-t-il ?
\item \textbf{Découpage en moments.} Diviser le texte en 2 ou 3 étapes logiques (moments). Chaque moment = une idée-force ou un argument. Noter en marge le résumé de chaque moment.
\item \textbf{Analyse des concepts.} Définir les \cle{mots-clés} du texte (ex. : \cle{devoir}, \cle{liberté}, \cle{conscience}, \cle{loi}, \cle{raison}) en mobilisant les notions du cours.
\item \textbf{Repérage des arguments et exemples.} Pour chaque moment : quel raisonnement ? quel exemple ou contre-exemple ? quelle référence implicite ?
\item \textbf{Mise en relation avec le concret.} Chercher une situation camerounaise (vie scolaire, familiale, professionnelle, citoyenne) qui illustre la thèse.
\end{enumerate}
\textbf{Réflexe à retenir :} un texte philosophique = une \cle{thèse} + des \cle{arguments} + des \cle{exemples}. L'analyse vise à dégager cette structure.
\end{methode}

\begin{exoresolu}[Analyse d'un texte sur le devoir]
\textbf{Énoncé.} Lisez le texte suivant, puis dégagez la thèse de l'auteur et les moments du texte.

\emph{« Le devoir n'est pas une contrainte extérieure qui pèserait sur l'homme comme une charge. Celui qui accomplit son devoir par peur du châtiment agit comme un esclave. Le véritable accomplissement du devoir suppose au contraire que l'homme se donne à lui-même sa propre loi : il obéit alors à sa raison, et c'est en cela qu'il est libre. Ainsi, loin de s'opposer à la liberté, le devoir en est la plus haute expression. »}

\tcblower
\textbf{Solution détaillée.}

\textbf{Étape 1 : lecture globale.} Le texte porte sur le rapport entre le \cle{devoir} et la \cle{liberté}. Le thème appartient à la réflexion morale (unité « Devoir et liberté »).

\textbf{Étape 2 : repérage de la thèse.} La thèse est dans la dernière phrase : « loin de s'opposer à la liberté, le devoir en est la plus haute expression ». L'auteur soutient que le devoir authentique \emph{réalise} la liberté au lieu de la nier.

\textbf{Étape 3 : découpage en moments.}
\begin{itemize}
\item \emph{Premier moment} (« Le devoir n'est pas... comme un esclave ») : critique de la conception hétéronome du devoir (contrainte extérieure, obéissance par peur = esclavage).
\item \emph{Deuxième moment} (« Le véritable accomplissement... il est libre ») : exposition de la conception autonome : le devoir = obéissance à la loi que la raison se donne à elle-même (autonomie).
\item \emph{Troisième moment} (« Ainsi... la plus haute expression ») : conclusion synthétique : le devoir est l'accomplissement de la liberté.
\end{itemize}

\textbf{Étape 4 : analyse des concepts.}
\begin{itemize}
\item \cle{Devoir} : obligation morale, ce qu'on \emph{doit} faire, non par contrainte physique mais par respect pour la loi morale.
\item \cle{Liberté} : ici au sens d'\cle{autonomie} (se donner sa propre loi), non d'indépendance (faire ce qu'on veut).
\item \cle{Raison} : faculté de se prescrire des lois universelles ; source de la loi morale en soi.
\end{itemize}

\textbf{Étape 5 : arguments.} Argument central : l'opposition esclave / homme libre. Obéir par peur = hétéronomie (loi venue d'autrui) = esclavage. Obéir à sa raison = autonomie (loi venue de soi) = liberté.

\textbf{Étape 6 : illustration concrète (Cameroun).} Un fonctionnaire qui refuse un pot-de-vin non par peur de la Cour suprême, mais parce que sa raison lui dit que la corruption détruit le bien commun, accomplit son devoir en homme libre.
\end{exoresolu}

\courssub{Les réponses aux questions}

\begin{methode}[Répondre aux questions posées sur le texte]
Les questions portent sur le sens, les concepts, les arguments ou la portée du texte.
\begin{enumerate}
\item \textbf{Analyser la consigne} : souligner le verbe opératoire (\emph{expliquez}, \emph{définissez}, \emph{dégagez}, \emph{montrez}, \emph{justifiez}, \emph{quelle conception... critique-t-il ?}). Chaque verbe impose une réponse type.
\item \textbf{Localiser l'appui textuel} : repérer la/les phrase(s) exacte(s) qui permettent de répondre. La réponse doit \textbf{toujours} s'ancrer dans le texte.
\item \textbf{Rédiger en phrases complètes} : reprendre les termes de la question dans la première phrase. Pas de style télégraphique.
\item \textbf{Citer court et expliquer} : insérer une citation brève entre guillemets « ... », puis l'expliquer \emph{avec ses propres mots}. Une citation sans explication = 0 point.
\item \textbf{Vérifier l'adéquation} : la réponse correspond-elle exactement à la demande ? (Ni hors-sujet, ni développement inutile).
\end{enumerate}
\textbf{Formule type :} \textbf{Réponse} = Idée directrice (reprise de la question) + Citation textuelle (« ... ») + Explication personnelle (reformulation + sens).
\end{methode}

\begin{exoresolu}[Réponses aux questions sur le texte sur le devoir]
\textbf{Énoncé.} Reprendre le texte ci-dessus.

Question 1 : Quelle conception du devoir l'auteur critique-t-il ?

Question 2 : Expliquez la formule : « il obéit alors à sa raison, et c'est en cela qu'il est libre ».

\tcblower
\textbf{Solution détaillée.}

\textbf{Question 1.} La consigne « quelle conception... critique-t-il » appelle l'identification de la thèse adverse.

\textbf{Réponse rédigée :} L'auteur critique la conception du devoir comme \cle{contrainte extérieure} subie passivement. Comme il l'écrit : « le devoir n'est pas une contrainte extérieure qui pèserait sur l'homme comme une charge ». Cela signifie que le devoir ne doit pas être réduit à un poids imposé du dehors. L'auteur précise : « celui qui accomplit son devoir par peur du châtiment agit comme un esclave » : obéir seulement pour éviter une sanction, c'est se comporter en être hétéronome, dépourvu de liberté morale.

\textbf{Question 2.} La consigne « expliquez » demande l'explicitation du sens profond de la formule.

\textbf{Réponse rédigée :} Cette formule signifie que la véritable liberté est \cle{autonomie} : n'obéir qu'à la loi que notre propre raison a reconnue comme juste. Quand l'homme agit par devoir rationnel, il est l'auteur de sa propre loi ; il n'obéit qu'à lui-même, donc il est libre. Par exemple, un élève de Première technique qui révise ses cours non par crainte du zéro, mais parce qu'il comprend que la réussite conditionne son avenir professionnel, obéit à sa raison : il exerce sa liberté. La liberté n'est donc pas l'absence de règle, mais l'obéissance consentie à une loi rationnelle.
\end{exoresolu}

\courssub{L'essai personnel}

\begin{methode}[Construire et rédiger l'essai personnel]
L'essai personnel (souvent noté sur 6 à 8 points) demande de \emph{discuter} la thèse : donner un avis argumenté, pas un résumé.
\begin{enumerate}
\item \textbf{Reformuler la thèse} de l'auteur en une phrase claire (preuve de compréhension).
\item \textbf{Prendre position explicitement} : « Je partage / Je ne partage pas / Je partage partiellement cette thèse car... ». Annoncer la couleur dès l'introduction.
\item \textbf{Développer au moins deux arguments personnels}, chacun en un paragraphe : (Idée) $\rightarrow$ (Explication) $\rightarrow$ (Exemple concret camerounais : école, famille, marché, administration, coutume, actualité).
\item \textbf{Envisager une objection} (antithèse) : « On pourrait objecter que... » puis la \textbf{réfuter} ou en \textbf{nuancer} la portée (synthèse).
\item \textbf{Conclure} : rappeler sa position, résumer la démarche, ouvrir sur une question plus large (ex. : « Cela invite à repenser l'éducation civique... »).
\item \textbf{Soigner l'expression} : connecteurs logiques (\emph{d'abord, ensuite, en effet, cependant, par conséquent, enfin}), phrases complètes, orthographe soignée.
\end{enumerate}
\textbf{Structure type :} Introduction (thèse + position) -- Développement (Arg 1 + Ex 1 ; Arg 2 + Ex 2) -- Objection / Réponse -- Conclusion (bilan + ouverture).
\end{methode}

\begin{exoresolu}[Rédaction d'un essai personnel]
\textbf{Énoncé.} « L'auteur affirme que le devoir est la plus haute expression de la liberté. Partagez-vous ce point de vue ? Justifiez par des arguments et des exemples. »

\tcblower
\textbf{Solution détaillée (modèle d'essai personnel).}

\emph{Introduction -- Reformulation et position.} L'auteur soutient que le devoir, loin d'aliéner, réalise la liberté lorsque l'obéissance naît de la raison (autonomie). \textbf{Nous partageons cette thèse} : la liberté sans loi se retourne contre elle-même.

\emph{Premier argument -- La maîtrise de soi.} En effet, faire « tout ce qu'on veut » livre l'homme à ses pulsions. Le jeune qui passe ses nuits sur les réseaux sociaux au lieu de dormir se croit libre, mais devient esclave de sa fatigue et de l'addiction ; il prépare son échec au Probatoire. À l'inverse, l'élève qui s'impose un planning de révision obéit à sa raison : il se rend libre de choisir son orientation. \textbf{Exemple :} un apprenti soudeur à Douala qui respecte les consignes de sécurité non par peur de l'amende, mais par conscience du danger, protège sa vie et son avenir : il est libre.

\emph{Deuxième argument -- La condition sociale de la liberté.} Ensuite, la vie en société prouve que le respect des devoirs communs garantit la liberté de tous. Sur la route Yaoundé--Douala, le conducteur qui respecte le code de la route accomplit un devoir ; grâce à cette règle partagée, chacun circule en sécurité. Si chacun refusait toute contrainte au nom d'une « liberté absolue », ce serait le chaos et l'insécurité pour tous (loi du plus fort). Le devoir civique est le socle de la liberté publique.

\emph{Objection et réponse.} On objectera que certains devoirs imposés sont injustes (lois discriminatoires, coutumes oppressives). \textbf{C'est vrai}, mais l'auteur précise que le devoir authentique vient de la \cle{raison}, non de la \cle{peur}. Résister à une loi injuste au nom de la justice (ex. : lutte pour les droits de l'homme au Cameroun) est alors un \emph{devoir de conscience} librement accompli. L'objection ne contredit pas la thèse, elle en précise la portée : seul le devoir \emph{rationnel} libère.

\emph{Conclusion.} Ainsi, le devoir fondé sur la raison n'étouffe pas la liberté : il la constitue. La vraie liberté n'est pas l'arbitraire (« faire n'importe quoi »), mais l'autonomie (« se donner une loi juste »). Éduquer à la liberté, c'est éduquer au devoir.
\end{exoresolu}

\courssub{Application guidée : un sujet complet corrigé}

\begin{exoresolu}[Sujet complet type Probatoire technique]
\textbf{Énoncé.} Texte : \emph{« L'homme naît libre, mais partout il est dans les fers. Se croire le maître des autres, c'est ne pas cesser d'être plus esclave qu'eux. »} (Rousseau, \emph{Du contrat social}, Livre I, chap. 1)

Question 1 : Quelle est la thèse de ce texte ?

Question 2 : Expliquez l'expression « partout il est dans les fers ».

Question 3 : Pensez-vous que la société soit un obstacle à la liberté ? Justifiez votre réponse.

\tcblower
\textbf{Solution détaillée.}

\textbf{Question 1 -- Thèse.} La thèse de Rousseau est double : (1) L'homme possède une liberté naturelle originelle (« l'homme naît libre ») ; (2) La vie en société historique a aliéné cette liberté (« partout il est dans les fers ») par l'instauration de rapports de domination. Pis, le dominateur lui-même est esclave de sa propre domination (« se croire le maître des autres, c'est ne pas cesser d'être plus esclave qu'eux »). La société actuelle est donc un état de fait illégitime qu'il faut refonder sur un contrat légitime.

\textbf{Question 2 -- Explication de « partout il est dans les fers ».} L'expression est une \cle{métaphore} : les « fers » désignent les chaînes du prisonnier. Rousseau veut dire que dans \emph{toutes} les sociétés existantes (partout), l'individu est soumis à des lois, coutumes, autorités, inégalités économiques qu'il n'a pas choisies. \textbf{Exemple concret :} un jeune Camerounais doit payer des impôts, respecter le code du travail, subir la hiérarchie patronale, se conformer aux attentes familiales : il n'a pas la liberté « sauvage » de l'homme naturel qui ne dépendait que de lui-même.

\textbf{Question 3 -- Essai personnel : La société, obstacle ou condition de la liberté ?}

\emph{Introduction.} Rousseau dénonce les fers de la société factice. \textbf{Nuance :} la société limite la liberté naturelle, mais nous soutenons qu'une société \emph{juste} (État de droit) est la \emph{condition} de la liberté réelle (civile).

\emph{Argument 1 -- La contrainte réelle (donner raison à Rousseau).} Certes, la société impose des contraintes : lois, horaires, discipline scolaire, impôts. L'élève en classe ne peut pas sortir quand il veut ; le citoyen ne peut pas rouler à contresens. En ce sens, Rousseau a raison : la vie sociale « enchaine » la volonté particulière. \textbf{Exemple :} le couvre-feu sanitaire (Covid-19) a restreint la liberté de circulation au Cameroun au nom de la santé publique.

\emph{Argument 2 -- La protection de la liberté (dépasser Rousseau).} Cependant, sans société organisée, c'est la loi du plus fort : le faible n'a \emph{aucune} liberté. Grâce aux lois, à la police, à la justice, chaque Camerounais peut étudier, travailler, voter, circuler en sécurité. La règle commune (contrat social) \emph{protège} la liberté de chacun contre l'arbitraire. \textbf{Exemple :} le Code du travail protège l'employé contre le licenciement abusif ; la Constitution garantit la liberté d'expression.

\emph{Objection -- Les sociétés oppressives.} On objectera que des États tyranniques transforment la société en prison (dictatures, corruption généralisée). \textbf{Réponse :} Cela prouve qu'il faut distinguer \emph{société factice} (domination) et \emph{société légitime} (contrat). Le remède n'est pas le retour à l'état de nature (impossible), mais la transformation politique pour que la loi exprime la \cle{volonté générale}.

\emph{Conclusion.} La société enchaîne si ses lois sont injustes (fait), mais elle libère si ses lois sont justes (droit). La liberté véritable n'est pas l'indépendance sauvage, mais l'autonomie citoyenne dans un État de droit. « Obéir à la loi qu'on s'est prescrite, c'est la liberté » (Rousseau, \emph{Émile}).
\end{exoresolu}

\begin{aretenir}[Points clés pour réussir l'exercice sur texte au Probatoire]
\begin{itemize}
\item \textbf{Lire 3 fois} : globale $\rightarrow$ structure $\rightarrow$ concepts.
\item \textbf{Thèse + Moments + Concepts + Arguments} = squelette de l'analyse.
\item \textbf{Questions} : Consigne (verbe) $\rightarrow$ Citation $\rightarrow$ Explication personnelle.
\item \textbf{Essai} : Position claire $\rightarrow$ 2 arguments + 2 exemples concrets (Cameroun) $\rightarrow$ Objection traitée $\rightarrow$ Conclusion ouverte.
\item \textbf{Forme} : Aération, numérotation, relecture, français correct.
\end{itemize}
\end{aretenir}

\begin{attention}[Les 5 erreurs qui coûtent des points au Bac]
\begin{enumerate}
\item \textbf{Recopier le texte au lieu de l'expliquer.} Le correcteur évalue \emph{votre} compréhension. Citez court (une phrase max entre guillemets) et expliquez toujours la citation.
\item \textbf{Répondre hors sujet.} Beaucoup racontent le cours au lieu de répondre à la question précise. \textbf{Soulignez le verbe de la consigne} et n'en sortez pas.
\item \textbf{Oublier de prendre position dans l'essai.} « L'auteur a raison » ne suffit pas. Il faut dire « \textbf{Je pense que...} parce que... » et justifier par \emph{vos} arguments.
\item \textbf{Arguments sans exemples concrets.} Un argument abstrait non illustré (vie scolaire, famille, quartier, actualité camerounaise) est fragile. \textbf{1 argument = 1 exemple daté/localisé.}
\item \textbf{Négliger la langue et la présentation.} Fautes d'accord, phrases sans verbe, copie illisible, questions non numérotées, pas de saut de ligne entre parties : perte de points assurée. \textbf{Relisez-vous 10 min.}
\end{enumerate}
\end{attention}

\newpage

\coursec{Exercices}

\courssub{Je vérifie mes connaissances}

\begin{exercice}[QCM : les étapes de l'exercice sur texte]
Pour chaque question, une seule réponse est correcte. Entoure la bonne lettre.
\begin{enumerate}
\item La première étape du travail sur un texte philosophique est :
\begin{enumerate}
\item[a)] rédiger l'essai personnel ;
\item[b)] la lecture analytique du texte ;
\item[c)] la recherche d'exemples ;
\item[d)] la rédaction de la conclusion.
\end{enumerate}
\item La \emph{thèse} d'un texte est :
\begin{enumerate}
\item[a)] le sujet général dont parle l'auteur ;
\item[b)] la position que l'auteur défend et cherche à établir ;
\item[c)] la question que le texte se pose ;
\item[d)] le résumé des exemples du texte.
\end{enumerate}
\item Répondre à une question de compréhension consiste à :
\begin{enumerate}
\item[a)] donner son opinion personnelle ;
\item[b)] recopier le texte intégralement ;
\item[c)] dégager le sens d'un passage en s'appuyant sur le texte ;
\item[d)] critiquer l'auteur.
\end{enumerate}
\item L'essai personnel doit :
\begin{enumerate}
\item[a)] paraphraser le texte ;
\item[b)] discuter la thèse de l'auteur avec des arguments et des exemples personnels ;
\item[c)] ignorer le texte ;
\item[d)] se limiter à une phrase.
\end{enumerate}
\end{enumerate}
\end{exercice}

\begin{exercice}[Vrai ou faux ? Justifie ta réponse]
Indique si chaque affirmation est vraie ou fausse, puis justifie en une ou deux phrases.
\begin{enumerate}
\item La problématique d'un texte est la question à laquelle la thèse de l'auteur répond.
\item Dans les réponses aux questions, il est interdit de citer le texte.
\item L'essai personnel peut prendre position contre la thèse de l'auteur, à condition d'argumenter.
\item Le thème et la thèse d'un texte sont deux expressions de la même chose.
\end{enumerate}
\end{exercice}

\begin{exercice}[Définitions]
Définis en deux ou trois phrases chacune des notions suivantes, telles qu'elles s'appliquent à l'exercice sur texte :
\begin{enumerate}
\item le \cle{thème} ;
\item la \cle{thèse} ;
\item la \cle{problématique} ;
\item l'\cle{essai personnel}.
\end{enumerate}
\end{exercice}

\begin{exercice}[Repérage rapide]
On donne le passage suivant : \og L'homme qui agit sans réfléchir ressemble à la machine : il exécute. Mais l'homme qui délibère avant d'agir affirme sa conscience, et c'est en cela qu'il est libre. \fg
\begin{enumerate}
\item Quel est le thème de ce passage ?
\item Quelle est la thèse défendue ?
\item Formule la problématique sous forme d'une question.
\item Relève un mot-clé et explique son sens dans ce contexte.
\end{enumerate}
\end{exercice}

\courssub{Je m'entraîne}

\begin{exercice}[Identifier thème, thèse et problématique en 10 minutes]
Lis le texte suivant, extrait adapté des écrits de Bergson sur la conscience.

\og La conscience est la lumière qui éclaire nos actions. L'animal agit par instinct, sans se représenter ce qu'il fait ; l'homme, lui, se connaît agissant. Il peut se retourner sur lui-même, juger ses propres pensées, hésiter entre plusieurs conduites. Cette faculté de se saisir soi-même est ce qui distingue l'esprit humain. Sans elle, point de choix véritable, point de responsabilité : nous serions emportés par le courant des habitudes comme une feuille par le vent. \fg

En 10 minutes maximum :
\begin{enumerate}
\item dégage le thème du texte ;
\item formule la thèse de l'auteur en une phrase ;
\item construis la problématique sous forme interrogative ;
\item trace l'arbre de la structure argumentative du texte (thèse au sommet, arguments et exemples en branches).
\end{enumerate}
\end{exercice}

\begin{exercice}[Expliquer un passage]
Dans le texte de Bergson ci-dessus, explique en cinq à huit lignes le passage suivant, en dégageant son sens et sa fonction dans l'argumentation : \og Sans elle, point de choix véritable, point de responsabilité : nous serions emportés par le courant des habitudes comme une feuille par le vent. \fg
Tu veilleras à : 1) reformuler l'idée sans paraphraser mot à mot ; 2) expliquer la comparaison de la feuille ; 3) montrer à quelle question ce passage répond.
\end{exercice}

\begin{exercice}[Correction d'une copie d'élève]
Voici la réponse d'un élève à la question : \og Quelle est la thèse du texte de Bergson sur la conscience ? \fg

\og Dans ce texte, l'auteur parle de la conscience. Il dit que la conscience est la lumière qui éclaire nos actions et que l'animal agit par instinct. L'homme se connaît agissant et il peut se retourner sur lui-même. Moi je pense que la conscience c'est important dans la vie, par exemple quand on va au marché il faut réfléchir avant d'acheter. \fg

\begin{enumerate}
\item Repère dans cette copie au moins trois erreurs méthodologiques (paraphrase, confusion entre thème et thèse, réponse hors de la question posée, absence de formulation claire).
\item Réécris une réponse correcte en trois à cinq lignes.
\end{enumerate}
\end{exercice}

\begin{exercice}[Rédiger une réponse justifiée]
À partir du texte de Bergson sur la conscience, réponds en dix lignes maximum à la question suivante : \og Selon l'auteur, en quoi la conscience distingue-t-elle l'homme de l'animal ? \fg
Ta réponse devra : 1) annoncer l'idée directrice ; 2) s'appuyer sur deux éléments précis du texte (courtes citations entre guillemets) ; 3) se conclure par une phrase de synthèse.
\end{exercice}

\courssub{Je me prépare au Bac}

\begin{exercice}[Sujet type Probatoire technique : la conscience]
\textbf{Texte} (extrait adapté de Bergson) :

\og La conscience est la lumière qui éclaire nos actions. L'animal agit par instinct, sans se représenter ce qu'il fait ; l'homme, lui, se connaît agissant. Il peut se retourner sur lui-même, juger ses propres pensées, hésiter entre plusieurs conduites. Cette faculté de se saisir soi-même est ce qui distingue l'esprit humain. Sans elle, point de choix véritable, point de responsabilité : nous serions emportés par le courant des habitudes comme une feuille par le vent. Mais la conscience demande un effort : le plus souvent, nous préférons l'automatisme de l'habitude, qui nous dispense de réfléchir. Vivre pleinement, c'est maintenir éveillée cette attention de l'esprit à lui-même. \fg

\textbf{Questions de compréhension et d'analyse}
\begin{enumerate}
\item Quel est le thème de ce texte ?
\item Quelle est la thèse défendue par l'auteur ? Justifie ta réponse par une courte citation.
\item Explique le sens du passage : \og le plus souvent, nous préférons l'automatisme de l'habitude, qui nous dispense de réfléchir \fg.
\item Quelle est la portée de la thèse de l'auteur : s'applique-t-elle à tous les hommes, ou seulement à certains ? Justifie à partir du texte.
\end{enumerate}

\textbf{Essai personnel}

La conscience fait-elle la grandeur de l'homme ? Tu discuteras cette question en prenant clairement position, avec des arguments et des exemples tirés de ta vie quotidienne ou de la société camerounaise.
\end{exercice}

\begin{exercice}[Sujet type Baccalauréat technique : le devoir et la liberté]
\textbf{Texte} (extrait adapté) :

\og On croit souvent que le devoir est l'ennemi de la liberté : obéir à une règle, n'est-ce pas renoncer à faire ce que l'on veut ? Mais cette opposition est trompeuse. Celui qui cède à tous ses désirs n'est pas libre : il est esclave de ses caprices, comme le commerçant esclave de la mode du moment. Le devoir, au contraire, est la règle que la raison se donne à elle-même. Accomplir son devoir, c'est donc agir non pas sous la contrainte d'autrui, mais selon une loi que l'on s'est librement prescrite. L'élève qui travaille alors que ses camarades l'invitent à la distraction n'est pas moins libre qu'eux : il est plus libre, car il est maître de lui-même. \fg

\textbf{Questions}
\begin{enumerate}
\item Dégage le thème et la problématique de ce texte.
\item Quelle est la thèse de l'auteur ? Contre quelle opinion commune s'élève-t-elle ?
\item Explique le sens du passage : \og Celui qui cède à tous ses désirs n'est pas libre : il est esclave de ses caprices \fg.
\item Quelle est la portée de l'exemple de l'élève dans l'argumentation ?
\end{enumerate}

\textbf{Essai personnel}

Accomplir son devoir, est-ce renoncer à sa liberté ? Tu illustreras ta réflexion par au moins deux exemples concrets de la vie courante camerounaise (école, travail, famille, vie citoyenne).
\end{exercice}

\begin{exercice}[Entraînement chronométré : sujet complet en 45 minutes]
\textbf{Texte} (extrait adapté) :

\og Beaucoup confondent le bonheur avec le plaisir. Pourtant le plaisir est passager : il naît, brille un instant et s'éteint, laissant souvent un vide plus grand que lui. Le bonheur, lui, est un état durable de l'âme, qui ne dépend pas des biens extérieurs mais de la manière dont nous jugeons notre vie. Celui qui attend le bonheur des choses — l'argent, les honneurs, les fêtes — court après une ombre. Le sage, au contraire, trouve le contentement en lui-même, dans l'accord entre ses actions et sa conscience. \fg

\textbf{Consigne} : en 45 minutes, réponds aux trois questions suivantes, puis auto-évalue ta copie à l'aide de la grille ci-dessous.
\begin{enumerate}
\item Quelle distinction l'auteur établit-il entre le plaisir et le bonheur ?
\item Explique : \og Celui qui attend le bonheur des choses court après une ombre \fg.
\item Essai personnel (15 à 20 lignes) : le bonheur dépend-il uniquement de nous-mêmes ?
\end{enumerate}

\textbf{Grille d'auto-évaluation} (attribue 0, 1 ou 2 points par critère) :

\begin{center}
\begin{tabularx}{\linewidth}{|l|X|c|}\hline
\rowcolor{popblueL} \textbf{Critère} & \textbf{Attendu} & \textbf{Note} \\ \hline
Exactitude & Le thème, la thèse et le sens des passages sont correctement dégagés & /2 \\ \hline
Appui sur le texte & Les réponses citent ou désignent précisément le texte & /2 \\ \hline
Justification & Chaque réponse est argumentée, sans paraphrase & /2 \\ \hline
Essai personnel & Position claire, arguments, exemples concrets & /2 \\ \hline
Expression & Langue correcte, présentation soignée & /2 \\ \hline
\end{tabularx}
\end{center}

Calcule ton total sur 10 et indique les deux points que tu dois améliorer en priorité.
\end{exercice}

\newpage

\coursec{Corrigés des exercices}

\coursec{L'exercice sur texte}

\courssub{Généralités sur l'exercice sur texte}

\begin{methode}[Les quatre étapes incontournables]
L'exercice sur texte, épreuve majeure du \cle{Probatoire} et du \cle{Baccalauréat technique}, suit un ordre logique immuable :
\begin{enumerate}
\item \textbf{Lecture analytique} : comprendre le texte (thème, thèse, problématique, structure, vocabulaire clé).
\item \textbf{Questions de compréhension et d'analyse} : répondre précisément en s'appuyant \emph{exclusivement} sur le texte (citations obligatoires).
\item \textbf{Question de portée / synthèse} : dégager la portée de la thèse (universalité, limites, implications).
\item \textbf{Essai personnel} : discuter la thèse avec \emph{ses propres arguments et exemples} (vie courante, contexte camerounais).
\end{enumerate}
\emph{Astuce :} consacre 10 min à l'étape 1, 15 min à l'étape 2, 5 min à l'étape 3, 25 min à l'étape 4, 5 min à la relecture (total 60 min).
\end{methode}

\begin{definition}[Notions fondamentales]
\begin{itemize}
\item \textbf{\cle{Thème}} : sujet général du texte (ce dont il parle). Se formule en un mot ou groupe nominal (\og la conscience \fg, \og le devoir et la liberté \fg). On le repère par les mots récurrents.
\item \textbf{\cle{Thèse}} : position que l'auteur \emph{défend et cherche à établir}. Se formule en phrase affirmative complète : \og L'auteur soutient que... \fg. C'est la réponse à la problématique.
\item \textbf{\cle{Problématique}} : question à laquelle la thèse répond. Se formule sous forme interrogative : \og En quoi... ? \fg, \og Pourquoi... ? \fg. Elle relie thème et thèse.
\item \textbf{\cle{Essai personnel}} : discussion \emph{libre et argumentée} de la thèse. L'élève prend position (accord, désaccord, nuance) et justifie par des arguments et des exemples concrets (impératif : vie camerounaise).
\end{itemize}
\end{definition}

\begin{exemplebox}[Pièges classiques (QCM / Vrai-Faux)]
\begin{itemize}
\item \textbf{Thème $\neq$ Thèse} : \og La conscience \fg (thème) vs \og La conscience distingue l'homme de l'animal \fg (thèse).
\item \textbf{Problématique $\neq$ Thème} : \og La conscience \fg (thème) vs \og En quoi la conscience distingue-t-elle l'homme ? \fg (problématique).
\item \textbf{Compréhension $\neq$ Opinion} : une question sur le texte attend la pensée de l'auteur (avec citation), pas ton avis.
\item \textbf{Citation $\neq$ Paraphrase} : citer (court, entre guillemets) est \emph{recommandé} ; paraphraser (recopier en changeant des mots sans expliquer) est \emph{interdit}.
\item \textbf{Essai personnel $\neq$ Paraphrase} : l'essai va au-delà du texte ; il exige prise de position, arguments propres, exemples concrets.
\end{itemize}
\end{exemplebox}

\begin{exoresolu}[Entraînement rapide : QCM et Vrai-Faux]
\textbf{Énoncé :} Choisis la bonne réponse / Vrai ou Faux avec justification.
\tcblower
\textbf{QCM :}
\begin{enumerate}
\item L'exercice sur texte commence toujours par : a) l'essai personnel \textbf{b) la lecture analytique} c) les questions. \emph{Justification : sans analyse, pas de compréhension fiable.}
\item La thèse, c'est : a) le sujet général \textbf{b) la position défendue par l'auteur} c) la question posée. \emph{Justification : thème = sujet, thèse = réponse, problématique = question.}
\item Une question de compréhension demande : a) ton opinion b) de recopier le texte \textbf{c) de dégager le sens en s'appuyant sur le texte}. \emph{Justification : expliquer $\neq$ recopier ; opinion réservée à l'essai.}
\item L'essai personnel consiste à : a) résumer le texte \textbf{b) discuter la thèse avec arguments et exemples personnels} c) citer l'auteur. \emph{Justification : l'essai suppose la compréhension mais va plus loin.}
\end{enumerate}
\textbf{Vrai / Faux :}
\begin{enumerate}
\item \textbf{Vrai.} Problématique = question du texte ; Thèse = réponse de l'auteur. Exemple : Problématique : \og La conscience distingue-t-elle l'homme ? \fg $\rightarrow$ Thèse : \og La conscience est ce qui distingue l'homme. \fg
\item \textbf{Faux.} Citer est \emph{recommandé} (courte citation entre guillemets). Interdit : la paraphrase non commentée.
\item \textbf{Vrai.} L'essai personnel est une discussion libre (approuver, nuancer, réfuter) \emph{argumentée}.
\item \textbf{Faux.} Thème = sujet général (mots) ; Thèse = position de l'auteur (phrase complète).
\end{enumerate}
\end{exoresolu}

\courssub{La lecture analytique du texte}

\begin{methode}[Repérer thème, thèse, problématique en 10 minutes]
\begin{enumerate}
\item \textbf{Lis deux fois} : 1\textsuperscript{re} lecture globale (de quoi parle-t-il ?), 2\textsuperscript{de} lecture active (surligne connecteurs, répétitions, thèse explicite).
\item \textbf{Thème} : note les mots-clés qui reviennent $\rightarrow$ groupe nominal.
\item \textbf{Thèse} : cherche la phrase affirmative centrale (souvent en début/fin de paragraphe, ou après \og car \fg, \og donc \fg, \og c'est pourquoi \fg). Formule : \og L'auteur soutient que... \fg
\item \textbf{Problématique} : transforme la thèse en question : \og En quoi / Pourquoi / Comment [thèse] ? \fg
\item \textbf{Structure} : découpe en parties logiques (arguments, exemples, objections, nuances). Repère les \cle{mots-charnières} (\og or \fg, \og mais \fg, \og en effet \fg, \og par exemple \fg).
\end{enumerate}
\end{methode}

\begin{exoresolu}[Application : texte sur la liberté et la délibération (extrait type)]
\textbf{Énoncé :} \og L'homme qui délibère avant d'agir est libre ; celui qui exécute sans réfléchir est une machine. \fg
\tcblower
\begin{itemize}
\item \textbf{Thème} : la liberté et la réflexion (ou l'action réfléchie).
\item \textbf{Thèse} : L'auteur soutient que c'est la \cle{délibération} (réflexion avant l'action) qui rend l'homme libre.
\item \textbf{Problématique} : En quoi la réflexion avant l'action fait-elle la liberté de l'homme ? (Qu'est-ce qui distingue l'action libre de l'action machinale ?)
\item \textbf{Mot-clé} : \og \cle{délibère} \fg. Sens : peser les raisons d'agir, examiner les conduites possibles au lieu d'exécuter automatiquement. Porte l'opposition homme libre / homme-machine.
\end{itemize}
\end{exoresolu}

\begin{exoresolu}[Entraînement chronométré (10 min) : extrait de Bergson, \emph{L'Énergie spirituelle}]
\textbf{Énoncé :} Identifier thème, thèse, problématique et structure argumentative.
\tcblower
\begin{itemize}
\item \textbf{Thème} : la conscience (comme caractère propre de l'homme).
\item \textbf{Thèse} : Bergson soutient que la conscience, entendue comme la faculté de \og se connaître agissant \fg, distingue l'homme de l'animal et rend possible le choix et la responsabilité.
\item \textbf{Problématique} : En quoi la conscience constitue-t-elle la spécificité de l'esprit humain ?
\item \textbf{Arbre de la structure argumentative} :
\begin{itemize}
\item \textbf{Thèse} : La conscience distingue l'esprit humain.
\begin{itemize}
\item \emph{Argument 1 (opposition)} : L'animal agit par instinct sans se représenter ; l'homme \og se connaît agissant \fg.
\item \emph{Argument 2 (pouvoirs de la conscience)} : Se retourner sur soi, juger ses pensées, hésiter $\rightarrow$ choisir.
\item \emph{Argument 3 (par l'absurde)} : Sans conscience, pas de choix ni responsabilité $\rightarrow$ emportés par les habitudes.
\item \emph{Image illustrative} : La feuille emportée par le vent = homme privé de conscience.
\end{itemize}
\end{itemize}
\end{itemize}
\begin{attention}[Vigilance]
Ne confonds pas l'argument 3 (raisonnement par l'absurde : conséquences absurdes du contraire) et l'image de la feuille (illustration). L'argument \emph{prouve}, l'image \emph{éclaire}.
\end{attention}
\end{exoresolu}

\courssub{Les réponses aux questions}

\begin{methode}[Expliquer un passage (question de compréhension)]
\begin{enumerate}
\item \textbf{Reformule l'idée directrice} du passage en une phrase (sans paraphrase).
\item \textbf{Cite court} (guillemets) les segments clés.
\item \textbf{Explique} chaque citation : que signifie-t-elle ? quel rôle dans l'argumentation ? (argument, exemple, objection, nuance, conséquence).
\item \textbf{Replace} le passage dans l'économie du texte : à quelle question répond-il ? quel lien avec la thèse ?
\end{enumerate}
\emph{Barème type (3 pts) :} reformulation (1) + citation exploitée (1) + fonction argumentative (1).
\end{methode}

\begin{exoresolu}[Modèle : explication du passage sur la feuille (Bergson)]
\textbf{Passage :} \og Sans conscience... nous serions comme une feuille emportée par le vent. \fg
\tcblower
\textbf{Réponse rédigée :} Bergson tire les conséquences de sa thèse par l'absurde : si l'homme perdait la conscience, il perdrait le choix et la responsabilité. Choisir suppose de comparer des conduites, donc de se connaître et réfléchir ; sans cela, nos actes sont dictés par l'habitude (automatismes). La comparaison à la \og feuille emportée par le vent \fg rend sensible cette passivité : comme la feuille ne maîtrise pas sa trajectoire, l'homme sans conscience subit ses habitudes. Ce passage répond à : \og Que deviendrait l'homme sans conscience ? \fg et renforce la thèse : la conscience est condition de liberté et responsabilité.
\end{exoresolu}

\begin{methode}[Rédiger une réponse justifiée (question d'analyse)]
\begin{enumerate}
\item \textbf{Annonce l'idée directrice} (réponse synthétique à la question).
\item \textbf{Appuie sur le texte} : 2 citations précises minimum, chacune suivie de son explication.
\item \textbf{Articule} : montre comment les citations s'enchaînent pour démontrer l'idée.
\item \textbf{Conclus} par une phrase de synthèse reprenant la question.
\end{enumerate}
\emph{Interdiction :} opinion personnelle, paraphrase, citation sans commentaire.
\end{methode}

\begin{exoresolu}[Modèle : \og Selon Bergson, pourquoi la conscience distingue-t-elle l'homme de l'animal ? \fg]
\tcblower
\textbf{Réponse :} Selon Bergson, la conscience distingue l'homme de l'animal car elle est la faculté de se représenter sa propre action. L'animal \og agit par instinct, sans se représenter ce qu'il fait \fg : son comportement est déterminé, il ne peut ni hésiter ni se juger. L'homme, au contraire, \og se connaît agissant \fg : il peut \og se retourner sur lui-même, juger ses propres pensées, hésiter entre plusieurs conduites \fg. Cette hésitation ouvre l'espace du choix, là où l'animal exécute. Ainsi, la conscience est la marque propre de l'esprit humain, fondant choix véritable et responsabilité.
\textbf{Barème (4 pts) :} idée directrice (1) + 2 citations exploitées (1) + explication sans paraphrase (1) + synthèse (1).
\end{exoresolu}

\begin{exemplebox}[Correction d'erreurs fréquentes (copie élève type)]
\textbf{Sujet :} \og Formule la thèse de Bergson. \fg
\textbf{Copie élève :} \og L'auteur parle de la conscience. La conscience est la lumière qui éclaire nos actions. L'animal agit par instinct. Moi je pense que la conscience c'est important, par exemple au marché on doit être conscient. \fg
\textbf{Erreurs :}
\begin{enumerate}
\item \textbf{Confusion thème/thèse} : \og parle de la conscience \fg = thème, pas thèse.
\item \textbf{Paraphrase} : recopie \og lumière qui éclaire \fg, \og instinct \fg sans formuler l'idée directrice.
\item \textbf{Hors-sujet} : opinion (\og Moi je pense \fg) et exemple marché (essai personnel) dans une question de compréhension.
\item \textbf{Absence de formulation thèse} : pas de phrase \og La thèse est que... \fg
\end{enumerate}
\textbf{Modèle correct :} \og La thèse de Bergson est que la conscience, comprise comme la faculté de \og se connaître agissant \fg, distingue l'homme de l'animal : l'animal agit par instinct sans se représenter, l'homme peut se retourner sur lui-même, juger et choisir. La conscience est ainsi condition du choix et de la responsabilité. \fg
\end{exemplebox}

\courssub{L'essai personnel}

\begin{methode}[Structure type de l'essai personnel (25 min)]
\begin{enumerate}
\item \textbf{Introduction (3--4 lignes)} : Accroche (proverbe, fait de société) $\rightarrow$ Rappel thèse auteur $\rightarrow$ \textbf{Position claire annoncée} (\og Nous soutiendrons que... \fg / \og Nous nuancerons... \fg) $\rightarrow$ Annonce du plan (2--3 arguments).
\item \textbf{Développement (2--3 paragraphes)} : Un argument par paragraphe.
\begin{itemize}
\item Phrase d'introduction de l'argument.
\item Explication philosophique.
\item \textbf{Exemple concret camerounais obligatoire} (vie scolaire, familiale, professionnelle, sociale, politique).
\item Mini-synthèse / transition.
\end{itemize}
\item \textbf{Nuance / Objection (1 paragraphe)} : \og Certes... \fg / \og Cependant... \fg (montre l'esprit critique).
\item \textbf{Conclusion (2--3 lignes)} : Réponse définitive à la question + ouverture brève.
\end{enumerate}
\emph{Expression :} langue soignée, paragraphes visibles, connecteurs logiques (\og D'une part \fg, \og En outre \fg, \og En revanche \fg, \og Ainsi \fg).
\end{methode}

\begin{attention}[Points de vigilance Probatoire / Bac Tech]
\begin{itemize}
\item \textbf{Position claire dès l'intro} : une copie qui \og hésite \fg jusqu'au bout est pénalisée.
\item \textbf{Exemples concrets camerounais} : \emph{exigés} (au moins 2 au Bac, 1 au Probatoire). Pas d'exemples abstraits (\og l'homme \fg, \og la société \fg). Préfère : \og un élève de Première technique à Yaoundé \fg, \og un apprenti mécanicien à Douala \fg, \og un chef de famille à Bafoussam \fg, \og un électeur à Garoua \fg.
\item \textbf{Arguments $\neq$ Exemples} : l'exemple \emph{illustre} l'argument, ne le remplace pas.
\item \textbf{Respect du sujet} : ne pas changer de sujet, ne pas faire un cours général.
\end{itemize}
\end{attention}

\courssub{Application guidée : sujets complets corrigés}

\begin{exoresolu}[Sujet type Probatoire Technique : La conscience (Bergson)]
\textbf{Texte :} Extrait de \emph{L'Énergie spirituelle} (conscience, attention à soi, habitude, feuille/vent, sagesse).
\tcblower
\textbf{Questions de compréhension et d'analyse (10 pts)}
\begin{enumerate}
\item \textbf{Thème (2 pts) :} La conscience, envisagée comme \cle{attention de l'esprit à lui-même}.
\item \textbf{Thèse + Justification (3 pts) :} Bergson soutient que la conscience distingue l'homme de l'animal et que vivre pleinement consiste à maintenir éveillée cette attention. \textbf{Citation exigée :} \og Cette faculté de se saisir soi-même est ce qui distingue l'esprit humain \fg.
\item \textbf{Explication passage (3 pts) :} L'auteur affirme que la conscience exige un effort et que nous lui préférons souvent la facilité de l'habitude (automatisme). Ce passage nuance : la conscience n'est pas un état permanent, mais une vigilance à conquérir contre la routine. \textbf{Analyse de l'image :} l'ombre / la fuite (selon passage exact).
\item \textbf{Portée de la thèse (2 pts) :} Universalité de la \emph{capacité} (propre à \og l'esprit humain \fg), inégalité de l'\emph{exercice effectif} (\og le plus souvent nous préférons l'automatisme \fg). Réponse nuancée exigée.
\end{enumerate}

\textbf{Essai personnel (10 pts) : \og La conscience fait-elle la grandeur de l'homme ? \fg}
\textbf{Modèle rédigé :}
\textit{Introduction :} \og L'homme est un roseau pensant \fg (Pascal). La conscience, cette capacité de se savoir, est-elle ce qui fait la grandeur de l'homme ? Bergson la présente comme la marque de l'humanité. Nous soutiendrons, avec lui, que la conscience fait la grandeur de l'homme, tout en reconnaissant qu'elle est aussi sa charge.

\textit{Argument 1 : Grandeur par la liberté et la responsabilité.} La conscience élève l'homme au-dessus de la bête et de la machine. Elle permet de juger, choisir, se repentir. \textbf{Exemple :} Un élève de Première technique à Yaoundé qui, la veille du Probatoire, renonce à une veillée pour réviser, montre qu'il se connaît et se gouverne. C'est la conscience qui fonde la responsabilité morale/juridique : on ne juge pas un fou comme un conscient. Au village, le chef de famille qui rend compte à sa communauté, l'électeur qui vote en connaissance de cause, témoignent de cette dignité.

\textit{Argument 2 : Grandeur par la lucidité.} La conscience permet de donner du sens à sa vie, de ne pas subir. \textbf{Exemple :} L'artisan de Foumban qui perfectionne son art par réflexion sur ses erreurs, au lieu de répéter machinalement, construit son excellence.

\textit{Nuance : Le fardeau de la conscience.} La conscience fait hésiter, culpabiliser, souffrir (remords, angoisse) ; l'animal ignore cela. \og L'homme est un dieu tombé qui se souvient des cieux \fg (Lamartine).

\textit{Conclusion :} La conscience fait la grandeur de l'homme non parce qu'elle le rend heureux, mais parce qu'elle le rend \textbf{libre et responsable}. Elle est la condition de tout choix véritable. Mieux vaut une conscience qui tourmente qu'une existence de feuille emportée par le vent.
\textbf{Barème essai :} Position claire (2) + Arguments (3) + Exemples concrets camerounais (3) + Expression/Organisation (2).
\end{exoresolu}

\begin{exoresolu}[Sujet type Baccalauréat Technique : Le devoir et la liberté (Kant / Stoïcisme)]
\textbf{Texte :} Extrait opposant opinion commune (devoir = contrainte) et thèse (devoir = liberté vraie, obéissance à la loi rationnelle). Exemple de l'élève qui travaille vs camarades distraits. Comparaison commerçant esclave de la mode.
\tcblower
\textbf{Questions (10 pts)}
\begin{enumerate}
\item \textbf{Thème + Problématique (2 pts) :} Thème : rapport devoir/liberté. Problématique : Le devoir est-il contraire à la liberté ou en est-il l'expression véritable ?
\item \textbf{Thèse + Opinion réfutée (3 pts) :} Thèse : Le devoir, loin de détruire la liberté, en est la forme vraie (obéir à une loi que la raison se prescrit). Opinion commune réfutée : \og Le devoir est l'ennemi de la liberté \fg, obéir = renoncer à ses désirs.
\item \textbf{Explication passage (3 pts) :} Renversement : liberté $\neq$ satisfaction désirs. Céder aux caprices = dépendance (esclave des impulsions). \textbf{Image :} \og Commerçant esclave de la mode \fg = suivre désirs = subir contrainte intérieure/sociale. Vraie liberté = maîtrise de soi.
\item \textbf{Portée exemple élève (2 pts) :} L'exemple vérifie la thèse sur un cas concret accessible : l'élève qui accomplit son devoir scolaire contre le plaisir immédiat prouve qu'il est \og maître de lui-même \fg, donc plus libre que ceux qui suivent leurs envies. Fonction : illustrer et \emph{vérifier} la thèse par l'expérience courante.
\end{enumerate}

\textbf{Essai personnel (10 pts) : \og Accomplir son devoir, est-ce renoncer à sa liberté ? \fg}
\textbf{Modèle rédigé :}
\textit{Introduction :} L'opinion oppose devoir et liberté. Nous défendrons, comme le texte, que le devoir accompli librement est le signe d'une liberté véritable.

\textit{Argument 1 : Refuser le devoir = esclavage des désirs.} \textbf{Exemple :} Apprenti mécanicien à Douala qui dépense tout chaque soir, ne refuse aucune invitation $\rightarrow$ enchaîné aux habitudes, sans avenir. À l'inverse, élève d'un lycée technique de Yaoundé qui révise au lieu de sortir : contrainte apparente = libération réelle (Bac, métier).

\textit{Argument 2 : Devoir citoyen = liberté rationnelle.} \textbf{Exemple :} Contribuable qui paie ses impôts, parent qui scolarise ses enfants : obéissent à une règle que la raison approuve (bien commun + bien propre). Obéir à une loi juste qu'on se reconnaît = être raisonnable, pas esclave.

\textit{Nuance :} Devoirs imposés injustement / soumission aveugle $\neq$ liberté. Confirme : seul le devoir \emph{assumé par la raison} libère.

\textit{Conclusion :} Accomplir son devoir, ce n'est pas renoncer à sa liberté : c'est l'exercer, devenir maître de soi au lieu d'être esclave de ses caprices.
\textbf{Barème essai :} Position (2) + 2 exemples camerounais (3) + Arguments (3) + Expression (2).
\end{exoresolu}

\begin{exoresolu}[Entraînement chronométré (45 min) : Plaisir vs Bonheur (Sénèque / Épicure / Morale)]
\textbf{Texte :} Opposition plaisir (passager, extérieur, vide) / bonheur (durable, intérieur, contentement, accord conscience). Image de l'ombre pour biens extérieurs. Figure du sage.
\tcblower
\textbf{Question 1 -- Distinction (modèle) :} Trois points : 1) Durée : plaisir \og passager \fg (\og naît, brille, s'éteint \fg) vs bonheur \og état durable \fg. 2) Source : plaisir = biens extérieurs (argent, honneurs) vs bonheur = \og manière dont nous jugeons notre vie \fg (intérieur). 3) Effets : plaisir laisse \og vide plus grand \fg vs bonheur = \og contentement \fg stable (\og accord actions/conscience \fg). Plaisir = moment subi ; bonheur = état construit.

\textbf{Question 2 -- Explication image ombre (modèle) :} Chercher bonheur dans l'extérieur est vain : comme une ombre, ces biens semblent proches mais fuient quand on les saisit, ne rassasient jamais. Argent/honneurs $\rightarrow$ nouveaux désirs. Image = illusion + fuite perpétuelle. Réponse à : \og Où chercher le bonheur ? \fg $\rightarrow$ En soi (figure du sage).

\textbf{Question 3 -- Essai : \og Le bonheur dépend-il uniquement de nous-mêmes ? \fg}
\textbf{Modèle :}
\textit{Thèse :} Le texte dit : bonheur dépend de notre jugement, non des biens. Nous nuancerons : \textbf{en grande partie de nous, pas uniquement.}

\textit{Argument 1 (Pour) :} À conditions égales, notre jugement fait le bonheur. \textbf{Exemple :} Deux élèves même lycée technique Bafoussam, même concession modeste : l'un se lamente, l'autre se réjouit de ses progrès à l'atelier et fierté famille. Contentement du sage (accord actes/conscience) est en notre pouvoir. Jeune honnête au marché dort en paix ; malhonnête riche vit dans la crainte.

\textit{Argument 2 (Contre / Nuance) :} Dire \og uniquement \fg est injuste. Maladie grave, guerre, pauvreté extrême, deuil pèsent sur l'âme quelle que soit la sagesse. Nier cela = accuser les malheureux.

\textit{Conclusion :} Le bonheur dépend \textbf{d'abord} de nous (jugement, choix, conscience) mais \textbf{non exclusivement} : la sagesse nous rend maîtres de notre bonheur \emph{dans les limites que la vie impose}.

\textbf{Auto-évaluation (grille) :}
\begin{center}
\begin{tabularx}{\linewidth}{|l|X|c|}\hline
\rowcolor{popblueL} \textbf{Critère} & \textbf{Attendu} & \textbf{Note} \\ \hline
Exactitude & Thème, thèse, sens passages corrects & 2/2 \\ \hline
Appui sur texte & Citations précises dans réponses & 1/2 \\ \hline
Justification & Argumentées, sans paraphrase & 2/2 \\ \hline
Essai personnel & Position claire, arguments, exemples concrets & 2/2 \\ \hline
Expression & Langue correcte, présentation soignée & 1/2 \\ \hline
\end{tabularx}
\end{center}
\textbf{Total : 8/10.} \textbf{À améliorer :} 1) Citer davantage (1 citation/réponse). 2) Relecture 5 min (orthographe, présentation).
\begin{attention}[Gestion du temps (45 min)]
10 min lecture/questions 1-2 -- 25 min essai -- 5 min relecture. Position claire annoncée d'emblée. Expliquer l'\emph{image} (ombre), pas seulement l'idée.
\end{attention}
\end{exoresolu}

\begin{savaistu}[L'exercice sur texte au Cameroun : enjeux et repères]
\begin{itemize}
\item \textbf{Coefficient :} Philo coefficient 2 (séries techniques) au Probatoire et au Baccalauréat.
\item \textbf{Durée :} 3h (Bac) / 2h (Probatoire). L'exercice sur texte est souvent le \emph{seul} sujet proposé (pas de choix dissertation).
\item \textbf{Attentes correcteurs MINESEC :} Maîtrise des 4 notions (Thème, Thèse, Problématique, Structure) + Citations \emph{exploitées} + Exemples \emph{camerounais} datés/localisés + Français correct (orthographe, syntaxe, paragraphes).
\item \textbf{Ressources :} Manuels officiels (Njock, Mvondo, etc.), annales Probatoire/Bac Tech (disponibles sur OnBuch+), textes au programme (Bergson, Kant, Sénèque, Aristote, Descartes, Rousseau, auteurs africains : Mudimbe, Wiredu, etc.).
\end{itemize}
\end{savaistu}

\newpage

\coursec{Fiche bilan}

\begin{popfigure}
\centering
\begin{center}\tcbox[colback=popblue,colframe=popblue,coltext=white,arc=9pt,boxsep=4pt,left=6pt,right=6pt]{\bfseries\small Exercice sur texte}\end{center}
\vspace{-2pt}
\begin{tcbraster}[raster columns=3,raster equal height=rows,raster column skip=6pt,raster row skip=6pt,enhanced,blankest]
\begin{tcolorbox}[enhanced,colback=white,colframe=poporange,boxrule=0.9pt,arc=3pt,title={1. Généralités},fonttitle=\bfseries\scriptsize,coltitle=white,colbacktitle=poporange,left=4pt,right=4pt,top=2pt,bottom=2pt,boxsep=1pt,toptitle=1.5pt,bottomtitle=1.5pt,halign=left]
\begin{itemize}[nosep,leftmargin=*,itemsep=1pt,font=\scriptsize]
\item Nature : commentaire ou explication
\item Objectif : montrer la compréhension
\item Structure : 3 parties
\end{itemize}
\end{tcolorbox}
\begin{tcolorbox}[enhanced,colback=white,colframe=popgreen,boxrule=0.9pt,arc=3pt,title={2. Lecture analytique},fonttitle=\bfseries\scriptsize,coltitle=white,colbacktitle=popgreen,left=4pt,right=4pt,top=2pt,bottom=2pt,boxsep=1pt,toptitle=1.5pt,bottomtitle=1.5pt,halign=left]
\begin{itemize}[nosep,leftmargin=*,itemsep=1pt,font=\scriptsize]
\item Identification (auteur, œuvre, date)
\item Thèse centrale + arguments
\item Vocabulaire clé + concepts
\item Plan du texte (mouvements)
\end{itemize}
\end{tcolorbox}
\begin{tcolorbox}[enhanced,colback=white,colframe=poppurple,boxrule=0.9pt,arc=3pt,title={3. Questions},fonttitle=\bfseries\scriptsize,coltitle=white,colbacktitle=poppurple,left=4pt,right=4pt,top=2pt,bottom=2pt,boxsep=1pt,toptitle=1.5pt,bottomtitle=1.5pt,halign=left]
\begin{itemize}[nosep,leftmargin=*,itemsep=1pt,font=\scriptsize]
\item Type 1 : Explication (sens, portée)
\item Type 2 : Analyse (arguments, méthode)
\item Type 3 : Discussion (critique, ex.)
\item Règle : 1 idée = 1 paragraphe
\end{itemize}
\end{tcolorbox}
\begin{tcolorbox}[enhanced,colback=white,colframe=poppink,boxrule=0.9pt,arc=3pt,title={4. Essai personnel},fonttitle=\bfseries\scriptsize,coltitle=white,colbacktitle=poppink,left=4pt,right=4pt,top=2pt,bottom=2pt,boxsep=1pt,toptitle=1.5pt,bottomtitle=1.5pt,halign=left]
\begin{itemize}[nosep,leftmargin=*,itemsep=1pt,font=\scriptsize]
\item Sujet lié au texte
\item Dissertation (intro, dev., concl.)
\item Thèse propre + arguments + ex.
\end{itemize}
\end{tcolorbox}
\begin{tcolorbox}[enhanced,colback=white,colframe=popgold,boxrule=0.9pt,arc=3pt,title={5. Méthode globale},fonttitle=\bfseries\scriptsize,coltitle=white,colbacktitle=popgold,left=4pt,right=4pt,top=2pt,bottom=2pt,boxsep=1pt,toptitle=1.5pt,bottomtitle=1.5pt,halign=left]
\begin{itemize}[nosep,leftmargin=*,itemsep=1pt,font=\scriptsize]
\item Étape 1 : Lecture active (brouillon)
\item Étape 2 : Rédaction questions (1h30)
\item Étape 3 : Dissertation (2h)
\item Étape 4 : Relecture (orthographe, cohérence)
\end{itemize}
\end{tcolorbox}
\end{tcbraster}

\legende{Carte mentale de la leçon 21 : L'exercice sur texte (5 étapes clés).}
\end{popfigure}

\begin{bilan}
\courssub{Généralités sur l'exercice sur texte}
\begin{itemize}
  \item \textbf{Nature de l'épreuve :} Épreuve écrite du Probatoire (coefficient variable selon séries). Durée totale : 3h30 à 4h.
  \item \textbf{Structure imposée :} Deux parties distinctes mais liées :
        \begin{enumerate}
          \item \textbf{Questions sur le texte} (environ 1h30) : Explication, analyse, discussion d'extraits précis.
          \item \textbf{Essai personnel (Dissertation)} (environ 2h) : Sujet de philosophie en lien thématique avec le texte.
        \end{enumerate}
  \item \textbf{Objectif évalué :} Capacité à \cle{comprendre} un texte philosophique, à en \cle{restituer} la pensée propre (sans hors-texte), à l'\cle{analyser} (logique, arguments) et à \cle{problématiser} pour produire une réflexion personnelle argumentée.
\end{itemize}

\courssub{La lecture analytique (sur brouillon — 30 à 40 min)}
\begin{itemize}
  \item \textbf{1. Identification :} Auteur, œuvre, date, contexte historique/philosophique (si connu).
  \item \textbf{2. Thèse centrale :} Formuler en une phrase affirmative la réponse de l'auteur au problème traité.
  \item \textbf{3. Vocabulaire \& Concepts :} Repérer les termes techniques (ex: \emph{liberté}, \emph{conscience}, \emph{État}, \cle{technique}, \cle{aliénation}) et leur sens précis dans le texte.
  \item \textbf{4. Structure logique (Plan du texte) :} Découper en \textbf{mouvements} (séquences logiques : thèse, arguments, objections, nuances, conclusion). Numéroter les lignes.
  \item \textbf{5. Repérage des citations :} Surligner les phrases-clés qui seront citées dans les réponses.
\end{itemize}

\courssub{Réponses aux questions : Types \& Règles de rédaction}
\begin{center}
\begin{tabularx}{\linewidth}{|l|X|X|}\hline
\textbf{Type de question} & \textbf{Consigne type} & \textbf{Attendus / Méthode} \\ \hline
\textbf{Explication} & « Expliquez le sens de cette phrase... » / « Que veut dire l'auteur par... ? » & Reformuler \textbf{avec ses mots} le sens précis. Définir les concepts. Montrer le lien avec la thèse. \textbf{Ne pas donner son avis.} \\ \hline
\textbf{Analyse} & « Montrez comment l'auteur démontre que... » / « Quels sont les arguments... ? » & Reconstituer le \textbf{raisonnement} : Prémisses $\to$ Inférence $\to$ Conclusion. Identifier le type d'argument (logique, fait, autorité, analogie). \\ \hline
\textbf{Portée / Discussion} & « Cette thèse est-elle convaincante ? » / « Discutez l'idée selon laquelle... » & \textbf{Nuancer} : Accorder du crédit (points forts) + Critiquer (limites, contre-exemples, objections). \textbf{Exemples personnels obligatoires} (concrets, camerounais si possible). \\ \hline
\end{tabularx}
\end{center}
\textbf{Règles d'or :} Une réponse = un paragraphe structuré (Idée directrice + Développement + Citation courte + Mini-conclusion). Répondre \textbf{dans l'ordre} des questions. Ne pas recopier le texte, l'\cle{exploiter}.

\courssub{L'essai personnel (Dissertation) : Structure type}
\begin{itemize}
  \item \textbf{Analyse du sujet (15 min) :} Définir les termes, trouver la \cle{problématique} (tension, paradoxe), annoncer le \cle{plan} (dialectique : Thèse / Antithèse / Synthèse \textbf{ou} Thématique : Causes / Conséquences / Solutions).
  \item \textbf{Introduction (rédaction soignée) :} Accroche $\to$ Définitions $\to$ Reformulation du sujet $\to$ Problématique $\to$ Annonce du plan.
  \item \textbf{Développement (2 à 3 parties, 2 sous-parties chacune) :}
        \begin{itemize}
          \item \textbf{Argument} (idée philosophique, référence auteur/œuvre).
          \item \textbf{Explication} du mécanisme conceptuel.
          \item \textbf{Exemple concret} (fait social, historique, littéraire, expérience vécue au Cameroun : ex. \emph{la corruption, le bilinguisme, l'exode rural, les NTIC, la chefferie traditionnelle}).
        \end{itemize}
  \item \textbf{Conclusion :} Réponse synthétique à la problématique $\to$ Ouverture (autre domaine, actualité, auteur).
\end{itemize}

\courssub{Gestion du temps \& Checklist technique}
\begin{itemize}
  \item \textbf{Total : 3h30--4h.} Lecture/Analyse : 35 min. Questions : 1h15--1h30. Dissertation : 1h45--2h. Relecture : 15 min.
  \item \textbf{Présentation :} Sauter des lignes entre intro/dev/concl et entre chaque grande partie. Encadrer les citations (guillemets + n° lignes). Écrire lisiblement. Pas de correcteur liquide abusif (rature propre).
\end{itemize}
\end{bilan}

\begin{aretenir}[Checklist avant le Bac -- Exercice sur texte]
\begin{itemize}
  \item $\square$ Je sais distinguer \textbf{Explication} (restitution fidèle), \textbf{Analyse} (logique interne) et \textbf{Discussion} (évaluation critique + mon avis).
  \item $\square$ Je sais identifier la \textbf{thèse centrale} et le \textbf{plan logique} (mouvements) d'un texte en 3 lectures actives.
  \item $\square$ Je maîtrise la structure \textbf{Introduction / Développement / Conclusion} pour la dissertation.
  \item $\square$ Je sais construire une \textbf{problématique} (question philosophique ouverte) et non une simple question de cours.
  \item $\square$ Je prépare 3 à 5 \textbf{exemples concrets camerounais} réutilisables (éthique, politique, technique, savoir).
  \item $\square$ Je cite \textbf{précisément} le texte (numéros de lignes) dans mes réponses aux questions.
  \item $\square$ Je gère mon temps : \textbf{ne pas dépasser 1h30 sur les questions} pour garder 2h pour la dissertation.
  \item $\square$ Je relis pour corriger la syntaxe, l'orthographe et vérifier la \textbf{cohérence} entre le texte et mon essai.
  \item $\square$ Je n'utilise \textbf{jamais} de connaissances extérieures pour \emph{expliquer} le texte (seulement pour la discussion/dissertation).
  \item $\square$ Je structure chaque paragraphe : \textbf{1 idée force = 1 paragraphe} (argument + explication + exemple).
  \item $\square$ Je connais les \textbf{grands auteurs au programme} (Platon, Aristote, Descartes, Kant, Marx, Bergson, Sartre, Foucault, auteurs africains : Mudimbe, Oyono, Eboussi Boulaga) pour les références.
  \item $\square$ Je sais définir rigoureusement les \textbf{concepts clés} du programme (Liberté, Conscience, Devoir, Bonheur, Justice, Technique, Science, État, Culture, Religion).
\end{itemize}
\end{aretenir}

\findecours

\end{document}
